% 本文件由 LaTeX 排版 Agent 根据有效 translation.json 自动生成。
% author=Augustine of Hippo; work_id=1508; title=论灵魂及其起源

\chapter{\WorkTitle{论灵魂及其起源}}

% structure: p0001 source=h1 target=omitted reason=page-title-already-rendered-by-parent

\SourceRecord{Translated by Peter Holmes and Robert Ernest Wallis, and revised by Benjamin B. Warfield.；Nicene and Post-Nicene Fathers, First Series；Vol. 5.；Edited by Philip Schaff.；Buffalo, NY: Christian Literature Publishing Co.,；1887.；<http://www.newadvent.org/fathers/1508.htm>.}

% structure: p0001 source=h1 target=section reason=page-title

\section{论灵魂及其起源（卷一）}

\Emph{致隐修士\Person{雷纳特斯}。}\par

\Emph{从\Person{雷纳特斯}那里收到\Person{味增爵·维克多}的两卷书，此人反对\Person{奥古斯丁}关于灵魂本性的意见，以及他对灵魂起源的犹豫，\Person{奥古斯丁}指出这位年轻的反对者如何因其自负，妄图在如此高深的问题上下判断，从而陷入了不可容忍的错误。然后他继续表明，\Person{维克多}认为可以证明人类灵魂非由生育传承，而是天主在每个人出生时重新注入的那些圣经段落，是模棱两可的，不足以证实他的这种主张。}\par

% structure: p0004 source=h2 target=subsection reason=relative-heading-rank; base=h2

\subsection{第一章 \Person{雷纳特斯}曾送给他那些曾寄给他的书卷，这对他是一份好意。}

最亲爱的弟兄\Person{雷纳特斯}，你的真诚、你的兄弟般的仁慈，以及我们彼此之间的相爱之情，我们早已有了明证；但如今你给了我更明确的证据，因为你寄给我两卷书，作者是我素昧平生的人——尽管并不因此就该轻视——名叫\Person{味增爵·维克多}（因为我确实他的作品标题下看到了这样署名的形式）：这是去年夏天你寄出的；但由于我不在家，直到秋末它们才辗转到我手里。以你对我极大的关爱，你怎会不在有机会或愿意时，将这类任何人写的作品（只要落入你手中，即使它们是写给别人的）让我知晓呢？何况当我的名字被提及和读到——而且是在反驳我曾在一些短论中发表的某些话语的上下文里。你做这一切的方式，正是你作为我最诚挚、最亲爱的朋友所必然采取的方式。\par

% structure: p0006 source=h2 target=subsection reason=relative-heading-rank; base=h2

\subsection{第二章 他以和善忍耐的心情接受一位年轻无经验者的书卷，此人以傲慢的口吻反对他。\Person{味增爵·维克多}脱离罗加提亚派而归依。}

然而，我有点难过，因为您的圣德对我的了解远不如我所期望的；您以为我收到您的信时会认为您是在伤害我，因为您让我知道了别人的所作所为。您可以从这个事实看出这种情绪离我有多远：我甚至没有就从他那里受到的委屈提出任何抱怨。因为，当他持有与我不同的观点时，他是否应该保持沉默？这无疑甚至应该使我感到愉快，因为他打破了沉默，使我们能够阅读他所说的话。我当然认为，他应该写信给我，而不是通过别人来写关于我；但由于他不认识我，他不敢亲自冒昧地反驳我的话。在他看来，没有必要就一个他自认为毫无疑义、反倒持有完全明朗和确定意见的问题向我请教。此外，他还听从了一位朋友的劝告，他告诉我们，是这位朋友迫使他写作的。如果他在争论中表达了对我不敬的言辞，我更愿意认为他这样做并不是有意要对我无礼，而是由于持有与我不同看法的必要性。因为，在所有对某人\Emph{心意}不确定和未知的情况下，我认为最好假设更善意的动机，而不是指责一个未明了的动机。也许他也是出于对我的爱，因为他知道他所写的可能最终会传到我手中；同时他不愿意我在一些他特别确信自己没有错误的观点上犯错。因此，我应该感谢他的好意，尽管我觉得必须反对他的意见。因此，就他持有不正确观点的那些方面而言，在我看来，他应该接受温和的纠正，而不是严厉的指责；尤其是，如果我得到的信息正确，他最近成为了公教徒——这件事是值得祝贺的。因为他已经脱离了此前卷入的多纳提斯派（或更确切地说，罗加提亚派）的分裂和错误；如果他按照应有的方式理解公教的真理，我们真可以为他的归依而欢喜。\par

% structure: p0008 source=h2 target=subsection reason=relative-heading-rank; base=h2

\subsection{第三章 \Person{味增爵}的雄辩，其危险与可容忍性。}

因为他有一种雄辩，能够解释他所想的。因此，我们必须相应地对待他；并希望他能持有正确的看法，并且不让无用的东西成为欲望的对象；这样他就不会把他用雄辩表达的任何东西当作真理提出来。但在他的直言不讳中，他有很多要纠正和删减冗词赘语的地方。实际上，他的这种特点已经冒犯了您，您是一个严肃的人，正如您自己的作品所表明的。这个缺点要么容易纠正，要么，如果轻浮的人喜欢它，而严肃的人容忍它，也无损于他们的信仰。因为我们中间已经有言辞浮夸但信仰纯正的人。因此，我们不必绝望，即使在他身上，这种品质（即使它持久，也是可以忍受的）也可能被调和、净化——事实上，可能被扩展或回归到完整而坚实的标准；尤其是他据说还年轻，所以勤奋可以弥补他缺乏经验所拥有的任何缺陷，而成熟的年龄可以消化他那粗疏的饶舌所难以消化的事物。令人烦恼、危险和有害的是，当愚蠢因雄辩所得的称赞而显得光鲜，当毒药从珍贵的杯中喝下时。\par

% structure: p0010 source=h2 target=subsection reason=relative-heading-rank; base=h2

\subsection{第四章 \Person{味增爵·维克多}书卷中包含的错误。他说灵魂来自天主，但既非由无中生有，亦非由任何受造物所造。}

我现在要继续指出在他的争辩性论述中应当主要避免哪些事。他说灵魂确实由天主所造，但不是天主的一部分或天主的本性——这完全是真实的陈述。然而，当他拒绝承认灵魂是由无中造成的，又不提它是从别的受造物造成；并且如此地使天主成为它的作者，以至于必须认为天主创造它，既不是出于任何不存在的东西（即出于无），也不是出于任何存在而非天主的东西，而是出于祂自己时；他几乎没有意识到，在他的思想回转中，他已回到他自以为已避免的立场，即灵魂无非是天主的本性；因而确实有某物由同一个天主从祂的本性造成，为制造这个，祂用来制造的材料就是祂自己；并且这样，天主的本性是可变的，而因变得更坏，天主自己的本性本身被同一个天主所谴责！这一切离你明智的信德所当相信的多么遥远，离公教信徒的心多么疏远，多么应当避免，你可以轻易看出。因为灵魂要么是由嘘气造成的，要么天主的嘘气如此进入它，以致它不是由祂造成，而是由祂从无中造成。确实不像人的情况：当人呼气时，他不能从无中形成气息，而是将他从空气中吸入的气息归还给空气。我们或许可以如此设想：某些空气围绕着神圣者，祂在吸气时吸入了一点，在呼气时又呼出，当祂向人的脸上嘘气时，就这样为他形成了一个灵魂。若是如此，祂所呼出的就不能出于祂自己，而必须出于周围的空气质料。然而，我们绝不能说全能者不能从无中造出生命的气息，使人类成为有生命的灵魂；并且把自己逼入如此困境，以致要么认为已有某物而非祂自己存在，祂从中形成气息，要么认为祂从自己形成了我们所见是受制于变化的东西。现在，凡出于祂的，必然与祂同一本性，因此是不变的；但灵魂（如众人所承认）是可变的。因此它不是出于祂，因为它不像祂那样不变。然而，如果它不是由任何其他东西造成的，则无疑是由无造成的——但由祂自己。\par

% structure: p0012 source=h2 target=subsection reason=relative-heading-rank; base=h2

\subsection{第五章 [V.]——维克多的另一个错误：灵魂是有形体的。}

但关于他的主张\Quote{灵魂不是神体，而是身体}，他还能说明什么呢，无非是我们不是由灵魂和身体组成，而是由两个或甚至三个身体组成？因为他说我们由神体、灵魂和身体组成，并断言三者都是身体；这随之而来，他认为我们由三个身体组成。这个结论多么荒谬，我想与其是向你，不如是向他证明。但这对于还没有发现存在着某种非有形的事物，却可能具有某种形似的相似性的人来说，并非不可容忍的错误。\par

% structure: p0014 source=h2 target=subsection reason=relative-heading-rank; base=h2

\subsection{第六章 [VI.]——另一个错误：出自他的第二卷，论灵魂因肉体而应受玷污。}

但他在他的第二卷中几乎令人无法容忍，当他试图解决一个关于原罪极其困难的问题时，即如果灵魂不是由父母遗传而是由天主重新嘘入人内，那么原罪如何属于身体和灵魂。为了解释这个棘手而深奥的问题，他这样表达他的观点：\Quote{通过肉体，灵魂适当地恢复它似乎通过肉体逐渐失去的原初状态，为能借此开始被那使它应受玷污的肉体所重生。}你观察此人，如此大胆地承担超过他能力的事，竟跌入如此峭壁，以致说灵魂应受肉体玷污；尽管他绝不能说明在穿上肉身之前，灵魂从何给自己招来这功过。因为如果它首先从肉体得到了罪的功过，就请告诉我们（如果他能的话），在犯罪之前，它从何得到被肉体玷污的功过。因为将这个功过标识它进入有罪的肉体而被其玷污，它当然要么从自身得来，要么（这更令我们反感）从天主得来。它在穿上肉身之前，肯定不能从那个肉体得到使它进入肉体而被玷污的恶功。现在，如果它从自身得到恶功，它如何得到，既然它在接受肉身之前没有犯罪？但如果声称它从天主得到恶功，那么我问，谁能听这样的亵渎？谁能忍受？谁能允许它被任意宣扬？因为这里产生的问题，请记住，不是审判灵魂在降生后应被定罪的恶功是什么，而是它在肉体之前被判定穿上肉体以被它玷污的恶功是什么？请他向我们解释这个，如果他能的话，既然他敢说灵魂应受肉体的玷污。\par

% structure: p0016 source=h2 target=subsection reason=relative-heading-rank; base=h2

\subsection{第七章 [VII.]——维克多将自己缠入一个极其困难的问题。天主的预知并非罪的原因。}

在另一处，他就同一问题（他本人曾深陷其中）进行解释时，以某些反对者的口吻说：\Quote{他们质问：既然灵魂因与肉身的结合，本来不可能有罪的那一位开始有罪，天主为何竟将如此不义的惩罚加于灵魂，宁愿将它贬入肉身？}如今，在这问题如暗礁般的海洋中，他本当警惕失事，切勿冒险陷入那些无法指望安然渡过的险境，唯有退却——也就是说，唯有悔改——才能安全。他试图藉天主的预知自救，却是徒劳。因为天主的预知只不过预先标出祂决意医治的罪人。倘若祂从罪恶中释放那些祂自己曾使之与无罪纯洁的灵魂纠缠于罪中的灵魂，那么祂所医治的是祂自己加给我们的创伤，而非祂在我们身上发现的创伤。然而，愿天主禁止，愿我们断然远离这种说法：当天主藉重生之洗洗涤婴孩的灵魂时，祂所纠正的是祂自己为它们造作的邪恶，即祂将无罪之灵魂与有罪之肉体混合，使它们被原罪沾染。至于这位诽谤者所宣称应当受肉体玷污的灵魂，他完全无法告诉我们，在它们与肉体结合之前，竟应得如此巨大的邪恶。\par

% structure: p0018 source=h2 target=subsection reason=relative-heading-rank; base=h2

\subsection{第八章 [VIII]——维克托的错误意见：灵魂应得成为有罪的}

他徒然以为自己能从天主预知解决此问题，却继续挣扎着说：\Quote{倘若那本来不可能有罪的灵魂应得成为有罪的，但它并未停留在罪中，因为，正如在基督身上所预表的，它本不必在罪中，正如它不可能在罪中一样。}那么他说\Quote{本来不可能有罪的}或\Quote{不可能在罪中}是什么意思呢？我想，除非是指它若未进入肉身？因为，若不从父母传来，它当然不能藉原罪而有罪，或丝毫涉及原罪，除了藉着肉体。我们看到它藉恩宠从罪中获得解放，但我们看不到它如何应得被卷入罪中。那么，他说的\Quote{倘若灵魂应得成为有罪的，但它并未停留在罪中}是什么意思？因为若我问他为何它未停留在罪中，他会很恰当地回答：因为基督的恩宠将它从中解救出来。既然他告诉我们一个婴孩的灵魂如何从罪性中解放，就让他进一步告诉我们它如何应得成为有罪的。\par

% structure: p0020 source=h2 target=subsection reason=relative-heading-rank; base=h2

\subsection{第九章——维克托完全无法解释无罪灵魂如何应得成为有罪的}

他在引言中说自己遭遇了什么？因为他在提出那个问题之前，并作为引言，他断言：\Quote{还有些辱骂性的言论隐藏在那些诋毁我们的人抱怨的怨言之下，我们如同在飓风中颠簸，一次又一次被猛撞在巨大的岩石上。}若我用这种口气谈论他，他大概会生气。这些话是他自己的；他先说出这些话，然后提出他的问题，仿佛向我们展示他撞击并失事的那岩石本身。因为他被带到如此地步，被冲向、漂向并撞上如此可怕的礁石，若不退却——即纠正他所说的话——便绝无可能逃脱；因为他无法指出灵魂因何种功德成为有罪的；尽管他不怕说，在灵魂有任何自己的罪之前，它已应得成为有罪的。然而，谁能在未犯任何罪的情况下，应得如此巨大的惩罚，以致在离开母胎之前就怀在别人的罪中，而后不再脱离罪？但从这惩罚中，天主的白白恩宠解放了那些在基督内重生的婴孩的灵魂，没有任何他们先前的功德——否则恩宠便不再是恩宠。\BibleRef{罗 11:6}至于这位极其聪明、且因我们谨慎（若非博学，至少审慎）而不悦的深奥智慧人物，让他告诉我们（若能的话），那使灵魂陷入如此惩罚（恩宠白白地从中解救，没有任何功德）的功德是什么。让他说话，若可能，让他以某种理由捍卫他的断言。我本不会如此要求他，若他自己没有宣称灵魂应得成为有罪的。让他告诉我们那功德是什么——是好功德还是坏功德？若是好的，善功如何导致恶？若是坏的，在犯任何罪之前，恶功从何而来？我亦需指出，若有善功，则灵魂的解放便不是白白恩宠，而是因先前的功德应得的，如此\Quote{恩宠便不再是恩宠}。然而，若有恶功，我就问它是什么。难道灵魂进入肉身是真实的，除非那毫无罪愆的那一位亲自差遣它，它就不会如此进入？因此，除非越陷越深，他将永远无法成立这个观点，即他断言灵魂应得成为有罪的。至于那些婴孩，在其圣洗中原罪被洗去，他总算找到了话可说（以某种方式）——即说，被卷入另一个人的罪中不可能对他们有害，因为他们在天主的预知中被预定得永生。这或许可以有相当好的意义，如果他未曾陷入自己的那个公式中，即他断言灵魂应得成为有罪的：他只有撤回自己的话并后悔曾说出它们，才能从这困境中脱身。\par

% structure: p0022 source=h2 target=subsection reason=relative-heading-rank; base=h2

\subsection{第十章 [IX]——维克托的另一错误：未受圣洗而死的婴孩可能获得天国。另一错误：必须为未受圣洗而死去的婴孩献上基督身体的祭献。}

但是，当他试图就那些因死亡而未能首先在基督内领受圣洗的婴孩作出恭敬的回答时，他竟大胆地许诺他们不仅得享乐园，而且得享天国——因为他找不到其他办法来避免这样说：天主将那些没有任何先前罪过却引入有罪肉身的无辜灵魂定为永死。然而，他多少看到了自己说出的话的邪恶，即暗示婴孩的灵魂无需基督的任何恩宠即可被救赎得享永生和天国，并且在他们的情形中，原罪可以在没有基督的圣洗（其中实现了罪的赦免）的情况下被取消：看到这一切，以及他在自己的沉船之海中陷入何等的深渊，他说：\Quote{我认为，确实必须由圣洁的司铎不断为他们奉献祭品和牺牲。} 你可以在此看到另一个危险，他除了后悔并收回自己的话之外，永远无法逃脱。因为谁能将基督的身体献给任何不属于基督肢体的人呢？再者，自从他说：\BibleQuote{人除非由水和圣神而生，不能进天主的国；}\BibleRef{若 3:5} 又说：\BibleQuote{谁为我的缘故丧失自己的性命，必得着性命。}\BibleRef{玛 10:39} 以后，没有一个人能成为基督的肢体，除非是通过在基督内的圣洗，或是为基督而死。\par

% structure: p0024 source=h2 target=subsection reason=relative-heading-rank; base=h2

\subsection{第十一章——为基督殉道补足圣洗的位置。与基督同钉的强盗的信德被视为殉道，因而代替圣洗。}

因此，那个在十字架之前并非主的追随者、却在十字架上成了主的宣认者的强盗——他的案例有时被用来对圣洗圣事提出异议或尝试攻击——被圣西彼廉列在那些以自己的血领受圣洗的殉道者之中，如同在激烈迫害时期许多未领受圣洗者所经历的那样。因为他宣认了被钉的主，那位知道如何衡量并评估此类证据的主赋予它如此大的份量和如此有效的价值，就好像他是为主被钉一样。那时，他的信德在十字架上兴盛，而门徒们的信德却衰败了，而且如果没有因那位死亡恐怖曾使其低垂者（指基督）的复活而重新绽放，就会无可挽回。他们在他临死时对他绝望——而他在与他同死时对他怀有希望；他们逃离生命之源——而他却向自己的受难同伴祈祷；他们哀悼好像一个人的死亡——而他相信死后他要成为君王；他们离弃自己救恩的担保者——而他尊敬自己十字架的同伴。在他身上发现了殉道者的全部衡量，他在那些注定成为殉道者的人跌倒之时相信了基督。这一切，确实在主的眼中是显明的，主立即将如此巨大的福乐赐予这位虽未领受圣洗，却如同在殉道的血中洗净的人。但我们自己，难道不能反思：一个在临死时向主祈求生命的人，在他活着时，会是怀着何等大的信德、何等的望德、何等的爱德为基督而死？除此之外，还有一个并非难以置信的报道：那个当时相信的强盗，当他悬挂在被钉的主身旁时，被从救主肋旁伤口流出的水所洒，如同领受了最神圣的圣洗。我不提这样一个事实：没有人能证明——因为我们没有人知道他是否在定罪之前领受过圣洗。然而，让每个人都按照自己喜欢的理解来看待这事吧；只是不要让任何关于圣洗的规则影响救主本人的诫命，从这强盗的例子中引出来；也不要让任何人为未领受圣洗的婴孩的情形，在永罚与天国之间，许诺一个任意选择的中等安逸和幸福之所。因为白拉奇异端恰恰许诺了这一点：他既不为婴孩惧怕永罚，因为他认为他们没有原罪，也不给他们天国的希望，因为他们没有领受圣洗圣事。然而，至于这个人（指本文批判的对象），虽然他承认婴孩带有原罪，却大胆地许诺他们即使没有圣洗也能得享天国。这甚至是白拉奇派也不敢做的，尽管他们断言婴孩完全无罪。那么，看他把自己缠绕在何等武断意见的罗网中，除非他后悔写下这些观点。\par

% structure: p0026 source=h2 target=subsection reason=relative-heading-rank; base=h2

\subsection{第十二章——殉道者圣伯尔伯都亚的兄弟狄诺克拉特，据说因圣人的祈祷而从受罚的状态中被解救出来。}

然而，关于圣伯尔伯都亚的兄弟狄诺克拉特，在正典圣经中没有记载；圣人本人，或无论谁写的记述，也没有说那个在七岁时去世的男孩是没有领受圣洗而死的；人们相信，在她的殉道即将来临时，她的祈祷被有效地垂听，使他从丧亡者的惩罚中被移去得享安息。现在，那个年龄的男孩既能说谎，也能说实话，既能宣认也能否认。因此，当他们领受圣洗时，他们背诵信经，并代替自己回答在考查中向他们提出的问题。那么，谁能知道那个男孩在领受圣洗之后，在迫害时期，是否被一个不虔诚的父亲引离基督而陷入偶像崇拜，并因此招致了永罚，而只有因基督的缘故，因他姐姐在临终时的祈祷，他才得以解脱？\par

% structure: p0028 source=h2 target=subsection reason=relative-heading-rank; base=h2

\subsection{第十三章——基督体血的祭献对未领受圣洗的人无效，也不能为大多数未领受圣洗而死的人奉献。}

但即使向此人让步（这绝不可能安全地允许于公教信仰和教会之规则），即基督圣体圣血的祭献可为任何年龄未受洗者献上，仿佛他们可因亲友这种敬虔的帮助而达至天国：那么对于我们的异议，即成千上万生于不敬虔父母之婴孩，从未因天主或人的任何仁慈落入敬虔亲友之手，且在如此年幼之际未得圣洗便离开那悲惨生命——他将有何言以对？让他告诉我们，若他能的话，那些灵魂是如何本该成为有罪的，以致此后绝不可能从罪中得释放。因为若我问他，他们若未受洗为何该受惩罚，他将正确回答我：因原罪。若我再追问，他们从何处得来原罪，他将回答：当然来自有罪的肉身。若我继续问，他们为何该受罚得一个有罪的肉身——他们在进入肉身之前未行任何恶事——且如此受罚以致经历他人之罪的污染，以致既无圣洗在他们生于罪中时重生他们，也无祭献在他们污秽中赎罪？让他找些话来回答吧！因为正是从这样的环境和这样的父母中，这些婴孩已经出生或仍在出生，以致他们不可能得到这种帮助。在此，所有论据都缺失了。我们的问题并非：灵魂在与有罪肉身结合后为何该受罚？而是问：灵魂在与有罪肉身结合之前毫无罪恶，为何该受罚去经历这种与有罪肉身的结合？他没有余地可说：\Quote{他们暂时分享他人之罪的污染并无损害，因为在天主的预知中，救赎已为他们预备。}因为我们现在谈论的是那些在受洗之前离开身体、无可救赎帮助的人。他也没有适当理由说：\Quote{圣洗未能洁净的灵魂，为她们所献的许多祭献将洁净她们。天主预知此事，并愿意他们暂时卷入他人之罪而不受永罚，且怀有永福的希望。}因为我们现在谈论的是那些生于不敬虔之人、出于不敬虔父母，绝不可能获得如此辩护和帮助的人。即使这些可被应用，也肯定不能有益于任何未受洗者；正如他所引用的\BibleBook{玛加伯}{书}中的祭献，对于为之献祭的有罪死者无用，因为他们未受割损。\BibleRef{加下 12:43}\par

% structure: p0030 source=h2 target=subsection reason=relative-heading-rank; base=h2

\subsection{第十四章——\Person{维克托}的两难：他必须说要么所有婴孩都得救，要么天主杀死无辜者。}

那么，让他若有能力回答这个问题：为何灵魂没有任何——无论是原罪还是本罪——应受罚去经历他人原罪，以致不能从中得释放；让他看看他将在两个选项中选哪一个：要么说那些未得圣洗便离世、且未为主圣体献祭的婴孩的灵魂，从原罪之链中得释放——尽管宗徒教导说\Quote{从一人而来，众人都被定罪}\BibleRef{罗 5:16}——当然，所有那些恩宠未及帮助的人，以便通过一人众人都得以进入救赎。要么说那些没有任何——自己的罪恶还是原罪——全然无辜、单纯、纯洁的灵魂，被公义的天主惩罚于永罚中，当祂亲自将他们插入有罪肉身而无任何救赎时。\par

% structure: p0032 source=h2 target=subsection reason=relative-heading-rank; base=h2

\subsection{第十五章[第十二节]——天主不按任何人若生命延长可能做的事审判他，而只按他实际所行的行为。}

就我而言，我确实断言，这些备选情况都不应被接受，也不应接受那第三种意见，即认为灵魂在肉体之前在其他某种状态中犯了罪，因此应被判定归于肉体；因为宗徒已经最明确地指出：\BibleQuote{孩子还没有出生，没有行善或作恶。}\BibleRef{罗 9:11}所以很明显，婴儿除了需要罪之赦免的原罪外，不可能有其他罪。此外，第四种立场，即那些未领受圣洗而死去的婴儿的灵魂，被公义的天主放逐并判定归入有罪的肉体，因为祂预知如果年长到能使用自由意志，他们会过邪恶的生活。但即使是他，尽管陷于这样的困惑中，也不敢如此断言。相反，他简短但明确地宣告，反对这虚妄的意见，说：\Quote{如果天主愿意审判任何尚未出生、未曾凭自己的自由意志行事的人，天主就是不义的。}这是他在处理一个反对某些人的问题时的回答，那些人问：为什么天主造人，既然祂预知这人不会好？如果祂因为预知这人不会好而不愿创造他，那祂就是在审判一个未出生的人。毫无疑问，如这个人自己所认为的那样，恰当的做法是全能者应在人完成其行为时审判他，而不是根据可能预见的行为，也不是根据可能允许将来发生的行为。因为如果一个人若活着将会犯的罪在他死后被定罪，即使这些罪并未犯下，那么当他被带走以使邪恶不能改变他的心意时，对他并无益处；因为对他的审判将根据可能在他身上发展的邪恶，而非根据实际在他身上找到的正直。而且，任何在领受圣洗后死去的人都不可能安全，因为即使在领受圣洗之后，人们——我不说会以某种方式犯罪——实际上甚至会犯背教。那么如何？假设一个在领受圣洗后被带走的人，如果他还活着，将会成为背教者，我们是否认为他被带走并从心意被邪恶改变的痛苦中得救对他也没有益处？我们是否要想象，由于天主的预知，他将被当作背教者，而不是基督的忠实肢体受审判？确实，如果罪不是按人已犯或人已图谋来惩罚，而是按全能者预知和将要发生的来惩罚，那么让第一对夫妇在堕落之前就被逐出乐园，从而如此神圣有福之地就不会发生罪，那该有多好！此外，当预知的事不会发生时，预知本身完全失效，又该怎么说？的确，怎能正确称其为对将要之事的预知，而事实上它不会发生？罪如何被惩罚——这些罪并不存在，也就是说，在取肉体之前未被犯下，因为生命尚未开始；在取肉体之后也未被犯下，因为死亡阻止了？\par

% structure: p0034 source=h2 target=subsection reason=relative-heading-rank; base=h2

\subsection{第16章 [XIII.]——维护灵魂不是由传生而来的意见的困难。}

那么，这种解决问题的方法——即灵魂被送入肉体直到它应从肉体中被释放——考虑到所假设的是婴儿的灵魂，其年龄尚不足以使意志变得自由——并不能揭示为什么它在未领受圣洗的情况下应受惩罚的原因，除非是原罪的原因。由于这个罪，我们不否认灵魂被公义地定罪，因为天主的公义法律为罪规定了惩罚。但我们要问，为什么灵魂被置于这种有罪的状态，如果它不是从那个在人类始祖中犯了罪的原初灵魂衍生而来的？因此，如果天主不谴责无辜者——如果祂不使那些祂视为无辜的人有罪——并且除了基督教会中基督的圣洗之外，没有任何事物能将灵魂从原罪或本罪中解放出来——并且罪在犯罪之前，更不用说从未犯过的罪，不能被任何公义的法律定罪：那么，这位作者就不能提出这四种情况中的任何一种；他必须（如果他能够）就婴儿的灵魂（他们未领受圣洗就离开生命，被送入惩罚）解释，他们从未犯过罪，是因什么功绩被交付于有罪的肉体，在那里找到确保他们公义定罪的罪。此外，如果他回避这四种健全教义所谴责的情况——也就是说，如果他没有勇气主张：灵魂即使没有罪，也被天主造成有罪；或者他们未通过基督的圣事就摆脱了内在的原罪；或者他们在被送入肉体之前在其他状态中犯了罪；或者他们从未犯过的罪在他们身上被定罪——如果，我说，他没有勇气告诉我们这些事，因为它们确实不值得提及，反而应当肯定婴儿不继承原罪，并且如果他们未领受重生圣事就离开，也没有理由被定罪，那么他将无疑，为定他自己的罪，陷入\Person{白拉奇}的可咒异端。为避免此，他与我一同对灵魂的起源持犹豫态度，不敢肯定他无法以人的理性理解也无力以神圣权威辩护的事，岂不是好得多？这样，他就不必在害怕承认自己无知时，被迫说出蠢话。\par

% structure: p0036 source=h2 target=subsection reason=relative-heading-rank; base=h2

\subsection{第17章 [XIV.]——他指出\Person{维克多}引用的圣经章节并不证明灵魂是由天主造成的，以致不由传生而来：第一段引文。}

或许他会说，他的见解有神圣权威的支持，因为他以为自己已用圣经经文证明：灵魂不是天主通过\OriginalTerm{繁殖}{propagation}的方式造的，而是通过独立的创造行动，重新吹入每个人内的。如果他能够证明这一点，我就承认，我从他那里学到了我极其渴望探求的道理。但他必定要去寻找其他的论据，而他或许找不到，因为他到目前为止所引用的经文并没有证明他的论点。因为他所有用于该主题的引文，在某种程度上无疑是恰当的，但对他所提出的关于灵魂起源的观点，只提供了可疑的证明。因为天主确实赐给了人气息和神，正如先知所证实的：\BibleQuote{上主这样说，祂创造了诸天，奠定了大地，并其中的万物；祂赐给地上的人民气息，给行走其上的人们精神。} \BibleRef{依 42:5} 他希望这段经文按他所维护的意义来理解；因此，\BibleQuote{祂赐给人民气息}这句话，可以被理解为：祂为人类创造灵魂，不是通过\OriginalTerm{繁殖}{propagation}，而是每次都以\OriginalTerm{吹气}{insufflation}的方式创造新的灵魂。那么，让他大胆地按此说法坚持：天主不赐给我们肉身，因为我们的肉身源自父母。同样，在宗徒所引的例子中，\BibleQuote{天主随祂的意愿给它一个身体} \BibleRef{格前 15:38}，如果他有胆量，让他否认麦粒出自麦粒，青草出自青草，从种子生出各从其类。如果他不敢否认这一点，那么他怎能知道\BibleQuote{祂赐给人民气息}这句话的含义呢？——是从父母遗传而来，还是重新吹入每个人内？\par

% structure: p0038 source=h2 target=subsection reason=relative-heading-rank; base=h2

\subsection{第十八章——\Quote{气息}有时意指圣神。}

此外，他如何知道，句子\BibleQuote{祂赐给地上的人民气息，给行走其上的人们精神}中的意念重复，不可以理解为两个表达指同一件事，即不是指人类本性赖以生活的生命或神，而是指圣神？因为如果\Quote{气息}不能表示圣神，主就不会在复活后向门徒\Quote{嘘了一口气}说：\BibleQuote{你们领受圣神吧！} \BibleRef{若 20:22} 在宗徒大事录中也不会这样记载：\BibleQuote{忽然从天上来了一阵响声，好像暴风\OriginalTerm{气息}{breath}猛刮，充满了他们所在的全屋；有些散开好像火的舌头，停留在他们每人头上，众人都充满了圣神。} \BibleRef{宗 2:2} 现在，假设先知所预言的就是这事，他说：\BibleQuote{祂赐给地上的人民气息}；然后，作为对他所说的\Quote{\OriginalTerm{气息}{breath}}的解释，他接着说：\Quote{并给行走其上的人们精神。} 肯定地，当他们一齐充满了圣神时，这个预言最清楚地应验了。然而，如果\Quote{人民}一词尚不适用于当时聚集在一处的120人，那么当信徒数目达到四五千，他们受洗时领受了圣神（\BibleRef{宗 4:31}），难道还有人怀疑那时领受圣神的人就是\Quote{人民}，甚至\Quote{行走在地上的人}吗？因为那属于人的本性的神，无论是通过\OriginalTerm{繁殖}{propagation}赐予，还是作为新东西吹入个人（我不决定这两种方式中哪一种应当被肯定，至少在两者之一能清楚确定之前），并不是在人们\Quote{行走在地上}时赐予，而是当他们还关闭在母腹中时就已赐予。\BibleQuote{因此，祂赐给地上的人民气息，并给行走其上的人们精神}，是在许多人一同成为信徒、一同充满圣神时发生的。祂将圣神赐给祂的子民，虽不是同时给所有人，而是给每人各按其时，直到祂子民的全体数目，通过离世和入世而被满足。因此，在这段圣经中，\Emph{气息}不是一样东西，\Emph{神}也不是另一样东西；而是同一意念的重复。正如\Quote{坐在天上的}不是一位，\Quote{上主}也不是另一位；\Quote{发笑}和\Quote{嗤笑}也不是两回事；在以下经文中只是同一意义的重复：\BibleQuote{坐在天上的必发笑，上主必嗤笑他们。} 同样，在经文中：\BibleQuote{我必将外邦人赐你为\OriginalTerm{嗣业}{inheritance}，地极为你\OriginalTerm{产业}{possession}}，当然不意味着\Quote{嗣业}是一回事，\Quote{产业}是另一回事；也不意味着\Quote{外邦人}是一回事，\Quote{地极}是另一回事；只是同一事物的重复。他若细细思想所读的，必会在圣经中发现无数这类表达。\par

% structure: p0040 source=h2 target=subsection reason=relative-heading-rank; base=h2

\subsection{第十九章——圣经中\Quote{气息}的意义。}

然而，希腊文版本中使用的术语\OriginalTerm{气}{\GreekText{πνοή}}，在拉丁文中有多种译法：时而译作\OriginalTerm{气息}{flatus}，即气息；时而译作\OriginalTerm{神}{spiritus}，即神灵；时而译作\OriginalTerm{灵感}{inspiratio}，即灵感。该术语出现在我们正在研究的段落中：\BibleQuote{祂将气息赐给其上的人民}，其中表示气息的词是\OriginalTerm{气}{\GreekText{πνοή}}。同一词也出现在人被赋予生命时的叙述中：\BibleQuote{天主向他的脸吹了生命的气息}。此外，在圣咏中同样出现该词：\BibleQuote{凡有气息的，都要讚美主}。同一词也出现在\BibleBook{约伯}{传}中：\BibleQuote{全能者的灵感教导人}。译者拒绝了flatus（气息）而采用\OriginalTerm{灵感}{adspiratio}，尽管他面对的正是我们所考察的先知书中的希腊文\OriginalTerm{气}{\GreekText{πνοή}}。我想，我们几乎可以肯定，在\BibleBook{约伯}{传}的这段经文中，圣神被指代。讨论的问题是智慧从何而来：\BibleQuote{它不来自年岁的多数；但圣神在凡人内，全能者的灵感教导人}。\BibleRef{约 32:7}通过这种重复，可以充分理解他说的不是人自己的神。他想表明人的智慧不是来自自己，所以用双重表达来解释：\BibleQuote{全能者的灵感教导人}。同样，同书另一处说：\BibleQuote{我嘴唇的领悟将默想纯洁。圣神塑造了我，全能者的气息教导我}。这里，他称为\OriginalTerm{灵感}{adspiratio}或\Quote{灵感}的，在希腊文中也是\OriginalTerm{气}{\GreekText{πνοή}}，与先知书中译为\OriginalTerm{气息}{flatus}的相同。因此，虽然断言\BibleQuote{祂将气息赐给其上的人民，将圣神赐给在上行走的人}指人的灵魂或神是鲁莽的——尽管更可信地理解为指圣神——但我还是问：凭什么胆敢断定先知说这些话是要暗示我们的本性具有生命力的灵魂或神不是由天主通过生育过程给予的？当然，如果先知说得很明白：\BibleQuote{祂将灵魂赐给地上的人民}，那么仍然要问：是否天主自己从先前的世代赐给灵魂，正如祂赐给身体来自先前的物质，并且不仅赐给人或牲畜，也赐给谷物的种子或任何其他身体，随祂喜欢？还是祂像第一个人那样，通过吹气将灵魂作为一种新恩赐给每个人？\par

% structure: p0042 source=h2 target=subsection reason=relative-heading-rank; base=h2

\subsection{第二十章——理解该段的其他方式}

也有些人对先知的话\BibleQuote{祂将气息赐给其上的人民}做出理解，认为这里的\Quote{气息}（\OriginalTerm{气息}{flatus}）等同于\Quote{灵魂}（\OriginalTerm{灵魂}{anima}）；而下文\BibleQuote{将圣神赐给在上行走的人}则指圣神；他们认为先知遵循了宗徒所说的次序：\BibleQuote{先有的不是属神的，而是属生灵的；后有的是属神的。}\BibleRef{格前 15:46}从这种观点出发，无疑可以形成优雅的解释，与宗徒的思想一致。\Quote{在上行走的人}在拉丁文中是\Quote{\OriginalTerm{践踏它}{calcantibus eam}}；字面意义是\Quote{践踏它}，暗示轻蔑。因为领受圣神的人在爱慕天上事物时轻视地上的事物。这些观点都不违反信仰，无论将这两个术语\OriginalTerm{气息}{flatus}和\OriginalTerm{神}{spiritus}视为属于人性，还是二者都指圣神，或者其中\Quote{气息}指灵魂，\Quote{神}指圣神。然而，如果这里指的是人的灵魂和神，既然它无疑应被视为天主的恩赐，那么我们必须进一步探究：天主如何赐予这恩赐？是通过生育过程，如同祂赐给我们身体肢体那样？还是通过吹气赐给每个人，不是通过生育，而是始终为新的创造？这些问题并不像此人所述的那样模棱两可；但我们希望用最可靠的圣经依据来辩护它们。\par

% structure: p0044 source=h2 target=subsection reason=relative-heading-rank; base=h2

\subsection{第二十一章——维克托引用的第二段经文}

以同样原则处理天主说的：\BibleQuote{因为我的圣神要从我出来，我创造了每一个气息。}前半句\BibleQuote{我的圣神要从我出来}应指圣神，关于祂，救主同样说：\BibleQuote{祂从父出来}\BibleRef{若 15:26}。但另一分句\BibleQuote{我创造了每一个气息}无疑指每个灵魂。然而，天主也创造人的整个身体；且如无人怀疑，祂通过生育过程造人的身体；因此，关于灵魂（它显然是天主的工作），自然仍可探究：祂是像造身体那样通过生育，还是像造第一个灵魂那样通过吹气而造灵魂？\par

% structure: p0046 source=h2 target=subsection reason=relative-heading-rank; base=h2

\subsection{第二十二章——维克托引用的第三段经文}

他接着赐给我们第三段引文，其中写道：\Quote{他在人内形成人的心神。}好像有人否认这一点！不；我们所有的问题都在于形成的方式。现在让我们取肉体的眼睛，问：除了天主，谁形成它？我想，祂不是在外部形成它，而是在它自身内形成，然而，确实是通过传生。既然祂也形成\Quote{在人内的人的心神，}问题仍然存在：它是每次都通过新的嘘气而来，还是通过传生而来。\par

% structure: p0048 source=h2 target=subsection reason=relative-heading-rank; base=h2

\subsection{第23章——他的第四段引文。}

我们读过所有关于玛加伯青年母亲的事迹，当孩子们受苦时，她在德行上比在他们出生时更为多产；她这样劝勉他们持守坚忍：\BibleQuote{我的儿子们，我不知道你们怎样进入我的胎中。因为不是我给了你们心神与灵魂，也不是我形成了你们每个人的肢体；而是天主，祂也创造了世界和其中万物；祂更形成了人类的世代，并搜寻一切行为；祂将因祂的大慈大悲把你们的心神与灵魂归还给你们。}\BibleRef{加下 7:22}这一切我们都知道；但我们不明白这如何支持此人的断言。因为哪个基督徒会否认天主赐给人灵魂与心神呢？但同样，我想他不能否认天主赐给人舌头、耳朵、手、脚和一切身体感官，以及所有肢体的形状和本性。因为他怎能否认这些都是天主的恩赐，除非他忘记了自己是基督徒？然而，既然显而易见，这些都是由祂所造，并通过传生赐予人；那么也必须提出这个问题：人的心神和灵魂是藉什么方式由祂形成？藉什么效率赐予人——来自父母，还是来自虚无，或者（如这个人所主张，其意义我们必须极力防范）来自天主嘘气的一种已有本性，不是从虚无中受造，而是从祂自身而来？\par

% structure: p0050 source=h2 target=subsection reason=relative-heading-rank; base=h2

\subsection{第24章 [XV]——灵魂是否由自然遗传（传生）而来，他所引用的经文未能证明这一点。}

既然他所引用的圣经章节丝毫没有表明他试图证明的内容（因为它们确实对当前直接问题毫无阐述），他这些话究竟是什么意思：\Quote{我们坚定地主张，灵魂来自天主的嘘气，而非自然生育，因为它是天主所赐？}好像，确实，身体能来自另一位，而不是来自创造它的那一位，\BibleQuote{万有出于祂，万有藉着祂，万有在祂内；}\BibleRef{罗 11:36}不是因为它们是祂的本性，而是因为它们是祂的化工。\Quote{也不是来自虚无，}他说，\Quote{因为它出自天主。}是否如此，我们必须说，不是此处要研讨的问题。同时，我们毫不犹豫地断言，他所提出的命题——灵魂既非来自遗传，也非来自虚无——肯定不真实：我说，我们毫不怀疑地断言这不真实。因为二者必居其一：如果灵魂不是来自父母的血缘遗传，它就出自虚无。假装它来自天主，以致成为祂本性的部分，简直是亵渎神明的不敬之言。但直至现在，我们恳求并寻找一些关于此点的清晰圣经章节，即灵魂是否不来自父母遗传；但我们不想要他所引用的那些经文，因为它们对当前的问题毫无阐明。\par

% structure: p0052 source=h2 target=subsection reason=relative-heading-rank; base=h2

\subsection{第25章——正如母亲不知道孩子在她腹中从何而来，我们也不知道灵魂从何而来。}

我多么希望，在如此深奥的问题上，只要他不知道该说什么，他就效法玛加伯青年的母亲！虽然她很清楚自己是从丈夫怀中受孕，而且万有的创造者为她创造了他们，无论是在身体，还是在灵魂和神魂，然而她却说：\Quote{我不能告诉你们，我的儿子们，你们是怎样进入我胎中的。}好吧，我只希望这个人能告诉我们她所不知道的事！当然，她（在我所提及的几点上）知道他们是如何进入她胎中的，就肉身实体而言，因为她不可能怀疑自己是由丈夫受孕的。她进一步承认——因为这一点她当然也很清楚——是天主赐予了他们灵魂和神魂，也是祂为他们形成了五官和肢体。那么，她所不知道的是什么呢？难道不可能是（我们也同样无法确定的）吗？即天主无疑赐予他们的灵魂和神魂，是从父母传生而来，还是像当初对第一个人那样，作为新创的嘘入？但无论是这个，还是关于人性构造的其他某个细节，她并不知道，她坦率地承认了自己的无知；并没有胆量随意辩护她所不知道的事。这个人也不会对她说，他竟不害羞地对我们说：\Quote{人享尊荣而不醒悟，如同牲畜一般，成为相似。}看，那女子说到自己的儿子时：\Quote{我不能告诉你们怎样进入我胎中。}然而她并没有被比作无知的牲畜。她说：\Quote{我不能告诉}；然后，好像他们问她为什么无知，她接着说：\Quote{因为不是我给了你们神魂和灵魂。}因此，谁赐予了他们那恩赐，就知道祂从何处造出祂所赐的，无论是由传生而来，还是作为新创造的嘘入——这一点（这个人说）我一无所知。\Quote{也不是我形成了你们每个人的五官和肢体。}然而，形成他们的那一位，知道祂是连同灵魂一起形成他们，还是在形成之后才将灵魂赐给他们。她对儿子们进入她胎中的方式，无论是这样还是那样，都毫无概念；她只确信一件事：赐予她所有一切的那一位，必会恢复祂所赐的。但这个人却要选择那女子所不知道的，关于我们本性如此深奥和隐晦的事实；只是，如果她错了，他不会判断她；如果她无知，他不会将她比作无知的牲畜。无论她所不知道的是什么，它当然属于人的本性；然而，任何人因这样的无知都是无可指责的。所以，我也从我这一方，关于我的灵魂说：我不知道它如何进入我的身体；因为不是我把它赐给我自己。赐予我的那一位知道，究竟是从我父亲那里传授给我，还是像对第一个人那样，为我新创造的。但就连我自己也将知道，当祂自己，在祂所定的时候，教导我的时候。但现在，我不知道；我也不像他那样羞于承认自己对不知道的事情的无知。\par

% structure: p0054 source=h2 target=subsection reason=relative-heading-rank; base=h2

\subsection{第二十六章——维克托所引用的第五段圣经}

他说：\Quote{你要学习，因为，你看，宗徒在教导你。}是的，我当然愿意学习，如果宗徒教导的话；因为唯独天主借着宗徒教导。但是，请问，宗徒所教导的是什么？他接着说：\Quote{你看，当他对\Place{雅典}人说话时，他是如何有力地陈明这个真理，说：\InnerQuote{他赐给众人生命和气息。}}好，谁会否认这一点呢？他说：\Quote{但你要明白，宗徒所陈述的是什么：\Emph{他赐给}，而不是\Emph{他赐给了}。他指的是持续而不确定的时间，并没有宣告过去和已完成的时间。如今他毫无间断地所赐的，他总是在赐给；正如赐给的那一位本身是永远存在的。}我正是按照在你寄给我的第二本书中所找到的那样，引用了他的原话。首先，我请你注意，他在竭力肯定自己所一无所知的事情时，已经走到了何等极端。因为他竟敢说，天主毫无间断地，不只在现在，而且直到永远，在人生时赐给灵魂。他说：\Quote{他总是在赐给，正如赐给的那一位本身是永远存在的。}我决不会说我不明白宗徒所说的，因为这已足够清楚。但此人所说的，他自己也应该知道，是与基督徒的信仰相悖的；他应当警惕，不要再继续这样的主张。因为，当然，当死人复活时，将不再有人出生；因此天主不再在任何一个出生时赐予灵魂；但那些他现在连同身体赐给人们的灵魂，他将审判。所以，他并非总是在赐给，尽管现在赐给的那一位是永远存在的。而且，这根本不是从宗徒的措辞\Emph{赐给}（不是\Emph{赐给了}）中可以推导出来的，而这位作者却想从中推论出：天主不是通过传生赐给人灵魂。因为灵魂仍然由他赐给，即使是通过传生；就如身体的各种禀赋，如肢体、感官、形状，实际上整个实体，都是由天主亲自赐给人的，尽管他是通过传生赐下它们。此外，因为主说：\BibleRef{玛 6:30}\BibleQuote{天主这样装饰田野的草，今天还在，明天就被投进炉中}（没有使用过去时\Emph{装饰了}，像他最初造成物质时那样；而是使用了现在时\Emph{装饰}，他确实仍在这样做），我们就因此说，百合花不是从它们自己种类的本源所产生的吗？那么，如果人的灵魂和神气也是同样由天主亲自赐予，只要它被赐予；并且也是由它自己种类的传生而赐予，又当如何？这是一个我既不主张也不反驳的立场。然而，如果必须加以辩护或反驳，我当然建议要用清晰而非模糊的证据来进行。我也不应该因为承认自己至今仍无法决定这个问题而被比作无知的牲畜，反而应该被比作谨慎的人，因为我不鲁莽地教导自己一无所知的事。但我自己不打算以辱骂还辱骂，把此人比作畜生；而是作为儿子警告他，要承认自己实在不知道所不知之事；也不要试图教导自己尚未学会的，以免他该被比作宗徒所提及的那些人，他们\BibleQuote{愿意做法律的教师，却不明白自己所说的和所主张的是什麼。}\BibleRef{弟后 1:7}\par

% structure: p0056 source=h2 target=subsection reason=relative-heading-rank; base=h2

\subsection{第二十七章［十七］——\Person{奥斯定}未敢就灵魂的传生作任何定义。}

因为他为何如此轻忽他所谈论的圣经，竟未注意到，当他读到人是从天主而来时，这不仅仅如他所争辩的，关乎人的灵魂和神魂，也关乎人的肉身？因为宗徒的话说：\BibleQuote{我们是祂的子孙，}\BibleRef{宗 17:28}此人以为这不应指肉身，而仅指灵魂和神魂。实际上，如果我们的肉身并非来自天主，那么圣经所说\BibleQuote{万物都是出于祂、借着祂、并在祂内}\BibleRef{罗 11:36}，就是假的。同样，关于同一宗徒的话：\BibleQuote{因为女人是由男人而出，同样男人也是由女人而生}\BibleRef{格前 11:12}，请他向我们解释，在这过程中他所愿指的是哪一种传生——是灵魂的传生，还是肉身的传生，还是两者兼有？但他不会承认灵魂是藉传生而来；因此，根据他和所有否认灵魂传生的人，宗徒说\BibleQuote{女人是由男人而出，同样男人也是由女人而生}时，所指的只是男女的肉身；女人是从男人被造，为的是以后男人能藉由出生从女人出来。因此，如果宗徒说这话时，并非旨在包括灵魂和神魂，而仅指两性的肉身，那么他为何紧接着说\BibleQuote{但一切都是出于天主}\BibleRef{格前 11:12}？除非肉身也是出于天主。因为他的整句话是：\BibleQuote{女人是由男人而出，同样男人也是由女人而生；但一切都是出于天主。}那么，请我们的辩论者确定这是指什么说的。如果是指人的肉身，那么当然，就连肉身也是出于天主的。那么，为什么每当此人读到圣经中说\Quote{ \Emph{出于天主} }，涉及人时，他坚持这应当理解为不关乎人的肉身，而只关乎人的灵魂和神魂？但如果\BibleQuote{一切都是出于天主}，这句话是指两性的肉身以及他们的灵魂和神魂而言，那么女人在一切方面都出于男人，因为女人从男人而来，男人藉女人而生；但一切都是出于天主。这里所说的\Quote{一切}是什么？除非是他所说的那些：那女人所从出的男人，以及那出于男人的女人，还有那藉女人而生的男人？因为那女人所从出的男人并非藉女人而生；只有那后来由男人藉女人所生的人，正如现在人所生的那样。因此，如果宗徒在引用我们从他那里引用的这句话时是指人的肉身，那么两性之人的肉身无疑是出于天主的。此外，如果他坚持认为人身上没有任何东西是出于天主，除了他们的灵魂和神魂，那么当然，女人即使在灵魂和神魂方面也是出于男人；这样，那些反对灵魂传生的人便毫无立足之地。但如果他想要这样区分，说女人在肉身方面出于男人，而在灵魂和神魂方面则出于天主，那么宗徒所说\BibleQuote{一切都是出于天主}，如何为真，如果女人的肉身是出于男人，以致不是出于天主？因此，既然宗徒更可能说真话，而此人作为权威不应凌驾于宗徒之上，那么女人是出于男人，不论是仅就肉身而言，还是就构成人性的整个整体而言（但我们对这些问题的说法并非绝对确定，仍在探求真理）；男人则是藉女人而生，不论他的完整人性是从父亲传给他并藉女人诞生，还是仅仅是肉身；关于这些，问题仍未确定。然而，\BibleQuote{一切都是出于天主}，关于这一点没有疑问；在这句话中包括了男人和女人的肉身、灵魂和神魂。因为即使他们并非由天主所生或来自天主，或作为祂本性的部分从祂而出，但他们仍是出于天主，因为凡被祂所创造、形成和造作的，其存在的实在性都是从祂而来的。\par

% structure: p0058 source=h2 target=subsection reason=relative-heading-rank; base=h2

\subsection{第二十八章 不可按字面强解自然的修辞格。}

他接着评论说：\Quote{但宗徒说：\BibleQuote{祂亲自赐给众人生命与灵，}然后又加上这些话：\BibleQuote{并且使全人类由一人血统所生，}\BibleRef{宗 17:25}将这灵魂与灵归诸创造者，就其起源而言，而将身体归诸流传。}确实，任何不愿随意否认灵魂流传的人，在清楚查明这见解是否正确之前，有理由从宗徒的话中理解到，他意指\Emph{由一人血统}这一措辞等同于\Emph{由一人}，这是一种以部分理解整体的修辞手法。那么，如果此人可以在这段经文中以部分理解整体：\BibleQuote{人成了有灵的生物，}\BibleRef{创 2:7}仿佛灵也被理解为包含在内，尽管圣经在那里并未提及灵，为什么不允许其他人对\Emph{由一人血统}这一措辞赋予同样广泛的意义，使得灵魂与灵可以被视为包含其中，理由是，由\Quote{\Emph{血}}这一术语所意指的人，不仅由身体构成，也由灵魂与灵构成？正如主张灵魂流传的辩论者，一方面不应过分苛责此人，因为圣经论及第一个人说：\BibleQuote{在祂内众人都犯了罪}\BibleRef{罗 5:12}（因为措辞不是\Quote{在祂内众人的肉体犯了罪}，而是\BibleQuote{众人，}即\BibleQuote{所有的人，}因为人不仅是肉体）——我再重申，他自己也不应被过分苛责，因为恰好写着\BibleQuote{所有的人，}如此可能他们被理解为仅仅就肉体而言；同样，另一方面，他也不应过分苛责那些持守灵魂流传的人，因为经句\Quote{全人类由一人血统，}似乎这段经文证明只有肉体是通过流传传递的。因为如果他们声称的为真，即灵魂不由灵魂生出，而只有肉体由肉体生出，那么\Quote{\Emph{由一人血统}}这一措辞，按照以部分代替整体的原则，并不意指整个的人，而仅指一个人的肉体；而另一措辞\BibleQuote{在祂内众人都犯了罪}必须被理解为仅指所有人的肉体，这肉体从第一个人传递下来，圣经以整体表示部分。反之，如果每个完整的人（由身体、灵魂与灵构成）由各自完整的个体流传为真，那么\BibleQuote{在祂内众人都犯了罪}这段经文必须按其本来的字面意义理解；而另一措辞\Quote{\Emph{由一人血统}}是比喻用法，以部分表示整体，即由灵魂与肉体构成的整体的人，或者更确切地说（正如此人喜欢说的），由灵魂、灵与肉体构成。因为圣经惯于使用两种表达方式，既以部分代替整体，又以整体代替部分。例如，部分暗示整体，如经上说：\BibleQuote{凡有血气的都要来到祢面前；}其中\Emph{血肉}一词意指整个的人。而整体有时也暗示部分，例如当说基督被埋葬时，实际上被埋葬的只是祂的肉身。至于宗徒证言中提到的说法，即\BibleQuote{祂赐给众人生命与灵，}我想经过前面的讨论，没有人会为此动摇。毫无疑问\BibleQuote{祂赐予；}事实无可争议；我们的问题是：祂如何赐予？是每次重新吹气，还是通过流传？因为完全恰当地说，祂被描述赐予人肉身的实质，尽管同时也不否认祂是通过流传的方式赐予的。\par

% structure: p0060 source=h2 target=subsection reason=relative-heading-rank; base=h2

\subsection{第二十九章 [XVIII.] — \Person{维克托}引用的第六段圣经经文}

现在让我们看看创世记的引文，其中女人是从男人的肋骨中造出来的，被带到他跟前，他说：\Quote{这是我的骨中之骨，肉中之肉。}\BibleRef{创 2:23} 我们的对手认为，\Quote{亚当本应说：\InnerQuote{我的灵魂中之灵魂，或我的神魂中之神魂}，如果这也从他而来。} 然而，事实上，主张灵魂传生观点的人觉得，他们的立场有一个更坚固的防御，即圣经记载告诉我们天主从男人肋骨中取出一根，形成了女人，却没有补充说祂向她的脸嘘入了生气；因此，正如他们所说，她已经从男人那里获得了灵魂。确实，如果她还没有获得灵魂，他们说，圣经肯定不会让我们忽视这个情况。至于亚当说：\Quote{这是我的骨中之骨，肉中之肉，}\BibleRef{创 2:23} 却没有加上\Quote{从我的灵魂或神魂，我的灵魂或神魂}，他们可以回答，正如已经指出的那样，\Quote{我的肉和骨}的表达可以用部分代表整体来理解，只是从男人取出的部分并非死的，而是有灵的；因为没有任何合理的理由否认全能者能做到这一切，其根据是没有人能够从男人的肉中切下一部分连同灵魂。然而，亚当接着说：\Quote{她要被称为女人，因为她是从男人取出来的。}\BibleRef{创 2:23} 那么，为什么他反倒不说（从而证实我们对手的意见）\Quote{因为她的肉是从男人取出来的}？事实上，持相反观点的人可能会争辩说，既然经文写的不是女人的肉，而是女人本身是从男人取出来的，那么她就应当被视为其整个本性都具有灵魂和神魂。因为虽然灵魂没有性别之分，但提到女人时，没有必要将她们与灵魂分开。否则，她们就不会这样被劝诫关于自我装饰：\BibleQuote{不以编发、黄金、珍珠或华贵的衣服为装饰，惟有（宗徒说）那以善行谈论虔敬的女子才相宜。}\BibleRef{弟前 2:9} 当然，\Quote{虔敬}是灵魂或神魂的内在原则；而她们仍被称为女人，尽管装饰涉及她们本性中无性别的那部分。\par

% structure: p0062 source=h2 target=subsection reason=relative-heading-rank; base=h2

\subsection{第三十章——从沉默论证的危险。}

现在，当辩论者如此交替争辩时，我这样判断他们：他们不得依赖不确定的证据；也不得对他们不知道的事情妄下断言。因为如果圣经说：\Quote{天主向女人脸上嘘入生气，她就成了有灵的生物，}即使那样也不能推论人的灵魂不是从父母传生的，除非同样的话也论及他们的儿子。因为有可能，从身体取出的无灵魂的肢体需要被赋予灵魂，而儿子的灵魂是从父亲传来，通过母亲传输而出。然而，在此点上完全缄默；它完全隐藏在我们的视野之外。没有否定，但也没有肯定。因此，如果圣经在某处可能不是完全沉默，那么这一点需要更清晰的证据来支持。由此得出，主张灵魂传生的人并没有从天主没有向女人脸上嘘气这一事实得到任何帮助；而那些因为亚当没有说\Quote{这是我的灵魂中之灵魂}而否认这一教义的人，也不应说服自己去相信他们一无所知的事情。因为正如圣经有可能对女人是否像男人一样通过天主的嘘气获得灵魂保持沉默，而没有解决我们面前的问题，反而留下开放的问题一样，同样地，尽管圣经对亚当是否说了\Quote{这是我的灵魂中之灵魂}保持沉默，同样的问题也可能保持开放和未解决。因此，如果第一个女人的灵魂来自男人，那么在他宣告\Quote{这是我的骨中之骨，肉中之肉}时，部分代表整体，因为不仅她的肉，而是整个女人都是从男人取出来的。然而，如果她不是从男人来的，而是像起初男人一样通过天主的嘘气而来，那么\Quote{她是从男人取出来的}这句话中，整体代表部分，因为假设中取出的不是她的整个自我，而是她的肉。\par

% structure: p0064 source=h2 target=subsection reason=relative-heading-rank; base=h2

\subsection{第三十一章——阿波利纳里派论证基督没有同样的人灵魂的论据。}

虽然如此，这些圣经章节并未解决这个问题，它们对目前所讨论之点确实没有决定性。但我确知这一点：那些认为第一个女人的灵魂不是来自她丈夫的灵魂的人，因为经上只说\Quote{我肉中的肉}，而没有说\Quote{我灵中的灵}，他们实际上是以与亚玻里拿留派及所有类似反对者相同的论证方式，来反对主的人性灵魂；他们否认主的灵魂，唯一理由是读到圣经说\Quote{圣言成了血肉}（\BibleRef{若 1:14}）。他们说，如果祂也有灵魂，就应该说\Quote{圣言成了人}。但这一伟大真理之所以用这些词语表达，实在是因为圣经习惯用\Quote{血肉}来指代整个人，如经上说\Quote{凡有血肉的，都要看见天主的救恩}。因为血肉若无灵魂，什么也不能看见。此外，圣经有许多其他章节明确表明，在基督这个人内不仅有血肉，而且有一个人的——即理性的——灵魂。因此，主张灵魂传生的人也可以理解，在\Quote{我骨中的骨，我肉中的肉}这话中，部分代表整体，灵魂也被包含在词语中，正如我们相信圣言成了血肉，并非没有灵魂。他们只需用明确的章节来支持他们关于灵魂传生的意见，正如其他圣经章节显示基督拥有人的灵魂一样。同样，我们也劝告另一方，即否认灵魂传生意见的人，他们应该为他们的断言——即每个新受造的人，灵魂都是由天主嘘气创造的——提出某些确凿的证据，然后坚持\Quote{这是我骨中的骨，我肉中的肉}这话不是作为部分代表整体的比喻说法（将灵魂包含在其意义中），而是仅仅按字面指血肉本身。\par

% structure: p0066 source=h2 target=subsection reason=relative-heading-rank; base=h2

\subsection{第32章 [XIX.]——维克多的自相矛盾：论灵魂的起源}

在这种情况下，我认为这篇论文现在必须结束了。它实际上包含了在我看来对当前讨论主题最主要的一切。阅读其内容的人将知道如何警惕，以免同意你寄给我的那两本书的作者，以致相信灵魂是由天主的嘘气所产生，却并非出于无中。事实上，假设这一点的人，即使口头上否认这个结论，实际上却断言灵魂具有天主的实体，并且是祂的后裔，不是由于恩赐，而是出于本性。因为一个人从谁那里获得其本性的起源，就必须认真承认他也从谁那里获得其本性的种类。但这位作者终究自相矛盾：他有时说\Quote{灵魂是天主的后裔——当然，不是由于本性，而是由于恩赐}；有时又说\Quote{它们不是由无中造成，而是来自天主}。因此，他毫不犹豫地将它们归于天主的本质，而这正是他先前否认的立场。\par

% structure: p0068 source=h2 target=subsection reason=relative-heading-rank; base=h2

\subsection{第33章——奥古斯丁不反对驳斥灵魂传生的意见，而支持灵魂嘘气的意见}

至于灵魂是由嘘气创造而非传生的意见，我们当然一点也不反对坚持它——只是要坚持的人必须已经成功发现一些新证据，或在正典圣经中找到了明确的见证，以解决这一极其棘手的问题，或在自己的推理中找到了不违背公教真理的证据。而不是像此人表现出的那样。他找不到任何值得说的事，同时又不愿停止其好辩的倾向，完全不估量自己的能力，为了避免无话可说，他大胆断言\Quote{灵魂应当被肉体玷污}，\Quote{灵魂应当成为有罪的}；尽管他在灵魂降生成人之前找不到任何善恶的功过。此外，\Quote{在未受洗就离开肉身的婴儿身上，原罪可能被赦免，并且应当为他们奉献基督身体的祭献}，这些婴儿尚未通过教会的圣事在基督内被结合，并且\Quote{他们离开今世而未受重生之洗，不仅能进入安息，甚至能获得天国}。他还提出了许多其他荒谬言论，若在此论文中收集并审视，显然会令人厌烦。如果灵魂传生的学说是假的，愿它的驳斥不要由这样的辩论者来进行；而愿为每次创造行动中嘘气新灵魂的对立原则的辩护，由更合适的人来做。\par

% structure: p0070 source=h2 target=subsection reason=relative-heading-rank; base=h2

\subsection{第34章——那些主张人的灵魂不是来自父母，而是由天主在每一例中重新嘘气的人必须避免的错误}

因此，所有希望坚持新灵魂是在人出生时被正确地说成是吹入而非源自父母的人，务必在我已提及的四个要点上处处谨慎。也就是说，不要让他们肯定灵魂因别人的原罪而成为有罪的；不要让他们肯定未领洗而死亡的婴儿能藉原罪的赦免以任何其他方式到达永生和天国；不要让他们肯定灵魂在其降生成人之前已在别处犯了罪，并因此被强迫引入有罪的肉身；也不要让他们肯定那些实际上并未在他们身上发现的罪，因预先知道而应受惩罚，尽管他们从未被允许进入能犯那些罪的生命。只要他们不肯定这些要点中的任何一个，因为每一个都是完全虚假和不敬的，他们就可以，若可能，提出圣经上任何决定性的证据关于这个问题；他们可以坚持自己的意见，不仅不受到我的任何禁止，甚至得到我的赞同和衷心感谢。然而，如果他们在神圣的圣经中发现任何非常明确的权威，并由于失败而不得不提出四个意见之一，那就让他们约束自己的想象，以免他们在困境中被驱使说出那个现在应受谴责且最近被定罪的贝拉奇异端，即婴儿的灵魂没有原罪。确实，一个人最好承认自己对一无所知的事无知，而不是陷入已被定罪的异端，或建立某种新异端，同时轻率地一再辩护那些只显示其无知的观点。此人还犯了一些其他荒谬的错误，甚至很多，在这些错误中他偏离了真理的常规之道，但尚未达到危险的地步；如果主愿意，我甚至想就他的著作写一些意见给他；或许我会向他一一指出，或指出其中许多，若我无法留意全部的话。\par

% structure: p0072 source=h2 target=subsection reason=relative-heading-rank; base=h2

\subsection{第三十五章 [XX.]——结论}

至于这本现在的论著，我认为最好只将其献给你，你对我们共同的信德以及我的品行都怀有善意而真诚的兴趣，作为真正的公教徒和好朋友。你会将本书交给任何你能找到对该主题感兴趣或认为值得信任的人阅读或抄写。在其中，我认为有必要压制并驳斥这位年轻人的僭越，但方式却要表明我爱他，希望他得到改正而非被定罪，并愿他在那伟大的家即天主教会中取得进步，天主的怜悯已带领他至此，使他能藉圣善的生活和纯正的教导，在那里\BibleQuote{成为贵重的器皿，圣洁的，合乎主用，预备行各样的善事}，见：\BibleRef{弟后 2:21}。但我还有进一步的话要说：如果我有必要向他表达我的爱，正如我真诚所做的，那么我更应该爱你，我的兄弟，我曾以最好的证明发现你对我怀有谨慎而清醒的感情和公教信德！你的忠诚的结果是，你以兄弟的真爱和责任，将那令你不悦、并在其中发现我的名字被以一种与你意愿相悖的方式处理的书籍，抄写并寄给了我。现在，我不仅不因你的这一慈善行为感到不悦，反而认为以友谊的真正要求，你若没有这样做，我本有权对你生气。因此我向你致以最诚挚的感谢。此外，我通过更清楚地表明了我接受你服务的态度，即在读过他的那些著作后，立即撰写了这本论著供你参考。\par

\SourceRecord{Translated by Peter Holmes and Robert Ernest Wallis, and revised by Benjamin B. Warfield.；Nicene and Post-Nicene Fathers, First Series；Vol. 5.；Edited by Philip Schaff.；Buffalo, NY: Christian Literature Publishing Co.,；1887.；<http://www.newadvent.org/fathers/15081.htm>.}

% structure: p0001 source=h1 target=section reason=page-title

\section{论灵魂及其起源（卷二）}

\Emph{以书信形式致司铎\Person{伯多禄}。}\par

\Emph{他劝告\Person{伯多禄}，不要因他可能使用\Person{维克多}写给他的关于灵魂起源的那些书而蒙上认可它们的名声，也不要将此人的鲁莽言论（违反基督信仰）当作天主教教义。他指出了\Person{维克多}的各种错误，且是极其严重的错误，并予以简要驳斥；最后他劝告\Person{伯多禄}本人努力说服\Person{维克多}改正其错误。}\par

致我亲爱的主内兄弟、同司铎\Person{伯多禄}阁下，\Person{奥斯定}主教，在主内致候。\par

% structure: p0005 source=h2 target=subsection reason=relative-heading-rank; base=h2

\subsection{第一章——堕落的雄辩是有害的才能。}

我收到了\Person{文森提乌斯·维克多}的两本书，是他写给阁下您的；它们是由我们的弟兄\Person{雷纳图斯}转交给我的，他虽是一位平信徒，但对关乎他自己和他所爱之人的信仰，有着审慎而虔敬的关切。读过这些书后，我发现它们的作者是一位口才极佳的人，他的言辞足够甚至太多；但关于他想要教导的题目，他却尚未得到充分的教导。然而，若蒙主恩赐，他也被赋予这种资格，他将对许多人有助益。因为他具备相当程度的能力来阐释和美化他的想法；所需要的只是他首先要小心正确地思考。堕落的雄辩是一种有害的才能；因为对于了解不足之人，它那流利的言辞总显得像是真理。我确实不知道您是如何接受他的书的；但如果我信息无误，据说您读过之后非常高兴，以至于您（虽是一位年长的司铎）竟吻了这位年轻平信徒的面庞，并感谢他教给您先前所不知道的事。在您的这一行为中，我并非不赞成您的谦逊；实际上，我甚为赞赏；因为您所称赞的不是这个人，而是通过他屈尊对您说话的真理本身；只是我希望您能向我指出，您通过他接受了什么真理。因此，如果您能在您对这封信的回信中告诉我他教了您什么，我将很高兴。我绝不耻于向一位司铎学习，既然您不羞于受教于一位平信徒，在宣扬和效法您的谦逊行为时，只要所学的课程是真实的。\par

% structure: p0007 source=h2 target=subsection reason=relative-heading-rank; base=h2

\subsection{第二章——他问\Person{维克多}传授了什么大知识。}

那么，您难道不明白，根据圣经所说\Quote{你将我的灵魂与我的灵分开}，有灵魂和灵两个东西？而且两者都属于人的本性，所以整个人是由灵、灵魂和身体组成的？然而，有时这两者被合称为\Emph{灵魂}；例如在\BibleQuote{人成了有生命的灵魂}这句话中。\BibleRef{创 2:7}此处也包含了灵。同样，在一些章节中，这两者被统称为灵，如所写的\BibleQuote{祂低下头，交付了灵}；\BibleRef{若 19:30}在这节经文中，也必须理解为灵魂。而且两者是同一实体？我想您已经知道这一切。但如果您不知道，那么您不妨知道，您并没有获得什么大知识，对此的无知不会带来多大危险。如果还要对这些点进行更精细的讨论，最好是与那人本人进行争论，我们已发现他言辞丰富。我们可以考虑的问题是：当提到灵魂时，是否也在术语中暗示了灵，使得两者包含灵魂，灵仿佛是它的一部分——事实上（如此人似乎认为的），在\Emph{灵魂}这个称谓下，整体是从部分命名的；或者，两者共同构成灵，真正称为灵魂的那部分是它的一部分；再或者，事实上，当\Emph{灵}这个术语被如此广泛地使用以致也包括灵魂时，整体是否不是从部分命名的，正如这个人所设想的。然而，这些只是细微的区别，对它们的无知当然不会带来什么大危险。\par

% structure: p0009 source=h2 target=subsection reason=relative-heading-rank; base=h2

\subsection{第三章——身体感官与灵魂感官的区别。}

我又想知道，此人是否教了您身体感官与灵魂感触之间的区别；以及您，在向此人学习之前已是相当年龄和地位的人，是否曾经认为，区分黑白的能力（麻雀也和我们一样能看到）与辨别正义与不义的能力（\Person{多俾亚}在失明后仍然能感知）是同一回事。如果您认为它们是同一回事，那么当您听到或读到\Quote{照亮我的眼睛，免得我沉睡致死}这句话时，您当然只想到了肉体的眼睛。或者如果这是一个隐晦的点，那么当您回忆起宗徒的话\BibleQuote{照亮你们的心眼}\BibleRef{弗 1:18}时，您一定以为我们在前额和脸颊之间某个地方有一颗心。然而，我远非这样认为您，所以您的这位教师不可能给了您这样的功课。\par

% structure: p0011 source=h2 target=subsection reason=relative-heading-rank; base=h2

\subsection{第四章——相信灵魂是天主的一部分是亵渎。}

倘若你在接受这位导师的教诲——你正庆幸已领受了——之前，竟偶然以为人的灵魂是天主本性的一部分，那你便不知道这见解是多么虚假、多么危险。倘若你只是从他那里学到灵魂并非天主的一部分，那么我劝你尽己所能地感谢天主，因为你在学习如此重要的功课之前没有被带离肉身。否则，你将以极大的异端者和可怕的亵渎者身份离开此生。然而，我从未相信你会如此，像你这样一位天主教的、地位不低的司铎，竟会以为灵魂的本性是天主的一部分。因此，我不得不向亲爱的你表达我的担忧：此人或许以某种方式教给你与你一向持守的信仰截然相反的东西。\par

% structure: p0013 source=h2 target=subsection reason=relative-heading-rank; base=h2

\subsection{第五章 [III.]——论受造物如何出于天主。}

正因为我不以为你——天主教会的一员——曾相信人的灵魂是天主的一部分，或灵魂的本性与天主的本性有任何同一，我才担心你或许被诱导赞同此人的见解，即\Quote{天主不是从无中创造了灵魂，而是灵魂如此出于祂，以致从祂流出}。他随同其他见解发表了这样的论断，使他在此议题上偏离常轨，坠入深渊。倘若他如此教导了你，我不愿你转而教我，反而愿你忘却所学。因为仅仅避免相信和说灵魂是天主的一部分是不够的。我们甚至不说圣子或圣神是天主的一部分，尽管我们肯定圣父、圣子、圣神都有同一本性。因此，仅仅避免说灵魂是天主的一部分是不够的，还必须说灵魂与天主并非同一本性。此人宣称\Quote{灵魂是天主的后裔，不是由于其本性，而是由于恩赐}，这是正确的；当然，不是所有人的灵魂，而是信徒的灵魂。但此后他又退回到他曾回避的论断，断言天主与灵魂有同一本性——虽未明言，却清楚而显明地如此意指。因为当他说灵魂出于天主，并非天主从另一本性或从无中创造，而是从祂自身创造时，除了他否认的恰恰是灵魂与天主自身同一本性之外，他还能让我们相信什么呢？因为每一本性，或是天主，祂无创造者；或是出于天主，以祂为创造者。那以天主为创造者、出于祂的本性，或非受造，或受造。非受造而出于祂的本性，或由祂所生，或由祂所发。所生的是祂的独生子，所发的是圣神，而此天主圣三具有同一本性。因为三者是一，每一位都是天主，三者共是一天主，不变、永恒，无时间的起始或终结。另一方面，受造的本性被称为\Quote{受造物}；天主是它的创造者，即有福的天主圣三。因此，受造物被称为\Emph{出于天主}，并非出于祂的本性。它被称为\Emph{出于祂}，是因为它在祂内有其存在的作者，并非由祂所生或所发，而是受祂创造、塑造、形成；有些从无其他物质——即绝对从无中，如天和地，或更确切地说，整个宇宙的物质在创造之时与宇宙同时存在——有些则从另一已受造存在的本性，如人出于尘土、女人出于男人、人出于父母。然而，每个受造物都出于天主——但出于天主作为它的创造者，或从无中，或从先前存在之物，而非作为从祂自身产生或发出者。\par

% structure: p0015 source=h2 target=subsection reason=relative-heading-rank; base=h2

\subsection{第六章——天主的本性会是可变的、有罪的、不敬的，甚至永远受罚的吗？}

然而，这一切我是对一位天主教徒说的：是在规劝他，而非教导他。因为我不认为这些事对您是陌生的，或您早已听说却不相信。我确信，您会如此阅读这封信，以致在它的陈述中也认出您自己的信仰，这信仰是藉着主的恩宠赐予，是我们在天主教会中共有的财富。所以，既然（如我所说）我现在是对一位天主教徒说话，请问您告诉我，您认为灵魂——我不说您的灵魂或我自己的灵魂，而是第一人的灵魂——是如何赐给他的？如果您承认它出自无，却是被天主所造并吹入他内，那么您的信仰与我的相符。反之，如果您认为它出自某个其他受造物，这受造物如同神圣工匠用以造灵魂的材料，正如尘土是形成亚当的材料，或肋骨是形成厄娃的材料，或水是创造鱼和鸟的材料，或地是形成地上动物的材料：那么这意见不是天主教的，也不是真实的。但若您认为——愿天主禁止——神圣造物主既非从无，也非从其他受造物，而是从祂自己，即从祂的本性，造了或正在造人的灵魂，那么您是从您的新导师那里学来的；但我不能因这发现而恭喜您或恭维您。您与他一起远远偏离了天主教信仰。更好的是，虽然这不真实，却更好，我说，更可容忍，您相信灵魂是由天主已经形成的其他受造实体所造，而非出自天主非受造的实体，如此，那可变的、有罪的、不虔敬的、且若坚持不虔敬到底必将遭受永罚的事物，不应以可怖的亵渎归诸天主的本性！远离吧，兄弟，我恳求您，远离此说，我不称它为信仰，而是可憎的不虔敬的错误。愿天主使您，一位稳重的司铎，免于被一个年轻的平信徒引诱的悲惨；并在您以为自己的意见是天主教信仰时，从信友的行列中失落。因为我不可像对待他那样对待您；这巨大的错误，若出于您，也不应得到与那年轻人同样的宽容，虽然您可能是从他那里学来的。\Emph{他}刚刚找到通往天主教会的路，以获得治愈和安全；\Emph{您}却在同一羊群的牧人中拥有地位。但我们不愿一只为躲避错误而来到主羊群的羊，其创伤被治愈的方式，竟是首先以他的传染性存在感染并摧毁牧人。\par

% structure: p0017 source=h2 target=subsection reason=relative-heading-rank; base=h2

\subsection{第七章．——认为灵魂是有形体的错误。}

但是如果您对我说，他没有教我这些；我也绝没有同意他的这个错误意见，尽管我被他那优雅动人的言谈的魅力所吸引；那么我衷心感谢天主。但我仍不禁要问，为什么，正如传闻所说，您甚至以亲吻表达了对他的感激，因为他教了您那些您在听他的讨论之前所不知道的事。如果这是一个虚假的报道，说您做了并说了那么多，那么我请您仁慈地给我这个保证，以便这无谓的谣言可以通过您亲笔写的书面声明而停止。然而，如果这是真的，您以如此的谦卑向这个人表示感谢，那么我确实会感到高兴，如果他并没有教您相信我已经指出为可憎且应小心避免的意见。我也不会责备[IV.]，如果您因与他讨论而获得了一些其他真实有用的知识而谦卑地感谢您的导师。但请问您，那是什么呢？是灵魂不是精神，而是身体？嗯，我确实认为，对此无知对基督徒的学习并无大碍；如果您沉溺于关于不同种类物质实体的更微妙的争论，我认为您获得的信息是困难多于实用。然而，如果主愿意我亲自写信给这个年轻人，正如我所愿，那么也许您的爱德会知道，您在何种程度上并非因他的教导而受益；尽管您为他所教的内容而高兴。现在我请求您不要因给我回信而感到烦恼；这样，那明显有用且与我们必不可少的信仰相关的事物，才不会碰巧变成别的东西。\par

% structure: p0019 source=h2 target=subsection reason=relative-heading-rank; base=h2

\subsection{第八章．——地狱中富人的口渴并不证明灵魂是有形体的。}

现在，关于他完全正确、观点极为健全地相信的那一点，即灵魂离开身体后，在它们必须接受最后审判（当身体归还它们时）之前，先受审判，并在它们曾在地上生活过的同一血肉中，或受折磨，或得荣耀；我问你，你真的不知道这事吗？有谁曾如此固执地反对福音，以至于没有听到这些真理，并在听到后相信它们，就是那个穷人死后被带到亚巴郎怀抱，以及那个富人在地狱中受折磨的比喻？但是，这个人可曾教导你，灵魂离开身体后，如何能从乞丐的指尖渴望一滴水？\BibleRef{路 16:24} 而他自己却承认，灵魂不需要身体的食物，除非为了保护包裹它的、正在朽坏的身体免于解体？这是他自己的话：\Quote{难道，}他问道，\Quote{因为灵魂渴望肉和饮料，我们就认为物质的食物进入它里面？} 然后不久他说：\Quote{由此可知并证明，肉和饮料的滋养不是灵魂所需要的，而是身体所需要的；同样，除了食物外，也提供衣服；所以，食物的供应似乎对于那个也适合穿衣服的本性是必要的。} 他相当清楚地阐述了这一观点；但他还添加了一些说明性的比喻，说道：\Quote{我们设想，一所房子的居住者在检查他的住所时会做什么？如果他看到那房子有摇晃的屋顶、倾斜的墙壁或薄弱的地基，他难道不会架起大梁、建起支撑物，以便通过他的关心和勤奋，支撑起那即将倒塌的建筑物，从而避免在住所的危险状态下，显然悬在居住者头上的危险吗？} 他说，\Quote{从这个比喻，看出灵魂如何渴望它的肉体，它无疑是从肉体中产生这种渴望本身。} 这个年轻人用这些非常清晰和适当的语言解释了他的想法：他断言，不是灵魂，而是身体需要食物；无疑是因为灵魂对身体的细心照顾，如同居住在一所房子里，通过明智的修缮防止了肉躯住所所面临的倒塌。那么，让他继续向你解释，这个富人的特殊灵魂，既然不再拥有一个会死的身体，却仍然感到口渴，并乞求穷人手指上的一滴水，它如此急切地要支撑以防止的是什么可能的倒塌？对这位教导老年人的精明教师来说，这是一个很好的棘手问题；让他自己探究，如果他能，就找出一个解决方案：那个在地狱中的灵魂乞求那么小一滴水的滋养，而它却没有一个破败的住所需要支撑，这是为了什么目的？\par

% structure: p0021 source=h2 target=subsection reason=relative-heading-rank; base=h2

\subsection{第九章 [五.]——无形体的天主怎能从自己呼出有形的实体？}

在相信天主确实是无形体这一点上，我祝贺他至少在此处没有受戴尔都良的狂言影响。因为\Emph{他}坚持说，正如灵魂是有形体的，天主也是如此。因此，特别令人惊讶的是，我们的作者在这个问题上与戴尔都良不同，却竭力说服我们，认为无形体的天主并非从无中造出灵魂，而是将灵魂作为有形的气息从自身呼出。这是何等奇妙的学习，使每一个时代都竖起耳朵聆听，并且为它的门徒赢得了年长者甚至司铎的追随！让这位杰出的人读一读他所写的，公开地读；邀请知名和不知名的人、有学问和没有学问的人来听。年长者，与你们年轻的导师们一起聚集；学习你们以往一无所知的事；现在听你们从未听过的事。看，按照这位经师的教导，天主创造一种气息，不是从某种已经存在的东西中，也不是从绝对不存在的东西中，而是从祂自己那完全无形体的本性中呼出一种形体，以至于祂实际上将自己无形体的本性改变为形体，然后这个形体才变成罪恶的形体。他难道不是说，当祂创造气息时，祂并没有从祂的本性中改变什么东西吗？那么，当然，祂不是从祂自己造出那气息：因为祂自己是一回事，祂的本性是另一回事。这个疯狂的人在想什么？但如果他说天主从祂自己的本性中创造气息，而祂自己却保持完全完整，那么问题就不在这里。问题是：那来自某种先前受造之物，也不是来自无，而是来自祂的东西，是否不是祂所是的，即相同的本性和本质？现在，在祂生了圣子之后，祂仍然完全完整；但因为祂从自己的本性生了圣子，祂并没有生出与祂自己所是的不同的东西。因为，撇开圣言取了人性成为血肉不谈，作为天主圣子的圣言是另一位，却不是另一物：也就是说，祂是另一位格，却不是不同的本性。这是怎么发生的，除非是因为祂不是从别的东西中受造，也不是从无中受造，而是从祂自己被生；不是为了成为比祂更好的，而是要完全成为\Emph{祂}所从出的那一位所是的，即同一本性、同等、同永恒、在各方面相似、同样不变、同样不可见、同样无形体、同样是天主；总之，要完全成为圣父所是的，除了祂实际上是圣子而不是圣父？但如果祂自己仍然是同一完美无缺的天主，却从自己造出某些不同于祂自己且更差的东西，不是从无中，也不是从其他受造物中，而是从祂自己；并且这某种东西作为形体从无形体的天主发出；那么，愿天主禁止一位公教徒接受这样的观点，因为它不是来自神圣的源泉，而是人类心智的虚构。\par

% structure: p0023 source=h2 target=subsection reason=relative-heading-rank; base=h2

\subsection{第十章 [VI.]——可能发现孩子与父母性情相似或不相似。}

再者，他是多么笨拙地努力使那他认为有形的灵魂，摆脱身体的情欲，提出了关于灵魂的婴孩时期；关于灵魂在麻痹和压抑下的情绪；关于身体肢体的截断，而不切割或分割灵魂的问题。但在处理这些要点时，我的责任是与他而非与您交涉；他应当努力为他所说的一切提出理由。这样，我们似乎就不会想就年轻人的作品对一位长者的庄重过于纠缠。至于在孩子们身上发现的与父母性情的相似，他并不争论是来自灵魂的种子。因此，这也是那些否认灵魂传生的人的意见；但持这种理论的反对方并不把他们的主张建立在这个基础上。因为他们也注意到孩子与父母性情不相似；他们推测，原因是同一个人自己常常拥有各种不同的、彼此不相似的性情——当然，不是他接受了另一个灵魂，而是他的生活变好或变坏了。所以他们说，灵魂不具有产生它的那个灵魂所拥有的同一性情，并非不可能，因为同一个灵魂在不同时期可能有不同的性情。因此，如果您认为您从他那里学到了这一点，即灵魂不是通过生育自然传给我们的——我只希望您从他那里发现了真相——我会非常乐意听从您，学习全部真理。但真正学习是一回事，自以为已经学会是另一回事。那么，如果您认为您已经学会了您仍然不知晓的东西，您显然没有学习，而是轻信了愉快的传闻。谬误在甜蜜中悄然潜入。我说这话，并不是因为我目前对认为灵魂是借天主嘘气全新创造的而非来自父母的论点感到确定；我认为这一点仍然需要那些能够教导它的人来证明。不；我这样说是因为这个人以这样一种方式讨论了整个主题，不仅没有解决仍然争论的点，甚至发表了无疑虚假的言论。为了证明有疑问的事情，他大胆地陈述了无疑该受谴责的事情。\par

% structure: p0025 source=h2 target=subsection reason=relative-heading-rank; base=h2

\subsection{第十一章 [VII.]——维克多暗示灵魂在降生成人之前有\Quote{状态}和\Quote{功绩}。}

你自己会犹豫不决，不谴责他关于灵魂所说的吗？他说：\Quote{你不愿意承认，}又说：\Quote{灵魂从有罪的肉体中汲取健康，你看出它通过肉体按时经过那圣洁状态，从而藉着它曾失去功绩的同一途径来改善自身状态？或者，因为圣洗洗涤身体，所以相信由圣洗所赋予的就不会传到灵魂或精神？因此，灵魂理当藉着肉体修复那它似乎已透过肉体逐渐失去的古老状态，以便藉着它曾当得被玷污的同一途径开始重生状态。}现在，请看这位导师陷入了何等严重的错误！他说：\Quote{灵魂藉着曾失去功绩的肉体来修复其状态。}那么，灵魂在肉体之前必然拥有某种状态和某种好的功绩，他希望灵魂透过肉体恢复这些，当肉体在重生之洗中被洁净时。因此，在肉体之前，灵魂曾在某处以良好的状态和功绩生活，这状态和功绩在它进入肉体时就失去了。他的话是：\Quote{灵魂藉着肉体修复它似乎已透过肉体逐渐失去的原始状态。}那么，灵魂在肉体之前拥有一个古老的状态（因为他用的词\Quote{\Emph{原始}}描述这状态的古老）；这古老状态除了是幸福可赞美的状态之外，还能是什么呢？现在，他断言这幸福是通过圣洗圣事恢复的，尽管他不承认灵魂的起源是通过传生来自那个曾在乐园中明显幸福的灵魂。那么，在另一段中他又说\Quote{他不断肯定灵魂不是由传生而存在，也不是从无中出来，也不是靠自己存在，也不是在身体之前存在}，这又怎么解释呢？你看他如何在这里坚持认为灵魂在身体之前就存在于某处，并且处于如此幸福的状态，以致同一幸福通过圣洗归还给它们。但是，仿佛忘记了自己的观点，他接着说到它的\Quote{开始重生状态藉着那——}意指肉体——\Quote{——藉此它当得被玷污。}在之前的陈述中，他指出了某种透过肉体而失去的善功；现在，他却谈到某种恶功，藉此灵魂不得不进入或被派到肉体中；因为他的话是：\Quote{藉此它当得被玷污；}如果它当得被玷污，那么它的功绩当然不可能是好的。请让他告诉我们，在被肉体玷污之前，它犯了什么罪，以致应得被肉体如此玷污。如果他能够，请让他向我们解释一件完全超出他能力的事，因为要发现关于这个问题有什么真实的话可以告诉我们，这显然远非他所能及。\par

% structure: p0027 source=h2 target=subsection reason=relative-heading-rank; base=h2

\subsection{第十二章 [VIII.]——灵魂如何当得降生成人？}

他随后又说：\Quote{因此，灵魂若当得成为有罪的，尽管它不能成为有罪的，却没有留在罪中；因为，正如在基督内所预表的，它注定不处于有罪的状态，正如它不能成为有罪的一样。}现在，我的弟兄，请问你真的这样想吗？无论如何，在阅读并仔细考虑了他的话，并且反思了是什么在听他宣读时引出了你的赞美，以及在他结束之后你表达感谢之后，你是否形成了这样的意见？我请你告诉我，这有何意：\Quote{虽然灵魂当得成为有罪的，但它不能成为有罪的。}他的措辞\Emph{当得}和\Emph{不能}是什么意思？因为它不可能当得它所谓的命运，除非它是有罪的；它也不会成为有罪的，除非它能够成为有罪的——以致在犯任何恶功之前的罪，它可能为自己建立一个境地，在那里，在天主舍弃之下，它可以推进到犯其他罪。当他说\Quote{不能成为有罪的}时，他的意思是不是，它除非进入肉体，否则不能成为有罪的？但是，如果它不可能在任何其他地方变得能够犯罪，除非进入那个特定的状态，那么它如何当得被派到一个它能成为有罪的状态呢？因此，请他告诉我们它如何如此当得。因为如果它当得变得能够犯罪，那么它必然已经犯了某种罪，因此它当得再次成为有罪的。然而，这些点或许看起来是晦涩的，或者可能被人讥讽地说成是这样的特性，但它们实际上是最明白清楚的。事实是，他不该说\Quote{灵魂当得藉着肉体成为有罪的}，因为他永远无法发现在肉体之前灵魂有任何善功或恶功。\par

% structure: p0029 source=h2 target=subsection reason=relative-heading-rank; base=h2

\subsection{第十三章 [IX.]——\Person{维克多}教导说天主阻挠自己的预定}

现在让我们继续讨论更明显的事情。因为当他被困在这些巨大困难中时，关于灵魂如何能被原罪的锁链束缚，当他们并非源自那首先犯罪的灵魂，而是造物主在每次出生时重新嘘气于有罪的肉体中——纯洁无染于罪的一切传染与传播——为了避免他的论点遭到反驳，即天主藉此嘘气使他们成为有罪的，他首先求助于来自天主预知的理论，即\Quote{祂已为他们预备了救赎}。婴孩藉此救赎的圣事领受圣洗，以便他们从肉体所染的原罪得以洗除，仿佛天主在补救祂自己的行为，因祂使这些灵魂受了玷污。但随后，当他谈到那些没有得到这种帮助、在领受圣洗前去世的人时，他说：\Quote{在此我不以权威自居，但我以一个猜测的例子呈献给你们。那么，我们说，对于婴孩，他们被预定领受圣洗，但由于此生的凋谢，在基督内重生之前就被匆匆接去。我们读到，}他补充说，\BibleQuote{记着说，关于这样的人：他不久便被接去，免得邪恶改变他的领悟，或欺诈欺骗他的灵魂。因此祂急忙将他从恶人中接去，因为他的灵魂悦乐了主；他在短时间内成了成全的，便满全了长久的岁月。}如今谁会不屑于有这样一位导师呢？那么，对于婴孩——人们往往希望他们在死前尽快领受圣洗——是否如果他们在此生被滞留极短的时间以便他们能领受圣洗、然后即刻去世，邪恶就会改变他们的领悟，欺诈就会欺骗他们的灵魂；而为了防止这事发生，匆忙的死亡来拯救他们，以致他们在领受圣洗之前就被突然接去了？那么，正是因着他们的圣洗，他们变得更糟，被欺诈所欺骗，如果他们是在领受圣洗之后被夺去的呢？哦，卓越的教导，值得赞赏和紧随！但他大大地依赖你们所有出席他朗读会的人的审慎，尤其是依赖你，他向你呈献这部论著并在朗读后亲手交给你的你，以为你们会相信他所引用的圣经是为了未领受圣洗的婴孩的情况，尽管它是论及所有那些圣人的年幼时期，这些圣人是愚昧之人认为受到苛刻对待的，每当我们突然被从今生移除，不被允许达到人们为自己贪图作为天主大恩的年岁。但是，他这些话是什么意思：\Quote{婴孩被预定领受圣洗，但由于此生的凋谢，在基督内重生之前就被匆匆接去}，仿佛某种命运、宿命或任何你喜欢的力量，不许天主成就祂所预定的事？而当他们悦乐了祂时，祂自己又怎能将他们匆匆接去呢？那么，祂果真预定他们领受圣洗，然后又亲自阻碍祂所预定之事的完成吗？\par

% structure: p0031 source=h2 target=subsection reason=relative-heading-rank; base=h2

\subsection{第十四章 [X.]——\Person{维克托}将那些未领受圣洗而死的婴孩送往乐园与天上的居所，但不送进天国。}

但我求你注意，他多么大胆，竟然对在如此深奥的问题上宁取谨慎而不取过分博学的犹豫不决感到不满：\Quote{我大胆说}（原话如此）\Quote{他们能获得原罪的赦免，但不能被接进天国。就如十字架上的盗贼，他告明了信仰却没有受洗，主没有给他天国，而是乐园；\BibleRef{路 23:43} 因此那句话仍然完全有效：\InnerQuote{除非人由水和圣神重生，不能进天国。}\BibleRef{若 3:5} 这尤其真实，因为主承认在他父的家里有许多住所，\BibleRef{若 14:2} 藉此指示居住者许多不同的功德；所以，在这些住所中，未受洗者得到赦免，受洗者得到藉恩宠为他预备的赏报。}你看到此人如何将乐园和父家的住所与天国分开，以便未受洗者也能在永恒福乐的地方得到丰富的供应。他这样说时，也没有看到，他不愿意将受洗婴儿的未来住所与天国区分开来，却毫不畏惧地将天主父的家本身或其各部分与天国分开。因为主耶稣没有说：在所有的受造物中，或在受造物的任何部分，而是说：\Quote{在我父的家里，有许多住所。}但未受洗的人如何能住在天主父的家里，除非他重生，否则他不可能有天主为父？他不应对天主忘恩负义，天主曾恩准将他从多纳徒派或洛加徒派中救出，却意图分裂天主父的家，将其一部分置于天国之外，使未受洗者能够居住。而且他本人以什么条件认为自己能进入天国，因为他将君王自己的家（无论君王愿意哪部分）从天国中排除？然而，从与他同钉十字架、将希望寄托于同钉之主的主旁边的盗贼的例子，以及从\Person{圣白小满}的兄弟\Person{迪诺克拉特}的例子，他论证说，即使对未受洗者，也可能给予罪的赦免和与真福者同住；仿佛任何不信（不信便是罪）的人已经向他显示盗贼和迪诺克拉特没有受洗。关于这些例子，我在写给我们的兄弟\Person{勒纳图斯}的书中已经更充分地解释了我的观点。你的仁爱之心若肯屈尊阅读那本书，将能查考；因为我确信，如果你向他借阅，我们的兄弟不会拒绝你。\par

% structure: p0033 source=h2 target=subsection reason=relative-heading-rank; base=h2

\subsection{第十五章 [XI.]——\Person{维克托}\Quote{裁定}应为未受洗死去的人奉献祭品。}

他仍然因犹豫不决而焦躁不安，几乎在他理论的可怕困境中窒息；因为很可能他比你以更敏锐的眼睛看出他所阐述的邪恶有多大，即婴儿的原罪在没有基督的圣洗圣事的情况下就被消除了。确实，为了在教会的圣事中多少寻找一些出路，虽然是迟来的，他说：\Quote{在他们身上，我确实决定，必须由圣洁的司铎不断奉献祭品和不停的牺牲。}那么，如果你愿意，你可以把他当作你的仲裁者，如果让他做你的导师还不够的话。让他决定，你必须为那些没有结合于基督的人奉献基督身体的祭献。这真是一个新奇的想法，与教会的纪律和真理的规范相悖；然而，当他敢于在书中提出时，他不是谦虚地说我宁愿认为，不是说我想，不是说我的意见是，也不是说我至少建议或提及，而是说我裁定；这样，如果我们（很可能）被他意见的新奇或乖僻所冒犯，我们可能会被他裁判的权威所慑服。我的兄弟，你自己要考虑如何能忍受他作为这些观点的导师。然而，具有正确情感的天主教司铎（你应在其中占据一席之地）决不能保持沉默——愿天主不许——并听此人宣读他的裁定，他们宁愿他恢复理智，既为持有这些意见，又为将它们付诸文字而表示后悔，并以最有益的补赎纪律惩罚自己。\Quote{现在，}他说；\Quote{正是根据玛加伯人在战斗中阵亡的例子，我得出必须这样做的必要性。当他们秘密地奉献了一些被禁止的祭品，并在战斗中阵亡之后，我们发现，}他说，\Quote{司铎们立即采取了这补救措施——献上祭品以解救他们的灵魂，这些灵魂因违禁行为的罪过而被束缚。}但他这么说，仿佛（根据他对故事的解读）那些赎罪祭是为未受割礼的人献上的，正如他决定我们的这些祭献必须为未受洗的人献上一样。因为割礼是那个时期的圣事，预表了今日的圣洗。\par

% structure: p0035 source=h2 target=subsection reason=relative-heading-rank; base=h2

\subsection{第十六章 [XII.]——\Person{维克托}应许未受洗者死后得乐园，复活后得天国，虽然他承认这与基督的话相悖。}

但是你的朋友，与他后来所显示出的相比，到目前为止所犯的错误还算勉强可以容忍。他显然有某种软化的倾向；当然，不是在他本应有所疑虑的事情上，即他竟敢断言未受洗者原罪也被解除，而且他们所有的罪都得到赦免，以致他们得以进入乐园，也就是一个极大幸福的地方，并在我们父家里拥有幸福的居所的权利；但他似乎对自己在天国之外也赐予他们较小福乐的住所表示了一些后悔。因此他接着说：\Quote{如果有人或许不愿相信，乐园是作为暂时和预备的恩赐赐予那个强盗或\Person{Dinocrates}的灵魂（因为他们在复活时仍可得到天国的赏报），虽然那主要经文阻碍着——\BibleQuote{人除非由水和圣神重生，不能进入天主的国。}\BibleRef{若 3:5}——他仍可毫不吝惜地赞同这一点；只要他赞美天主怜悯和预知的旨意和效果。}这些话是我在阅读他的第二卷书时抄录的。好了，现在，还有谁能在这一错误观点上表现出更大的胆大妄为、鲁莽或傲慢呢？他竟然引用并引起人注意主那句庄严的判决，把它放在自己的陈述中，然后说：\Quote{虽然这意见与\InnerQuote{主要经文}相悖——\BibleQuote{人除非由水和圣神重生，不能进入天主的国；}}他竟然敢抬起傲慢的头颅，去谴责君王的判决：\Quote{他仍可毫不吝惜地赞同这一点；}而且他解释他的观点是，未受洗者的灵魂有权暂时获得乐园作为恩赐；在这个类别中，他提到了垂死的强盗和\Person{Dinocrates}，仿佛他在指定，更确切地说，在预先判断他们的归宿；此外，在复活时，他要将他们迁到更好的住处，甚至使他们获得天国赏报的奖酬。\Quote{虽然，}他说，\Quote{这与君王的判决相抵。}现在，我请求你，我的弟兄，请认真思考这个问题：一个人若将自己与君王本身的权威相悖的赞同强加给任何人，他配得上君王的什么判决呢？\par

% structure: p0037 source=h2 target=subsection reason=relative-heading-rank; base=h2

\subsection{第十七章。—— 悖逆的怜悯与有怜悯的悖逆被摒弃。以殉道替代圣洗。}

新近出现的白拉奇异端者被公会议和宗座权威公义地定罪，因为他们竟敢将安息和救恩的处所赐予未受洗的婴儿，甚至在天国之外。如果他们不否认婴儿有原罪以及需要藉圣洗圣事赦免原罪，他们就不敢这样做。然而，此人承认在这点上的大公信仰，承认婴儿被原罪的锁链束缚，但他却不用重生的圣洗就解除了他们的锁链，死后又慈悲地将他们接入乐园；并且以更大的慈悲，在复活后甚至将他们引入天国。撒乌耳在他饶了天主命令杀死的君王时，就曾表现出这样的慈悲；\BibleRef{撒上 15:9}然而，他那悖逆的慈悲（或如果你愿意，有慈悲的悖逆）被公义地谴责和定罪，好叫人警惕，不要违背造人之主的判决而向他人施以仁慈。真理藉着其降生成人的口，如雷声般宣告说：\BibleQuote{人除非由水和圣神重生，不能进入天主的国。}\BibleRef{若 3:5}为了将殉道者排除在这条判决之外——他们因耶稣之名被杀，未来得及在基督的圣洗中受洗——他在另一处说：\BibleQuote{为我的缘故丧失自己生命的，必获得生命。}\BibleRef{玛 10:39}而且，使徒也大声宣告，没有任何未在基督教信仰的圣洗中重生的人可以免除原罪：\BibleQuote{因一人的过犯，众人都被定了罪。}\BibleRef{罗 5:18}为平衡这定罪，主只提供了祂救恩的帮助，说：\BibleQuote{信而受洗的必要得救；但不信的必被定罪。}\BibleRef{谷 16:16}在婴儿的情况下，这信仰的奥秘完全由代他们接受圣洗的担保人的回答来实现；若不实现，他们都因一人的过犯而进入定罪。然而，与如此清楚的真理宣告相对，有一种比可怜更愚蠢的虚荣站出来说：不是如此；婴儿不仅不进入定罪，尽管没有基督教信仰的圣洗解除他们原罪的锁链，他们甚至死后还能暂时享受乐园的福乐，复活后还将拥有天国的幸福。现在，此人若没有冒昧地试图解决一个超出他能力的关于灵魂起源的问题，他岂敢如此反对稳固的大公信仰呢？\par

% structure: p0039 source=h2 target=subsection reason=relative-heading-rank; base=h2

\subsection{第十八章 [XIII.]—— 维克托的困境与堕落。}

因为他被那些提出自然质疑的人逼到了可怕的困境中：\Quote{为什么天主将如此不公正的惩罚加于灵魂，竟愿意将它贬入一个罪恶的身体之中，既然灵魂与那开始成为罪恶的肉身结合，否则它本不可能成为有罪的？}当然，他们说：\Quote{如果天主没有使灵魂与有罪的肉身有分，灵魂就不可能成为有罪的。}然而，我这位对手无法发现天主如此行事的公义，尤其是对于那些未藉圣洗除去原罪而死亡的婴孩之永恒惩罚；他同样无法找出，为何良善正义的天主将那些祂预知不会从基督恩宠的圣事中获得益处的婴孩灵魂，藉由将他们送入得自亚当的身体，以原罪之链捆锁他们（灵魂本身毫无传生的玷污），并由此使他们应受永恒惩罚。他也不愿承认这些灵魂同样是从那第一个原始灵魂获得有罪的起源。因此，他宁愿以可悲的信德破灭来逃避，也不愿在他的争论航程中收帆稳桨，以明智的忠告遏制其致命鲁莽的进程。在我们年长者的眼中，他年轻的目光毫无价值，仿佛这最麻烦最危险的问题更需要雄辩的激流而非谨慎的忠告。连他自己也预见到了这一点，但毫无用处；因为，仿佛为了陈述反对者对他的责难，他说：\Quote{在此之后，其他侮辱性的指责又加到了那些辱骂我们之人的抱怨声中；我们仿佛被卷入旋风中，反复撞击在巨岩上。}说了这话之后，他为自己提出了我们已经讨论过的那非常危险的问题，他在其中摧毁了公教信仰，除非他藉由真正的悔改修复他所粉碎的信仰。我自己避开了那旋风和那些岩石，不愿将我脆弱的小舟托付给它们的危险；当我写作这一主题时，我表达的方式是解释我犹豫的原因，而非展示臆断的鲁莽。我这部小作品在他于你家中遇见时引起了他的嘲笑，他全然不顾后果地冲向了暗礁：他的行为显示了更多的血气而非智慧。然而，那种过度的自信将他引向何处，我想你现在自己也能察觉。但我更热切地向天主感恩，因为在此之前你就已看穿了它。因为当他拒绝在自己的航程结果尚不确定时刹住狂奔时，他着陆于他那可怜的企图，并坚持认为，对于未受基督重生而死亡的婴孩，天主立即赐予他们乐园，并最终赐予天国。\par

% structure: p0041 source=h2 target=subsection reason=relative-heading-rank; base=h2

\subsection{第19章——维克托依靠模棱两可的圣经章节}

诚然，他为了证明天主并非藉传生从原始灵魂衍生人的灵魂，而是如同最初那样藉吹气在每个人内形成灵魂所引用的圣经章节，是如此不确定和模棱两可，以至于它们极容易以一种与他所赋予的不同意义被理解。这一点，我想我在写给上文提到的那位朋友的书中已经足够清楚地证明了。他用来证明的章节告诉我们，天主赐予、或制造、或塑造人的灵魂；但它们关于祂从何处赐予它们，或用什么制造或塑造它们，却什么也没说：它们没有触及灵魂究竟是藉传生从第一个灵魂而来，还是如同第一个灵魂那样藉吹气而来的问题。然而，这位作者仅仅因为他读到天主\Quote{赐予}灵魂、\BibleRef{依 42:5}\Quote{制造}灵魂、\Quote{塑造}灵魂，就以为这些说法否定了灵魂的传生；而同一圣经的证言却表明，天主也\Emph{赐予}人的身体，或\Emph{制造}它们，或\Emph{塑造}和\Emph{造形}它们；尽管无人怀疑这些身体是由天主藉种子传生而赐予、制造和塑造的。\par

% structure: p0043 source=h2 target=subsection reason=relative-heading-rank; base=h2

\subsection{第20章——维克托引用圣经的沉默，并忽略圣经用法}

至于那节经文，它肯定\Quote{\BibleQuote{天主\OriginalTerm{由一个人}{of one blood}造了全人类}}（\BibleRef{宗 17:26}），以及\Person{亚当}所说的\Quote{\BibleQuote{这才真是我骨中的骨，肉中的肉}}（\BibleRef{创 2:23}），由于一处没有说\Quote{\Quote{\Emph{由一个灵魂}}}，另一处没有说\Quote{\Quote{\Emph{我灵魂的灵魂}}}，他便认为这就否定了子女的灵魂来自父母，或第一个女人的灵魂来自她的丈夫，好像如果经文如所提议的那样写成\Quote{\Quote{由一个灵魂}}而不是\Quote{\OriginalTerm{由一个人}{of one blood}}，所能理解的除了整个人类以外还能是别的，而并未否定身体的传生。同样，如果经文曾写成\Quote{\Quote{我灵魂的灵魂}}，当然也不会否定那显然是从男人取出的肉体。圣经经常以部分指整体，以整体指部分。的确，如果在这人所引以为证的经文中说人类被造不是\Quote{\OriginalTerm{由一个人}{of one blood}}而是\Quote{\OriginalTerm{由一个人}{of one man}}，那也不会损害那些否定灵魂传生之人的意见，尽管人并非只是灵魂，也非只是肉体，而是两者。因为他们可以这样回答：这里圣经可能以整体指部分，也就是说，以整个人来指肉体而已。同样，主张灵魂传生的人争辩说，在\Quote{\OriginalTerm{由一个人}{of one blood}}这节经文中，\Quote{血}字是以部分指整体的方式表示人。因为正如一方似乎因\Quote{\OriginalTerm{由一个人}{of one blood}}而非\Quote{\OriginalTerm{由一个人}{of one man}}的说法得到帮助，另一方显然也因明文记载的\Quote{\BibleQuote{故此，就如罪恶藉着一个人进入了世界，死亡藉着罪恶也进入了世界；这样死亡就传到了众人，因为众人都犯了罪}}（\BibleRef{罗 5:12}）而非\Quote{众人都因一个人的肉体犯了罪}得到支持。同样，一方似乎因圣经说\Quote{\BibleQuote{这才真是我骨中的骨，肉中的肉}}而得到帮助，理由是部分涵盖整体；另一方却从经文紧接着的\Quote{\BibleQuote{她应称为女人，因为她是由男人取出的}}得到一些优势。按他们的主张，后半句应写成\Quote{因为她的肉体是由男人取出的}，如果不是整个人（连同灵魂）而是仅她的肉体从男人取出的。然而，整个事情的真相是，在听了双方的论点之后，任何不存党派偏见的人立即看出，在这争论中模糊的引证是无用的；因为反对坚持灵魂传生意见的一方，不能引用那些只提到部分的经文，因为圣经可能在这些经文中以部分指整体；例如，当我们读到\Quote{\BibleQuote{圣言成了血肉}}（\BibleRef{若 1:14}）时，我们当然理解的不仅是血肉，而是整个人；同样，反对否认灵魂传生教义的另一方面，引用那些不提及人的部分而提及整体的经文也是无用的，因为在这些经文中圣经可能以整体指部分；正如我们承认\Person{基督}被埋葬了，而只有他的肉体被安放在坟墓里。因此我们说，基于这些理由，既不能轻率地建立，也不能同样轻率地推翻灵魂传生论；但我们补充建议，应当恰当寻找其他没有歧义的经文。\par

% structure: p0045 source=h2 target=subsection reason=relative-heading-rank; base=h2

\subsection{第二十一章 [XV.]——\Person{维克托}的困惑与失败。}

基于这些理由，我至今仍无法发现这位导师教了你什么，以及你有什么理由对他如此感激。因为关于灵魂起源的问题仍然悬而未决：天主是否从祂吹入第一个人的那一灵魂，藉由传生而赐予、塑造并造作人的灵魂；抑或祂在每一情况下都是亲自吹入，如同祂为第一个人所做的那样。因为天主\Emph{确实}塑造、造作并赐予人灵魂——基督信仰毫不犹豫地如此宣称。然而，当此人试图解决这问题而不估量自己的能力时，他否定了灵魂的传生，并断言造物主将灵魂吹入人内，这些灵魂纯洁无染，不受任何罪的沾染——不是从无中，而是从祂自己——这便以可耻的方式将可变性归之于天主，从而羞辱了天主的本性本身，这一归责必然无法立足。随后，为了避免任何可能导致天主被视为不义的暗示——即祂将那些纯洁无染于任何本罪的灵魂，若未藉基督信仰的重生而得赎，便以原罪之绳捆绑——他说出了一些话和想法，但愿他没有教给你。因为他竟将如此大的幸福和救恩赐予未受洗的婴儿，连白拉奇异端也不敢这样做。然而，尽管如此，当问题涉及到成千上万由不虔敬者所生、并在不虔敬者中死去的婴儿时——我并非指那些因慈善之人虽愿施洗却无力相助的，而是指那些无人能够或将会想到为其施洗，也无人为其奉献过或将奉献祭献——正如你的这位导师所想，这祭献甚至应为未受洗者奉献——他却没有找到解决之道。若有人就这些婴儿问他：他们的灵魂凭什么应受天主的处置，使祂将它们卷入有罪的肉身而招致永罚，永远不能在洗礼的水池中被洗净，也不能由基督体血的祭献而得赎？那么，他将要么感到完全茫然，因而以虽迟却真实的善意来看待我们的犹豫；要么会断定，必须为全地所有未受基督信仰洗礼而死的婴儿奉献基督的身体（他们的名字从未被听说过，因为他们在基督的教会中不为人知），尽管它们尚未被纳入基督的身体。\par

% structure: p0047 source=h2 target=subsection reason=relative-heading-rank; base=h2

\subsection{第二十二章 [XVI.]——\Person{伯多禄}在\Person{维克多}事件中的责任。}

我兄弟，千万不要让这样的看法令你喜悦，也不要因懂得这些而高兴，更不要妄自去教导它们。否则，就连那人也远比你更好。因为他在第一卷书开头写了这样谦逊卑微的序言：\Quote{虽然我渴望满足你的请求，但我不过是在清楚地证明我的冒昧。}稍后他又说：\Quote{事实上，我绝无把握能证明我所提出的论点；此外，我总是不愿坚持自己的意见，若发现它不可信；并且我诚心希望，若我的判断被否定，我便热切地遵循更好、更真的观点。因为让你容易接受对一主题的更真观点，显示了最佳意向和可赞扬的目的；而拒绝迅速转向理性的道路，则表明固执败坏的心智。}既然他真诚地说这一切，并且仍如所说的那样感受，那么无疑他对于正确的结局抱有很有希望的感觉。他以类似的语调结束他的第二卷书：\Quote{你切莫以为，}他说，\Quote{我让你作我言语的审判官，会有任何可能不光彩地反噬于你。唯恐某些好奇的读者锐利的目光有机会翻阅并遭遇我非法书页上可能残留的任何原始错误的踪迹，我请求你，若有必要，毫不留情地撕毁一页又一页；在对我施以批判性指责之后，再以涂抹掉那赋予我无价值之词语以形状的墨水来进一步惩罚我；这样，你得以充分利用机会，防止因你对我所持的强烈好感或我著作中潜伏的错误而引起的任何嘲笑。}\par

% structure: p0049 source=h2 target=subsection reason=relative-heading-rank; base=h2

\subsection{第二十三章 [XVII.]——哪些人不会因阅读有害书籍而受伤害。}

既然他以这样的保障开始并结束了他的著作，并将校正和修订它们的宗教责任放在了你的肩上，我只希望他能在你身上找到他所祈求的一切，使你能够\Quote{在怜悯中公义地纠正他、责备他；而罪人的油膏抹他的头}——无非是谄媚者不当的顺从和奉承者欺骗性的宽纵——却远离你的双手和双眼。然而，若你看到需要修正之处却拒绝加以纠正，你就得罪了爱；但若你因认为他的意见正确而觉得他无需纠正，那么你就是在与真理为敌。因此，他（因他过于乐意被纠正，只要有真正的批评者在旁）比你更好——如果你明知他犯错却藐视讥笑他，或不知他偏离正道却紧随其错误。因此，凡你在他寄给你并题献给你的书中发现的任何内容，我恳请你清醒且警觉地加以思量；你或许会比我本人发现更多应受指责的言论。至于其中值得赞扬和认可的内容——你从中学习到的任何好处，藉他的指导，或许是你以前确实不知道的——请明确告诉我们是什么，好让众人都知道，你是为此特别益处而向他表达感谢，而非因他书中许多应受谴责的言论。我是说，\Emph{所有}像你一样听他朗读其作品的人，或是后来亲自阅读这些作品的人，免得他们在你的事例（尽管不是你的榜样）之下，从他华丽的风格中饮下毒药，犹如从精美的酒杯中饮下，因为他们并不确切知道你自己饮了什么、什么你没有尝过，而且因你的崇高品格，他们以为从这泉源中饮下的任何东西都会有益健康。因为听、读和将大量内容存入记忆，若不是饮酒的过程，又是什么？然而，主曾预言有关祂忠信的追随者，甚至\BibleQuote{若他们喝了任何致命之物，也绝不受害。} \BibleRef{谷 16:18}因此，那些有判断力地阅读、按信德准则赞赏一切可称赞之事、并否定那些应被弃绝之事的人，即使他们将应受谴责的言论存入记忆，也不会因这些句子的毒害和败坏本质而受害。至于我自己，因主的仁慈，我绝不至于有丝毫懊悔，因为出于我们先前之爱，我已就这些要点向您——可敬且虔诚的您——提供了劝告和警告，无论您以何种方式接受这劝诫，我都认为您首先有权向我提出此要求。若我这封信能够发现或使您的信仰既免于我从这人的书中能指出的那些败坏和错误的观点，又在公教会正统中健全，我必将向祂——信赖祂的怜悯最为有益——献上丰厚的感谢。\par

\SourceRecord{Translated by Peter Holmes and Robert Ernest Wallis, and revised by Benjamin B. Warfield.；Nicene and Post-Nicene Fathers, First Series；Vol. 5.；Edited by Philip Schaff.；Buffalo, NY: Christian Literature Publishing Co.,；1887.；<http://www.newadvent.org/fathers/15082.htm>.}

% structure: p0001 source=h1 target=section reason=page-title

\section{论灵魂及其起源（第三卷）}

\Emph{致味增爵·维克多}。\par

\Emph{奥古斯丁向味增爵·维克多指出他若想成为天主教徒，在其关于灵魂起源的著作中应作的修正。那些在前两部写给雷纳图斯和伯多禄的书中已驳斥的意见，奥古斯丁在这第三部写给维克多本人的书中也作了简要的责难；此外，他还将这些意见归为十一项错误之目。}\par

% structure: p0004 source=h2 target=subsection reason=relative-heading-rank; base=h2

\subsection{第一章 [I.]—奥古斯丁写作的目的。}

关于我认为有义务写给您的，我亲爱的儿子维克多，我将这一点置于您心智中一切思想之上：如果我似乎轻视了您，那绝非我的本意。同时我必须恳求您，不要如此滥用我们的俯就，以至于因未受我轻视而认为您已获得我的认可。因为我爱您，不是要依从您，而是要纠正您；既然我绝不怀疑您有改正的可能，我就不希望您因我无法轻视我所爱之人而感到惊讶。既然在您与我们联合之前，我就有责任爱您，好使您能成为天主教徒；那么，现在既然您已与我们联合，我更应当爱您，以免您成为新的异端者，并使您成为如此坚定的天主教徒，以致没有异端者能抵挡您！就天主慷慨地赐予您的才智而言，如果您不认为自己已经明智，并以虔诚、谦卑而恳切的祈祷向那使人明智者祈求，使您成为明智者，并且宁愿不被错误引入歧途，也不愿受到那些走入歧途者的奉承，您无疑会是一个明智的人。\par

% structure: p0006 source=h2 target=subsection reason=relative-heading-rank; base=h2

\subsection{第二章 [II.]—为什么维克多取了味增爵的名字。恶人之名绝不应被他人采用。}

首先引起我对你有些担忧的，是你书中以你之名出现的标题；因为我向那些认识你、可能与你意见相同的人打听文森提乌斯·维克托是谁，发现你曾是多纳徒派，或者更确切地说是罗加图派，但最近已与天主教会共融。正当我为此欣喜，如同人们看见有人从谬误体系中获救时自然感到高兴一样——而对你，我的喜悦更大，因为我看到你那深深吸引我的才华并没有留在真理的敌人那边——你的朋友们又给了我额外的消息，使我在喜悦中感到忧伤：据说你希望在自己的名字前加上文森提乌斯这个名，因为你仍怀着深情尊重那位继承罗加图的人，此人拥有这个名字，你视他为伟大而圣洁的人，因此你希望他的名字成为你的姓氏。还有人告诉我，你还曾夸耀他在某种异象中向你显现，并协助你编写那些我在这本小册子中与你讨论的著作，甚至到了他亲自口授你应当写什么主题和论据的程度。如果这一切属实，我就不再奇怪你能作出那些陈述了——只要你耐心听从我的劝告，并以天主教徒的关注仔细考虑和衡量那些书，你无疑会后悔曾经提出它们。因为按照宗徒的描述，那\BibleQuote{伪装成光明的天使}（\BibleRef{格后 11:14}）的，在你面前已化作你认为曾经或仍是光明天使的形状。的确，当他伪装成异端者而不是光明天使时，他欺骗天主教徒的能力较小；但如今你已是天主教徒，我不愿你被他迷惑。他必定因你得知真理而感到痛苦，而且他因曾说服你相信谬误而得到的快乐越大，这痛苦也越深。然而，为了使你不要爱一个死人——这种爱既对你自己无益，也对他无用——我劝你暂且思考这一点：他当然不是一个正义圣洁的人，因为你因他们的异端而脱离了多纳徒派或罗加图派的陷阱；但如果你认为他是正义圣洁的，那你便因与天主教徒共融而自取灭亡。你其实只是在伪装天主教徒，若你在思想上与所爱的人相同；而且你知道圣经对此说了什么可怕的话：\BibleQuote{圣神必远离欺诈的人。}（\BibleRef{智 1:5}）然而，如果你真诚地与我们共融，而非仅仅伪装天主教徒，你又怎能仍如此爱一个死人，甚至现在还要夸耀他的名字，而你已不再容许自己留在他所犯的错误中？我们确实不喜欢你有这样的姓氏，仿佛你是死去的异端者的纪念碑。我们也不喜欢你的书有这样标题，好像我们在他的墓碑上读到的会是假的。因为我们确信文森提乌斯不是\Emph{\OriginalTerm{胜利者}{Victor}}，而是\Emph{\OriginalTerm{被征服者}{Victus}}——但愿这结出丰硕的果实，正如我们希望你被真理征服一样！然而，你的想法是狡猾而巧妙的：你把你希望我们认为是受他默感而口授的那些书命名为\WorkTitle{文森提乌斯·维克托}，这无异于暗示，与其说是你，不如说是他愿意被冠以胜利的称号，因为他向你启示了你著作的内容，从而征服了谬误。但这对你有什么益处呢，我儿？我恳求你，做一个真正的天主教徒，而不是伪装的，免得圣神远离你，使那位文森提乌斯完全不能帮助你——最恶毒的谬误之灵已经伪装成他来欺骗你；因为所有这些邪恶的意见都来自那恶者，尽管狡猾的欺骗使你相信相反的情况。如果这一劝告能使你以敬畏天主之人的谦卑和天主教徒的和睦顺从改正这些错误，那么它们将被视为一个过于热心的年轻人的失误，他宁愿改正它们而不是坚持它们。但若他凭借影响使你固执地坚持这些意见（愿天主阻止），那么就有必要把意见及其作者都谴责为异端，正如教会牧养补救的职责所要求的，在致命的感染悄悄蔓延到漫不经心的群众中之前予以制止，而有益的纠正却被忽视——这如同以爱的名义却没有爱的实质一般。\par

% structure: p0008 source=h2 target=subsection reason=relative-heading-rank; base=h2

\subsection{第三章——他列举希望文森提乌斯·维克托的著作中改正的错误。第一个错误。}

若你问我具体是哪些错误，你可以阅读我写给我们的弟兄、天主的仆人雷纳图斯和司铎彼得的信。你自己认为有必要给后者写信，内容正是我们现在讨论的那些作品，据你所说，是\Quote{顺从，}据你所说，\Quote{出于他本人的意愿和请求。}现在，我毫不怀疑，如果你愿意，他们会把我的论文借给你阅读，甚至未经请求就会引起你的注意。但无论如何，我不会错过这个机会，告诉你我希望在你的这些作品以及你的信仰中做出哪些修正。首先，你坚持认为\Quote{灵魂并非由天主从无中创造，而是出自祂自身}。在这里，你没有思考必然得出的结论，即灵魂必定具有天主的本性；你当然非常清楚这种观点是多么不虔敬。为了避免这种不虔敬，你应当这样说：天主是灵魂的创造者，即灵魂是由祂所造，但不是出自祂。因为凡是出自祂的（例如，祂的独生子）都与祂自己具有同一本性。但为了使灵魂不与它的造物主具有同一本性，灵魂是由祂所造，但不是出自祂。那么，请你告诉我它从何而来，否则就承认它是从无中而来。你所说的\Quote{它是天主本性嘘气中的一个微粒}是什么意思？那么，你是想说，出自天主本性的那嘘气本身（你所指的微粒属于它）与天主自身不是同一本性吗？如果这是你的意思，那么天主就从无中创造了那嘘气，而你将灵魂视为它的一个微粒。或者，如果不是从无中，请告诉我天主用什么创造了它？如果祂从自身创造了它，那么结论就是祂自身（这是绝不能承认的）是祂自己作品的质料。但你接着说：\Quote{然而，当祂从自身发出嘘气或气息时，祂同时保持完整无缺；}这就好像蜡烛的光，当另一支蜡烛从它点燃时，原来的光也保持完整，并且具有相同的本性，而不是别的。\par

% structure: p0010 source=h2 target=subsection reason=relative-heading-rank; base=h2

\subsection{第四章（[IV.]）——维克托的比喻：天主藉嘘气创造而不分施祂的实体}

\Quote{但是，}你说道，\Quote{当我们给一个袋子充气时，我们的本性或品质并没有被注入袋子，而使袋子膨胀的气流所包含的气息本身从我们发出，却没有丝毫减损。}现在，你详尽阐述并坚持这些话，强调这个比喻对于理解以下问题十分必要：天主如何在不损害自己本性的情况下，从自身创造灵魂，以及当灵魂如此从自身被造时，它却不同于祂自身。因为你问道：\Quote{对袋子的充气是我们灵魂的一部分吗？或者我们在充气时创造了人类吗？或者当我们通过充气将气息施于不同物体时，我们是否在任何方面受到损害？但我们从自身将气息转移到任何东西时并未受到损害，我们也不记得因给袋子充气而对自己造成任何伤害，尽管在这一过程中我们的气息仍保持其全部性质和数量。}现在，无论这个比喻在你看来多么优雅和贴切，我恳请你考虑它多么严重地误导了你。因为你断言，非实体的天主呼出了实体性的灵魂——不是从无中造的，而是从祂自身造的——而我们自己呼出的气息是实体性的，尽管比我们的身体更微妙；我们并非从灵魂中呼出它，而是通过身体结构内部的功能从空气中呼出。我们的肺就像一对风箱，由灵魂驱动（灵魂的命令也驱动身体其他部分），用于吸入和呼出大气。因为，除了构成我们食物和饮料的固态或液态养料之外，天主还赐予我们这第三种养料，即我们呼吸的大气；而且效果如此之好，以至于我们可以在没有食物和饮料的情况下存活一段时间，但没有这第三种养料我们无法生存片刻，这养料由四面环绕我们的空气在呼吸时供给我们。正如我们的食物和饮料不仅需要被摄入体内，还需要通过为此目的形成的通道排出，以防止任何一方造成伤害（无论是因不进入还是不离开身体）；同样，这第三种气态养料（不允许留在我们体内，从而不会因滞留而腐败，而是一旦被吸入便排出）被赋予了相同的通道，而不是不同的通道，用于它的进入和排出，即口、鼻或两者同时。\par

% structure: p0012 source=h2 target=subsection reason=relative-heading-rank; base=h2

\subsection{第五章——检验维克托的比喻：人在嘘气时给出什么吗？}

现在请你自己来证明我所说的话，为了你自身在你自己的情况中的满足：通过呼气放出气息，看看你是否能够长时间不吸回气息；然后再通过吸气将它吸回，看看如果你不再呼出它，你会感到多么不适。现在，当我们按照你的指示给袋子充气时，我们实际上做的正是维持生命所需的事情，除了在人工实验中我们的吸气稍强一些，以便我们能够呼出更强的气息，从而通过压缩我们吹入的空气来充满并鼓起袋子，这更像是用力吹气，而不是普通呼吸和呼吸的温和过程。那么，你根据什么理由说：\Quote{无论何时我们把气息从自己转移到任何物体上，我们都不会受到伤害，也从未记得给袋子充气对我们自己造成任何损害，尽管经过了这一过程，我们自己气息的全部品质和整体数量仍然留在我们身上}？很明显，我儿，如果你曾经给袋子充过气，你就没有仔细观察过你自己的行为。因为由于立即恢复呼吸，你没有察觉到你在充气时损失了什么。但是，如果你宁愿这样做，而不是仅仅因为你已经说出了你的陈述就僵化地坚持它们——不是给袋子充气，而是让自己充满，并用你浮夸话语的空谈充满你的听众（你本应用真实的事实来建设和教导他们），那么你可以极其容易地学到这一切。在目前的情况下，我不送你去任何别的老师那里，只送给你自己。那么，向袋子里好好吹一口气；立刻闭上嘴，捏紧鼻孔，这样你就会发现我对你说的话是真的。因为当你开始忍受实验中伴随的不可忍受的不适时，你希望张开嘴并松开鼻孔来恢复什么？当然，如果你的假设是正确的——即你呼气时什么也没有损失——那就没有什么可恢复的了。观察一下，如果你没有通过吸气恢复你呼气时失去的东西，你会陷入什么样的困境。再看看，如果不是通过呼吸的修复和反应，吹气会造成多大的损失和伤害。因为如果不是你用来填满袋子的气息全部通过重新打开的通道返回，以执行滋养你自己的功能，我不知道还有什么会留给你——我不说再给另一个袋子充气，而是维持你的生命本身！\par

% structure: p0014 source=h2 target=subsection reason=relative-heading-rank; base=h2

\subsection{第六章.— 根据真理修正的比喻}

那么，你本应在写作时想到这一切，而不该用你那个最喜欢的比喻——充气和可充气的袋子——把天主呈现在我们眼前，好像祂从某个已经存在的其他本性中呼出灵魂，正如我们自己从周围的空气中制造我们的气息一样；或者至少，你不该以一种实际上与你的比喻截然不同、且充满不敬的方式，把天主描绘成要么在自身不受损害的情况下，却从自己的本体中产生某种可变之物；或者更糟，创造它时竟以自己作为自己作品的质料。然而，如果我们想从我们的呼吸中取一个比喻来合宜地说明这个主题，那么以下这个比喻更加可信：正如我们呼吸时，不是从自己的本性中，而是因为我们并非全能，从我们周围的空气中制造气息——这空气在我们呼吸时吸进呼出；而这气息既无生命也无感觉，尽管我们自己有生命有感觉——同样，天主能够——不是从祂自己的本性中，而是（作为如此全能以至于能创造任何祂所愿的）甚至从完全不存在的东西中，即从无中——制造一个活着有感觉的气息，但显然是可变的，尽管祂本身是不变的。\par

% structure: p0016 source=h2 target=subsection reason=relative-heading-rank; base=h2

\subsection{第七章 [V].— 维克多似乎也将创造的气息归于人}

但你加上关于\Person{圣厄里叟}举的例子，说因他朝死者脸上吹气而使死人复活，这意味着什么呢？\BibleRef{列下 4:34} 难道你真的以为厄里叟的气息成了那孩子的灵魂吗？我简直无法相信连你也会如此远离真理。如果那曾从活孩子身上取出导致他死亡的气息，后来又被复原给他使他活过来，那么请问：当你说\Quote{厄里叟没有损失什么}时，这话有什么相干呢？好像可以想象有什么东西从先知传给了孩子使他复活似的？但假如你的意思不过是说先知吹了气而仍完好无损，那么你有什么必要用厄里叟来举这个例子呢？这例子用随便哪个呼气而不使任何人复活的人身上不是一样合适吗？此外，你说话很不谨慎（尽管天主禁止你相信厄里叟的气息成了复活孩子的灵魂！），因为你暗示你的意思是要把天主起初所行的事与先知所行的事区分开，即天主只吹了一口气，而先知吹了三口气。这是你的原话：\Quote{厄里叟照着原始创造的方式，朝\Place{叔能}妇人的死去孩子脸上吹气。当先知的吹气使神圣的力量感动死去的肢体，使之恢复原有的活力时，厄里叟并没有损失什么，通过他的吹气，死尸恢复了灵魂和呼吸。只有这一点不同：主只朝人脸上吹了一口气，他就活了；而厄里叟朝死人脸上吹了三口气，他又活了。}这样，你的话听起来好像只是吹气的次数造成了全部区别，使我们不相信先知实际做了天主所做的事。那么，这话需要完全修正。天主的工作与厄里叟的工作之间有如此大的区别，前者吹了生命的气息，使人成了活灵魂；后者吹的气本身没有知觉，也没有赋予生命，而是具有象征意义，为了某种示意。先知并没有真正通过赋予生命使孩子复活，而是通过给予爱求得天主行这奇事。至于你声称他吹了三口气，要么是你的记忆（常会发生），要么是经文版本有误，误导了你。我何必多言？你不该寻求例子和论证来支持你的观点，而应改正和改变你的意见。我恳求你，不要相信、不要说、也不要教导\Quote{天主造人的灵魂不是出于虚无，而是出于自己的本体}，如果你愿意作一个天主教徒的话。\par

% structure: p0018 source=h2 target=subsection reason=relative-heading-rank; base=h2

\subsection{第八章 [VI.]——维克多的第二个错误。（见本书卷一 26 [XVI.]。）}

我恳求你，不要相信、不要说、也不要教导\Quote{因此天主在无限的时间里不断地赐予灵魂，正如赐予者自身永远存在一样}，如果你愿意作一个天主教徒的话。因为时候将到，天主不再赐予灵魂，尽管他自身不会因此而停止存在。你的说法\Quote{不断地赐予}可以被理解为\Quote{不停歇地赐予}，只要人生育后代，正如论到某些人所说他们\Quote{常常学习，终久不能明白真理}。\BibleRef{弟后 3:7} 因为这里\Quote{常常}一词并非指\Quote{学习永不止息}，因为他们一旦停止在肉身中存在，或开始遭受地狱的火刑，就不再学习了。然而，当你说\Quote{不断地赐予}时，你并没有允许你的话被这样理解，因为你认为它必须应用于无限的时间。这还算小事；仿佛有人请你更清楚地解释\Quote{不断地赐予}这一短语，你接着说：\Quote{正如赐予者自身永远存在一样。}天主教的正统信仰彻底谴责这一断言。愿我们远离这种想法：天主不断地赐予灵魂，正如赐予它们的天主自身永远存在一样。他自身永远存在到从不停止存在；但灵魂，他将不会不断地赐予；当生育的时代结束，不再有孩子出生需要赐予灵魂时，他无疑会停止赐予。\par

% structure: p0020 source=h2 target=subsection reason=relative-heading-rank; base=h2

\subsection{第九章 [VII.]——他的第三个错误。（见本书卷二 11 [VII.]。）}

再次，我恳求你，不要相信、不要说、也不要教导\Quote{灵魂因肉身而理应失去什么，尽管在肉身之前它具有善功}，如果你愿意作一个天主教徒的话。因为宗徒声明\Quote{还没有生下来的孩子，既没有行善，也没有作恶}。\BibleRef{罗 9:11} 因此，灵魂在参与肉身之前，怎么可能有什么善功呢，如果它没有行过任何善事？你难道敢断言它在肉身之前曾度过善生，而实际上你甚至不能向我们证明它曾经存在过？那么，你怎么能说：\Quote{你不允许灵魂从有罪的肉身获得健康；那么，你可以看到，为了改善它的状况，它通过那使它失去功绩的同一肉身，适当地过渡到这神圣的状态}？或许你不知道，这些观点（即承认人的灵魂在肉身之前具有善的状态和善功）已经被天主教会定罪，不仅针对某些古代异端（此处我不提其名），而且最近也针对普里西利安派。\par

% structure: p0022 source=h2 target=subsection reason=relative-heading-rank; base=h2

\subsection{第十章。——他的第四个错误。（见本书卷一 6 [VI.]和卷二 11 [VII.]。）}

既不要相信，也不要说出，更不要教导：\Quote{灵魂藉着肉身恢复其原初状况，并借着那曾使它应受玷污的同一方式重生}，如果你希望做一个公教徒。我本可详论你在下一句中所显现的自相矛盾，即你说灵魂藉着肉身应得地恢复其原初状况，而这状况它似乎曾藉着肉身逐渐丧失，为的是它能开始藉着那曾使它应受玷污的同一肉身得到重生。在此，你——正是那个刚刚说过灵魂藉着肉身恢复其状况（因这状况它曾失去其\Emph{善功}，而你所意指的只能是\Emph{善功}，当然通过圣洗在肉身中恢复）的人——在你思想的另一个转折中说：灵魂曾藉着肉身应受玷污（在这句话中，所意指的不再是善功，而是\Emph{恶功}）。多么明显的矛盾！但我将略过它，只满足于指出：相信灵魂在降生成人之前应得善或恶，是绝对不符合公教信仰的。\par

% structure: p0024 source=h2 target=subsection reason=relative-heading-rank; base=h2

\subsection{第十一章（第八节）——他的第五个错误。（见上文第一卷第八章[第八节]及第二卷第十二章[第八节]。）}

既不要相信，也不要说出，更不要教导，如果你希望做一个公教徒：\Quote{灵魂在犯任何罪之前就应成为有罪的。}的确，应成为有罪是一种极坏的功过。而且，它当然不可能在犯任何罪之前，尤其是在进入肉身之前，就招致如此恶劣的功过，因为那时它不可能拥有任何或善或恶的功过。那么，你怎能说：\Quote{因此，那不能成为有罪的灵魂，如果应成为有罪的，却并未停留在罪中，因为正如它在基督内所预像的，它注定不应处于有罪的状态，正如它不能处于那种状态一样}？现在，请稍稍思考你在说什么，并停止重复这样的陈述。灵魂如何应成为有罪的，又如何不能成为有罪的？我请问你，那从未罪恶地生活过的，如何应成为有罪的？我再次请问，那不能成为有罪的，如何成了有罪的？或者，如果你说\Quote{\Emph{不能}}是指脱离肉身时的无能为力，那么在这种情况下，灵魂如何应成为有罪的，又因何种功过被遣入肉身，那时它与肉身结合之前并不能成为有罪的，以致应得任何恶呢？\par

% structure: p0026 source=h2 target=subsection reason=relative-heading-rank; base=h2

\subsection{第十二章（第九节）——他的第六个错误。（见上文第一卷第十至十二章[第九、十节]及第二卷第十三、十四章[第九、十节]。）}

如果你希望做一个公教徒，就要避免相信、说出或教导：\Quote{那些在领受圣洗之前被死亡夺走的婴儿，仍可获得原罪的赦免。}因为那些误导你的例子——即十字架上承认主的强盗，或圣伯尔贝都亚的兄弟狄诺克拉特——在你辩护这一错误意见时对你毫无助益。至于那强盗，尽管在天主的审判中，他可能被算作那些因殉道宣认而得到净化的人，但你无法知道他是否未领受圣洗。因为暂且不提一种观点，即他可能被从主肋旁同时流出的血和水所洒，\BibleRef{若 19:34}，他挂在旁边的十字架上，从而被一种最神圣的洗礼所洗净，又如果他在狱中领受了圣洗，正如后来有些人在迫害期间私下获得的那样，那又如何？或者如果他是在入狱之前领受了圣洗，那又如何？的确，如果他已领洗，他那时从天主那里获得的罪的赦免，并不会保护他免受公法的判决，就身体死亡而言。又如果他已领洗，却犯了抢劫和无法无天的罪行并招致惩罚，但藉着悔改加上他的洗礼，获得了那虽已领洗却犯下的罪的赦免？因为毫无疑问，他的信德和虔诚在主心中是明显的，正如它们在我们眼中从他的话语中是明显的一样。如果我们断定所有那些没有记录其洗礼就离世的人都是未领洗而死，我们就是在诽谤宗徒们本人；因为我们不知道他们中任何一位是在何时领洗的，除了宗徒保禄。\BibleRef{宗 9:18}然而，如果我们能将主对圣伯多禄说的话\BibleQuote{洗过澡的人，不需要再洗，只要洗脚就够了}\BibleRef{若 13:10}视为他们确实领洗的证据，我们该如何看待其他人，即那些我们甚至连这一点都没有读到的人——巴尔纳伯、弟茂德、弟铎、息拉、费肋孟，以及福音作者马尔谷和路加，还有无数其他的人？关于他们的洗礼，我们不应有任何怀疑，尽管我们没有读到任何记录。至于狄诺克拉特，他是个七岁的孩子；既然如此年龄领洗的孩童如今能背诵信经并在通常的考核中为自己回答，我不知道为什么不能设想他是在领洗之后被他不信主的父亲召回异教崇拜的亵渎和污秽，并因此被判受那痛苦，后来因他姐妹的代祷而得到释放。因为在关于他的记述中，你从未读到过他从未是基督徒，或是作为慕道者而死。但就此事而言，我们拥有的关于他的记述本身并不出现在那部圣经正典中，而在这类问题上，我们的证据总应从圣经正典中提取。\par

% structure: p0028 source=h2 target=subsection reason=relative-heading-rank; base=h2

\subsection{第十三章（第十节）——他的第七个错误。（见上文第二卷第十三章[第九节]。）}

如果你愿意做一个天主教徒，就不要冒险相信、说出或教导说：\Quote{主预定领洗的人能从他的预定中被夺走，或在全能者所预定的事在他们身上成就之前就死去。}在这种教义中，所赋予偶然性的力量，反对天主的力量，比我所能说的更大，由于这些偶然事件的发生，他所预定的事不被允许实现。几乎不必花时间或认真言语来警告接受此错误的人，反对那将把他吞没的绝对混乱的漩涡，因为我只要简短地警告那准备接受纠正的谨慎人提防威胁的祸害，就足以应付情况了。现在这是你的话：\Quote{我们说，对于预定领洗却因生命终结而在基督内重生之前就被匆匆夺去的婴孩，必须采取某种这样的方法。}那么，真有预定领洗的人因生命终结而在领洗之前就被抢先夺走了吗？难道天主能预定任何他或是在预知中看见不会发生、或是在无知中不知道不能发生的事，以致挫败他的目的或损害他的预知吗？你看，关于这个问题本可以讲许多重要的话；但我因不久前已论述过它而有所约束，所以满足于这简短而顺便的告诫。\par

% structure: p0030 source=h2 target=subsection reason=relative-heading-rank; base=h2

\subsection{第十四章——他的第八个错误。（参看上文，第二卷13[IX]章）}

如果你愿意做一个天主教徒，就拒绝相信、说出或教导说：\Quote{圣经论及那些在基督内重生之前就被死亡抢先夺去的婴孩时说：\BibleQuote{他被迅速提去，免得邪恶改变他的心意，或虚伪诱惑他的灵魂；因此天主急忙将他从恶人中间提去，因为他的灵魂悦乐了主；他在短期内成为完全，完成了长久的岁月。}}（见：\BibleRef{智 4:11}）因为这段经文与你所应用的婴儿毫无关系，而是属于那些领洗后度虔诚生活、不准久留世上的人——他们不是因年岁，而是因天上智慧的恩宠而成为完全。然而，你这个错误——认为这段圣经是关于未领洗而死去的婴孩——对圣洗本身造成了不可容忍的伤害，如果一个婴孩本来可以在领洗后\Quote{被匆匆夺去}，却在此之前\Quote{被匆匆夺去}，原因是：\BibleQuote{免得邪恶改变他的心意，或虚伪诱惑他的灵魂。}好像如果不在之前被夺去，这\Quote{邪恶}和这\Quote{虚伪诱惑灵魂}使灵魂变坏，就被相信是在圣洗本身之中！简言之，既然他的灵魂悦乐了天主，天主急忙将他从恶中提去；他就不再稍等片刻，以在他身上成就他所预定的；反而宁愿违背他预定的目的，急忙进行，免得未领洗孩童中如此悦乐天主的东西被他的洗礼所消灭！好像垂死的婴孩会在他所应被我们怀抱里跑去救他免于灭亡的事上灭亡！因此，关于\BibleBook{}{智慧篇}的这些话，谁若只是适当反思，还能相信、说出、写下或引用它们是关于未领洗而死去的婴孩的吗？\par

% structure: p0032 source=h2 target=subsection reason=relative-heading-rank; base=h2

\subsection{第十五章[XI]——他的第九个错误。（参看上文，第二卷14[X]章）}

如果你愿意做一个天主教徒，我恳求你，不要相信、说出或教导说：\Quote{有一些住处是在天主的国之外，就是主所说在他父家里的。}因为他并非如你所引证的那样说：\Quote{在我父那里（\OriginalTerm{在我父那里}{apud Patrem meum}）有许多住处；}虽然，即使他那样说，这些住处也很难被认为是在他父亲的\Emph{家里}以外的其他地方；但他明确地说：\BibleQuote{在我父家里有许多住处。}（见：\BibleRef{若 14:2}）现在，谁竟如此轻率，将天主家中的一些部分从天主的国分离出来；以致于，地上的君王不仅在自己的家中，也不仅在自己的国家，而且远及海外地区统治，而那位创造天地的君王却不能算是统治他自己的整个家呢？\par

% structure: p0034 source=h2 target=subsection reason=relative-heading-rank; base=h2

\subsection{第十六章——天主无所不在治理万物：然而\Quote{天国}可能并非处处皆是。}

然而，你也许不无可能争辩说：一切事物确实属于天主的国，因为他在天堂为王，在地上为王，在深渊，在乐园，在地狱（因为何处他不在那里为王呢？既然他的权能处处至高无上？）；但天国是一回事，除非他们在重生的水泉中被洗涤，否则没有人能被允许进入，这是主真实而定论的话；而另一回事是地上的王国，或任何其他受造部分的王国，在其中可能有天主家中的一些住处；但这些住处虽属于天主的国，却不属于那个具有特殊卓越和福乐的天国；因此，虽不能粗暴地将天主家中的部分和住所从天主的国分离，但并非所有的住所都预备在天国里；而且，即使在不位于天国的地方，天主若愿意为他们安排这样的住所，他们也可以幸福地生活，即使他们未领洗。他们无疑是在天主的国里，尽管（因未领洗）他们不可能在天国里。\par

% structure: p0036 source=h2 target=subsection reason=relative-heading-rank; base=h2

\subsection{第十七章——应如何理解天主的国在何处。}

如今，说这话的人无疑自以为说得很有道理，因为他们对圣经只有肤浅且不经心的看法；他们也不明白我们在祈祷中说\Quote{愿祢的国来临}\BibleRef{玛 6:10}时，使用\Quote{天主的国}这一短语是什么意思；因为那被称为天主的国，在其中祂的全家将幸福而永远地与祂一同为王。现在，就祂拥有万有的权能而言，祂当然已经在为王。那么，当我们祈祷祂的国来临，岂不是希望我们配得与祂一同为王吗？但那些将要遭受永火痛苦的人也将处于祂的权下。那么，我们是否要断言这些不幸的人也在天主的国里呢？显然，承蒙天主的国的恩赐和特权是一回事，受其律法约束和惩罚是另一回事。然而，为使你得到非常明显的证据，证明一方面不应将天国分给受洗者，另一方面将天主的国的其他部分分给未受洗者（正如你似乎所确定的），我请你听主亲自说的话；祂并未说\Quote{除非人由水和圣神重生，他不能进入天国}；祂的话是\Quote{他不能进入天主的国}。祂与尼苛德摩就此话题的对话如下：\Quote{我实实在在告诉你：人除非重生，不能看见天主的国。}请注意，祂在此不是说\Emph{天国}，而是说\Emph{天主的国}。然后，尼苛德摩回答祂说：\Quote{人已年老，如何能重生？难道他能再入母腹而重生吗？}主为解释而更清楚、更公开地重复了先前的陈述：\Quote{我实实在在告诉你：除非人由水和圣神重生，他不能进入天主的国。}请再次注意，祂用了同样的短语\Emph{天主的国}，而非\Emph{天国}。\BibleRef{若 3:3}值得注意的是，祂在第二次解释时变换了两个词（因为祂先说\Quote{除非人由\Emph{重生}}，然后用更完整的表达\Quote{除非人由\Emph{水和圣神}出生}来解释；同样，祂用更完整的短语\Quote{他不能\Emph{进入}}来解释\Quote{他不能\Emph{看见}}），但在此处祂没有任何变化；祂第一次说\Quote{天主的国}，后来精确地重复了同一短语。现在不必提出并讨论天主的国与天国是否应当理解为含义不同，还是同一事物有两个名称。只要发现没有人能进入天主的国，除非他在重生的洗礼中洗净，就足够了。我想你此时会明白，将置于天主的家中的任何住处与天主的国分开是多么偏离真理。至于你所持有的想法，即认为在住处的多样性中（主告诉我们祂父家中有许多住处），会发现有些没有由水和圣神重生的人居住，我劝你（如蒙允许）不要推迟改正，以便你能持守公教信仰。\par

% structure: p0038 source=h2 target=subsection reason=relative-heading-rank; base=h2

\subsection{第十八章 [XII.]——他的第十个错误（见前文第一卷第十三章 [XI.] 和第二卷第十五章 [XI.]）}

再者，如果你愿意成为公教信徒，我恳求你：不要相信、不要说、也不要教导说\Quote{基督徒的祭献应当为那些没有受洗而离世的人奉献}。因为你未能证明你所引用的\WorkTitle{玛加伯书}中犹太人的祭献是为任何未受割礼而离世的人奉献的。\BibleRef{加下 12:43}在你这个违背整个教会权威和教导的新奇意见中，你使用了非常傲慢的表达方式。你说：\Quote{为这些人，我断然决定，必须由圣洁的司铎持续奉献祭品和不断的祭献。}在此你作为一名平信徒，不愿谦卑地接受天主司铎的教导；你也不与司铎们一同探讨（这是你至少能做的），反而以骄傲的决断将自己置于他们之上。我的孩子，请摒弃这一切虚张声势；人不能以如此骄傲的举止行走在谦卑的基督所教导祂就是道路的道路上。\BibleRef{若 14:6}没有人能怀有如此骄傲的心态进入祂的窄门。\par

% structure: p0040 source=h2 target=subsection reason=relative-heading-rank; base=h2

\subsection{第十九章 [XIII.]——他的第十一个错误（见前文第一卷第十五章 [XII.] 和第二卷第十六章）}

再者，你若愿意作一个天主教徒，就不要相信、不要说、也不要教导说：\Quote{有些未受基督圣洗而离开此世的人，暂时不进入天国，而是进入乐园；但在死人复活以后，他们也会获得天国的福乐。}连佩拉基亚异端也不敢如此大胆地允许他们，尽管该异端主张婴儿不染原罪。然而，你，作为一个天主教徒，承认他们生在罪中；但你却凭借某种不可理解的乖僻，在你提出的新奇见解中，断言他们未经拯救的圣洗，便能解除与生俱来的罪，并被接入天国。你似乎没有意识到，在这点上你的观点远低于佩拉纠本人。因为他被主的那句不让未受洗者进入天国的圣言所震慑，不敢把婴儿送到那里，尽管他认为他们毫无罪过；而你对经上所记\BibleQuote{人除非由水和圣神重生，不能进入天主的国}\BibleRef{若 3:5}竟如此轻慢，以至于（且不说使你鲁莽地将乐园与天国分离的错误）你竟毫不犹豫地向某些人——你作为天主教徒相信他们生而有罪——许诺，即使他们未受洗而死，也能得到罪的赦免和天国。好像你能成为真正的天主教徒，是因为你反对佩拉纠而建立了原罪论，却对主关于圣洗的话加以拆毁，从而显出你是一个反对主的新异端。对于我们来说，亲爱的弟兄，我们不愿这样战胜异端：用一种错误去战胜另一种错误，而且更糟的是，用更大的错误去战胜较小的错误。你说：\Quote{也许有人不愿承认，那盗贼和狄诺克拉特的灵魂暂时被赋予了乐园，而他们在复活时仍有天国的继承权，因为那句主要的经文\—\InnerQuote{人除非由水和圣神重生，不能进入天主的国}\—\阻碍了这个观点；那么，他可以保留我在这点上毫不吝惜的同意；只要他充分尊重天主仁慈和预知的效果和旨意。}这是你自己的话，你在其中表示你同意那个人的说法，即乐园暂时赐给某些未受洗者，而在复活时为他们保留了天国的赏报，这与那句\Quote{主要经文}相悖，那句经文决定了未由水和圣神重生的人不能进入天国。佩拉纠害怕反对这句福音的\Quote{主要经文}，他不相信任何（他仍不认为是罪人的）未受洗者会进入天国。而你，相反，承认婴儿有原罪，却没有通过重生的洗池使他们解脱，而是让他们暂时住在乐园，随后甚至允许他们进入天国。\par

% structure: p0042 source=h2 target=subsection reason=relative-heading-rank; base=h2

\subsection{第二十章 [XIV.]——奥古斯丁敦促维克多纠正他的错误。（参见以上卷二第22章[XVI.]。）}

现在，我恳求你，如果你拥有一个天主教徒的心智——换句话说，如果你说以下这些话时是完全真诚的：你并不过于自信你所发表的每一条陈述都能证明；你常有的目标是不坚持自己的见解，即使它显得不太可能；如果自己的判断被否定，你很高兴采纳并追求更好、更真实的观点——那么就请你纠正这些错误以及诸如此类的错误，或许在你更仔细、更从容地阅读你的著作时还能发现其他错误。好了，我亲爱的弟兄，请证明你说这些话并非虚假；这样，天主教会就能因你的才能和品格而喜乐，因为你不仅具有天赋，而且还有智慧、虔诚和节制，而不是由于你固执地坚持这些错误而点燃异端的疯狂。现在，你也有机会表明你在紧接着我刚才提到的令人满意的声明的那段话中，表达你的感受是多么真诚。你说：\Quote{因为，一切最高尚的目标和值得称赞的意图的标志，就是愿意迅速转向更真实的观点；而拒绝迅速回到理性的道路，则显示出一种堕落和固执的判断。}那么，请展示出你受这种高尚目标和值得称赞的意图所影响，将你的心智迅速转向更真实的观点；不要表现出堕落和固执的判断，拒绝迅速回到理性的道路。因为如果你的话是真诚坦率地说出的，如果它们不只是嘴唇的声音，如果你真的从心里感受到它们，那么你必然憎恶任何拖延完成改正自己这一伟大善举的行为。的确，你说拒绝回到理性的道路显示出堕落和固执的判断，这不算什么，除非你加上\Quote{\Emph{迅速}}。加上这个词，你就向我们表明，那种从不进行改正的人的行为是多么可憎；因为即便只是迟迟才进行改正的人，在你看来也理应受到严厉的谴责，足以被描述为表现出堕落和固执的心智。因此，请听从你自己的劝告，并好好利用你雄辩的丰富资源（主要且大量地），以便你能迅速回到理性的道路——比你曾从那里偏离时更加迅速，当时你正值不稳定的年龄，智慧和学识都不足。\par

% structure: p0044 source=h2 target=subsection reason=relative-heading-rank; base=h2

\subsection{第二十一章。——奥古斯丁称赞维克多的才能和勤奋。}

若要完整處理和討論所有我希望在你的著作中——或更確切地說，在你本人身上——得到修正的要點，並為每一項修正提供簡短的理由，那將會耗費我過長的時間。然而，你也不應因此輕視自己，而認為你的才能和口才無足輕重。我在你身上發現了對聖經不小的記憶力；但你的學識與你的天賦及你所付出的努力相比，並不夠相稱。因此，我希望你一方面不要因過於歸功於自己而變得虛榮；另一方面，也不要因沮喪或絕望而變得冷淡和漠不關心。我只希望我能與你一同閱讀你的著作，並以交談而非書寫的方式指出必要的修正。這件事情通過我們之間的口頭溝通比用書信更容易完成。如果要用書寫來處理整個主題，將需要多卷著作。然而，我立刻將那些我希望以一個確定的數目綜合概括的主要錯誤提請你注意，以便你不要推遲對它們的修正，而是將它們從你的宣講和信仰中徹底驅逐；這樣，你所擁有的卓越辯論才能，藉著天主的恩寵，便能為你所用，有益於建樹，而不致於損害和破壞健全純正的教義。\par

% structure: p0046 source=h2 target=subsection reason=relative-heading-rank; base=h2

\subsection{维克多错误的概要重述}

這些具體的錯誤是什麼，我已經盡我所能地解釋過了。但我將以簡短的總結再重述一遍。第一項是：\Quote{天主不是從虛無中造了靈魂，而是從祂自身中造的。}第二項是：\Quote{正如施予的天主自身永遠存在，祂也永遠在無限的時間中施予靈魂。}第三項是：\Quote{靈魂在肉身之前曾擁有某些功勞，但因肉身而失去了。}第四項是：\Quote{靈魂藉著肉身恢復其原有的狀態，並通過那使之應受玷污的同一肉身而重生。}第五項是：\Quote{靈魂在犯罪之前就應得犯罪。}第六項是：\Quote{那些在受洗前被死亡攔阻的嬰兒，仍可獲得他們原罪的赦免。}第七項是：\Quote{主預定要受洗的人，可能從祂的預定中被帶走，或在全能者所預定的事在他們身上成就之前就死去。}第八項是：\Quote{那些在基督內重生之前被死亡攔阻的嬰兒，聖經所說的\InnerQuote{他迅速被接去，免得邪惡改變他的悟性}（參\BibleBook{}{智慧篇}）以及其餘相同意思的經文，正是指他們。}第九項是：\Quote{在天主的國之外，還有一些主所說在祂父家裡的住所。}第十項是：\Quote{基督徒的祭獻應當為那些沒有受洗就離開肉身的人奉獻。}第十一項是：\Quote{有些未受基督洗禮就離世的人，暫時不是進入天國，而是進入樂園；然而在死人復活時，他們甚至可達至天國的真福。}\par

% structure: p0048 source=h2 target=subsection reason=relative-heading-rank; base=h2

\subsection{固執使人成為異端者}

那麼，對於這十一項命題，它們極其明顯地乖謬，並與公教信仰對立；因此，你若不希望我們不僅為你來到我們公教的祭台前而歡喜，更為你真正成為一個公教徒而歡喜，就不應再猶豫，要將它們從你的心思、你的言語和你的筆下連根拔除並丟棄。因為如果你頑固地分別堅持這些主張，它們可能會產生與其觀點數量一樣多的異端。所以，請你想想，當它們集中在一個人身上時是多麼可怕，而如果它們被不同的人分別持有，每一項在持有者身上都是該受懲罰的。然而，如果你本人不再為它們爭辯，反而以忠實的言語和著作轉而攻擊它們，那麼你作為自我譴責者所獲得的讚美，將超過你對任何其他人所作的正當批評；而你對自己錯誤的修正，將比從未犯過這些錯誤更令人欽佩。願主臨於你的心靈和思想，並以祂的聖神傾注於你的靈魂，賜予你如此謙遜的準備、真理的光明、愛德的甘飴和和平的虔敬，使你在真理中勝過自己的心靈，勝過任何以錯誤反對真理的人。但我絕不希望你認為，持有這些意見已使你偏離公教信仰——儘管它們無疑與公教信仰對立——只要你能在那位無誤地洞察每個人內心的天主面前，真誠並出自內心地回顧你所說的話：你說你對自己提出的觀點不至於過度自信，它們都是可以證明的；你的恆常目標是，如果這些觀點被證明不合理，就不堅持自己的意見；因為當你的任何判斷被責備時，你樂於採納並追求更好更真的思想。即使所說的可能因無知而不合公教形式，但這種心態本身——出於修正的意圖和準備——就是公教的。說了這些之後，我現在必須結束這一卷，讓讀者稍作休息，準備在我開始下一卷時重新關注後續內容。\par

\SourceRecord{Translated by Peter Holmes and Robert Ernest Wallis, and revised by Benjamin B. Warfield.；Nicene and Post-Nicene Fathers, First Series；Vol. 5.；Edited by Philip Schaff.；Buffalo, NY: Christian Literature Publishing Co.,；1887.；<http://www.newadvent.org/fathers/15083.htm>.}

% structure: p0001 source=h1 target=section reason=page-title

\section{论灵魂及其起源（第四卷）}

\Emph{致\Person{文森特·维克多}。}\par

\Emph{他首先表明，他在灵魂起源问题上的犹豫不决受到了不应有的责备，他被错误地与牲畜相比较，因为他在这个问题上没有轻率下结论。然后，关于他自己毫不犹豫的声明，即灵魂是精神，而非身体，他指出\Person{维克多}多么轻率地反对这一主张，尤其当他徒劳地试图证明灵魂在其本质上是物质性的，并且人里面的精神与灵魂本身是不同的。}\par

% structure: p0004 source=h2 target=subsection reason=relative-heading-rank; base=h2

\subsection{第一章——本书的个性特点}

现在，在我的论文后续部分，我必须请求你听我说说关于我自己——尽我所能；或者更确切地说，正如那位掌握我们和我们言语者所允许的。因为你两次责备我，甚至提到了我的名字。在你的书开头，你说自己完全意识到自己缺乏技巧，没有得到学问的支持；而当你提到我时，赐予我\Quote{最有学问的}和\Quote{最有技巧的}赞美之词。然而，与此同时，对于那些你自认为完全熟悉的话题——对这些话题我或承认无知，或没有失当地自由地假设有所了解——你，虽然年轻且是平信徒，却不犹豫地责备我，一个老人和主教，而且是你自己判断为最有学问和最有技巧的人。好吧，就我而言，我对自己的大学问和大技巧一无所知；不仅如此，我很确定我不具备如此卓越的品质；此外，我毫不怀疑，一个没有技巧和学问的人偶尔也有可能知道一个有学问和有技巧的人所不知道的事；在这方面，我明确称赞你，因为你将爱真理置于单纯的人的尊重之上——因为即使你没有理解真理，你至少认为它是这样的。你这样做无疑是鲁莽的，因为你以为你知道你实际上一无所知的事；而且毫无约束，因为你不偏袒任何人，选择了公开表达你心中的一切。因此，你应该明白，我们应当更加谨慎地将主的羊群从他们的错误中召回；因为连羊群也不应该向牧者隐瞒他们在牧者身上发现的任何缺点，这显然是错误的。但愿你所责备我的是真正值得责备的事！因为我不否认，在我的行为和著作中，有许多点可以被一个明智的法官责备而不鲁莽。现在，如果你选择其中任何一点来责备我，我或许能够通过这些向你展示，我希望你在那些你明智而公正地谴责的细节上如何表现；此外，我（作为长者对幼者，作为有权柄者对必须服从的人）将有机会在你矫正下树立一个榜样，这个榜样在我方面不应过于谦卑，而是对我们双方都有益。然而，对于你实际责备我的那些点，它们不是谦卑迫使我去纠正的，而是真理迫使我去部分承认和部分辩护的。\par

% structure: p0006 source=h2 target=subsection reason=relative-heading-rank; base=h2

\subsection{第二章——\Person{维克多}认为\Person{奥古斯丁}应受责备的地方}

这就是那些责难。第一，我不敢就那些已赐予或正被赐予人类灵魂的起源作出明确的陈述——自从第一个人以来，因为我承认我对此的无知；第二，因为我说我确信灵魂是精神体，而非躯体。然而，你在第二点下包含了两个责难的理由：其一，因为我拒绝相信灵魂是物质的；其二，因为我肯定它是精神体。因为在你看，灵魂同时是躯体，又不是精神体。所以我必须请求你注意我为自己的辩护，以反驳你的责难，并请你抓住我的自我辩护所提供的机会，学习你自己身上也有哪些需要修正之处。那么，请回想你那本书中首次提到我名字的那句话。你说：\Quote{我知道，}你说，\Quote{许多名望极高的人，在被人请教时，要么保持沉默，要么没有清楚承认任何东西，而是在开始阐述时便从讨论中撤除了一切确定的东西。我在你府上读到的那一位博学之士和著名主教（人称\Person{奥斯定}）的多种著作，也是这种性质。我想，实情是，他们带着一种过度的谦逊和畏缩探究了这主题的奥迹，并在自己里面消耗了自己论文的判断，声称自己无法对此确定任何东西。但是，我向你保证，在我看来，一个人对自己陌生，或一个被认为已获得万物的知识的人，竟认为自己不认识自身，这极端荒谬且不合理。因为如果一个人不知道如何讨论和决定自己的性质和本性，人与禽兽有何区别？因此，圣经的话可正当地用在他身上：\BibleQuote{人虽然受尊崇，却不了解，他就像牲畜，与它们相比。}因为当良善仁慈的天主以理性和智慧创造万物，并造人作为有理性的动物，具有理解能力，赋以理性，并充满感觉——因为藉他明智的安排，他给所有不分享理性能力的受造物分配了位置——那么，有什么比认为天主竟不让人认识自己更不协调的观念呢？此世的智慧，确实常常费力以求达到真理的知识；它的研究无疑达不到目标，因为它不能藉什么媒介允许真理被确知；然而，关于灵魂的本性，有一些与它试图辨别的真理相近（我甚至要说，类似）的东西。在这种情况下，一个信仰宗教的人竟对这主题没有明智的观点，或禁止自己获取任何观点，这是多么不体面甚至可耻的事！}\par

% structure: p0008 source=h2 target=subsection reason=relative-heading-rank; base=h2

\subsection{第三章 我们对身体的本性认识多少？}

那么，你对我们的无知所施加的这一极其清晰而雄辩的谴责，严格地将你置于这样的必要之中：你必须知道与人性相关的每一件可能之事，以至于，若你不幸对任何细节无知，你就必须（记住，这不是我，而是你自己造成了这种必要）与\Quote{牲畜}相比。因为，虽然当你引证\Quote{人纵然受尊荣，却不领悟}这一经文时，你似乎特别针对我们发出责难——因为我们（不像你）在教会中占有尊贵的地位；然而，你也占据着人性中过于尊贵的等级，不至于不被置于牲畜之上，而按照你自己的判断，如果你碰巧对任何显然属于你本性的要点无知，你将不得不与牲畜相比。因为你不仅以你的责难玷污了那些与我同样受无知之苦的人——即关于人类灵魂的起源（尽管我甚至对这一点也并非完全无知，因为我知道天主向第一个人的脸上吹了一口气，\Quote{人就成了一个有灵的生物}，\BibleRef{创 2:7}——然而，我若没有在圣经中读到这一真理，是决不可能凭自己知道的）；而且你还明确问道：\Quote{如果一个人不知道如何讨论和确定自己的本质和性质，人与野兽之间有什么区别？}并且你似乎如此明确地持有你的见解，以至于认为一个人应当能够如此清楚地讨论和确定他自己全部本质和性质的事实，以致没有关于他自己的任何事能逃过他的注意。现在，如果事实果真如此，那么若你不能告诉我你头上头发的确切数量，我现在就必须将你比作\Quote{牲畜}。但是，如果在这一生中无论我们进步多少，你都允许我们对属于我们本性的某些事实无知，那么我想知道你的让步有多远，以免它可能恰恰包括我们现在提出的这一点，即我们决不知道我们灵魂的起源；尽管我们知道——这属于信德——毫无疑问，灵魂是来自天主的恩赐，而且它仍然与天主本身的本质不同。再者，你是否认为每个人对自己本性的无知恰好与你对它的无知处于同一水平？每个人对这个问题所拥有的知识也必须与你所能达到的知识相等吗？因此，如果他不幸比你多一点点无知，你就必须将他比作牲畜；同样，如果有任何人在这个问题上比你稍微聪明一点，他将乐于以同样的公正将你比作上述的牲畜。因此，我必须请求你告诉我，你允许我们对我们的本性无知到何种程度，以便我们能远离可怕的牲畜；并且我请你适当反思，是否知道自己对问题的任何部分无知的人，比那些认为知道自己实际上不知道的事的人离牲畜更远。人的全部本质无疑是精神、灵魂和身体；因此，任何要将身体从人性中分离出来的人都是不明智的。然而，那些被称为解剖学家的医生们，甚至通过解剖活人的方法（只要人在检查者手中能保持些许生命）进行了仔细的审视和研究；他们的研究深入了肢体、血管、神经、骨骼、骨髓、内脏；所有这些都是为了发现身体的本质。但这些人中没有一个因为我们对他们的学科无知而将我们与牲畜相比。但也许你会说，是对灵魂而非对身体本质无知的人才被比作牲畜。那么，你一开始就不应像你所做的那样表达自己。你的话不是：\Quote{如果一个人对灵魂的本质和性质无知，人与牲畜之间有什么区别？}而是说：\Quote{如果一个人不知道如何讨论和确定自己的本性和性质。}当然，我们的性质和本性必须与身体一起被考虑，但同时，对我们所组成的各个要素的研究在每个情况下是分别进行的。就我自己而言，如果我想展示我有多大能力科学而明智地处理人性的整个领域，我将需要写许多卷；更不用说有许多题目我必须承认自己的无知了。\par

% structure: p0010 source=h2 target=subsection reason=relative-heading-rank; base=h2

\subsection{第四章[III.]——呼吸问题是关乎灵魂、身体还是什么？}

但依你之见，我们在前书中所讨论的关于人的气息一事，属于什么？——是属于灵魂的本性，因为它是在人内产生气息的；还是属于身体的本性，因为身体由灵魂推动而产生气息；还是属于这空气的本性，因为通过其交替作用，我们发现它产生气息；或者更确切地说，属于三者：即灵魂作为推动身体的，身体通过运动接收并呼出气息，以及周围的空气，其进入使身体鼓胀，其离开使身体收缩？然而，尽管你博学善辩，你显然对此一无所知，当你设想、说出、写出并在聚集听你意见的众人面前宣读，说我们是出于自己的本性吹鼓一个袋子，而我们的本性却丝毫不因这操作而减少；尽管你可以很容易地查明我们是如何完成这一过程的，不是通过繁琐地查阅人类或受默感的著作，而是通过简单考察你自己的身体行为，无论何时你愿意。既然如此，我怎么能相信你能教导我关于灵魂的起源——我承认自己无知的问题——而你却实际上对自己不断用鼻子和嘴做的事以及为何这样做一无所知？愿主使我得以被你所劝，接受而非抵抗如此明显且触手可及的真理。也愿你不要在袋子的充气问题上质问你的肺腑，以至于宁愿将它们充气来反对我，而不是顺从它们的教导，当它们以完全的真理回答你的询问——不是通过言语和争论，而是通过气息和呼吸。那么，我就能耐心忍受你纠正和责备我对灵魂起源的无知；不仅如此，我还会衷心地感谢你，如果你除了责备我之外，还能以真理说服我。因为如果你能教导我所不知的真理，那么我有责任以一切耐心忍受你可能加给我的任何打击，不仅言语上，甚至用手。\par

% structure: p0012 source=h2 target=subsection reason=relative-heading-rank; base=h2

\subsection{第五章 [IV.]——惟独天主能教导灵魂从何而来。}

关于我们之间的这个问题，我向你的仁爱承认，我非常渴望知道两件事之一，如果可能的话——要么是关于灵魂的起源，对此我一无所知，要么是这些知识在我们今生是否可得。因为我们的争论难道不就是涉及那些我们被命令的事项吗？\BibleQuote{不要求过于高深的事，也不要探究超出你能力的事；但主所命令和教导你的一切，你要永远思想。}\BibleRef{德 3:21}这就是我所希望知道的，或是从天主自己——祂知道祂所创造的一切，或是从一个足够博学、知道自己所说的是什么的人那里知道，而不是从一个对自己呼出的气息一无所知的人那里。并非人人都记得自己的婴儿时期；你以为一个人能在没有神圣教诲的情况下，知道自己是如何开始存在于母胎中的吗？——尤其当人性的知识完全逃避他，使他不仅对自己内在之物无知，而且对从外部添加到其本性的东西也一无所知。我最亲爱的兄弟，你还能教导我或任何其他人，人出生时如何获得灵魂吗？当你仍然不知道生命是如何由食物维持，以至于如果暂时断绝食物就必定死亡时。或者，你能否教导我或任何其他人，人从哪里获得灵魂，而当你实际上仍然不知道袋子被吹鼓时从哪里得到充填？我唯一的愿望是，既然你不知道灵魂的来源，那么让我方面知道这些知识在我今生是否可得。如果这是那些过于高深、我们被禁止探究或查究的事之一，那么我们就有充分的理由害怕，担忧我们可能犯罪，不是由于对它的无知，而是由于对它的追求。因为我们不应该以为，一个主题要归于过于高深的范畴，就必须属于天主的本性，而不是我们自己的本性。\par

% structure: p0014 source=h2 target=subsection reason=relative-heading-rank; base=h2

\subsection{第六章 [V.]——关于身体本性的问题已经足够神秘，但并不比关于灵魂的问题更高深。}

对于这一说法，你有什么要说的：在天主的工程中，有些比认识天主本身更难——尽管祂对我们而言可以成为认识的对象？因为我们得知天主是天主圣三；但直到今天，我们不知道祂创造了多少种动物，甚至不是那些能够进入\Person{诺厄}方舟的陆地动物，\BibleRef{创 7:8} 祂所创造的——除非你碰巧知道了这个事实。此外，在\BibleBook{智慧}{篇}中写着：\Quote{因为他们既然能如此得力，得以认识并测量世界，怎么更不容易找到它的主呢？} \BibleRef{智 13:9} 这是因为我们讨论的对象\Emph{在我们内}，所以对我们不算太高吗？必须承认，我们灵魂的本性比我们的身体更内在。仿佛灵魂借助身体的眼睛从外部探索身体本身，比用自身从内部探索做得更好。因为身体内部的哪些地方没有灵魂存在呢？然而，即使关于我们框架的这些内在和关键的部位，灵魂曾通过身体的眼睛检查和搜索它们；它成功了解到的一切都是通过身体的眼睛获得的；并且，毫无疑问，所有物质实体都在那里，即使灵魂不知道它。既然我们的内部器官没有灵魂就无法生存，那么灵魂更能够给予它们生命，而不是认识它们。那么，灵魂的身体是比灵魂自身更高的认识对象吗？因此，如果它希望探究和思考：人类的种子何时变成血液，何时变成固体肌肉；骨骼何时开始变硬，何时充满骨髓；有多少种静脉和神经；前者通过什么通道和回路为整个身体提供灌溉，后者提供捆扎；皮肤是否算作神经，牙齿是否算作骨骼（因为它们有些不同，因为它们没有骨髓）；指甲与两者有何不同，在硬度上相似，同时与头发共有可生长和可修剪的特性；还有，那些他们称之为\Emph{动脉}的、其中循环空气而非血液的静脉有什么用途——我再说一遍，如果灵魂希望了解这些及其类似的身体本性上的点，难道应该对人说：\Quote{不要寻求对你过高的事，也不要探索超出你力量的事？} 但是，如果探究的是灵魂自身的起源（这件事它一无所知），那么，这个题目竟然不算过高或超出力量而无法理解吗？你认为灵魂不知道它是由天主嘘入的还是由父母遗传的，是荒谬的、不合情理的，尽管它一旦过去就不再记得这件事，并认为它属于那些无法追忆的遗忘之事——比如婴儿期，以及出生后紧接的其他生命阶段，尽管它们发生时无疑伴随着感觉。然而，你却不认为它对自己所支配的身体无知是荒谬或不合理的，也不知道任何关于身体的不属于过去之事而是现在事实的事——比如，它是否为了给身体生命而启动静脉的搏动，为了通过肢体运作而启动神经；如果是这样，为什么它只在特别意愿时才运动神经，而却持续影响静脉的脉动，即使在不意愿时；那个被称为\OriginalTerm{主导部分}{\GreekText{ἡγεμονικόν}}（灵魂的权威部分、理性）从身体的哪个部位实施其普遍统治，是从心脏还是从大脑，还是通过分工：来自心脏的运动和来自大脑的感觉——还是来自大脑，既有感觉又有自愿运动，但来自心脏的是不自愿的静脉搏动；还有，如果它两项都是从大脑发出的，为何它在不意愿时也有感觉，却只在愿意时才运动肢体？那么，既然只有灵魂自己在身体中完成这一切，它如何不知道自己在做什么？它的能力又从何而来？它如此无知并不丢脸。那么，你难道认为如果它不知道自身如何被造，从何而来，就是可耻的吗？因为它并非自己造了自己。好了，那么，没有人知道灵魂如何在身体中完成所有行动；你难道不认为这也属于那些被称为\Quote{对我们太高，超乎我们力量}的事情吗？\par

% structure: p0016 source=h2 target=subsection reason=relative-heading-rank; base=h2

\subsection{第七章——我们往往更需要关于我们最内在之事的教导，而非关于离我们更远之事}

但我必须向你提出一个由我们的主题引发的更广泛的问题。为什么只有极少数人知道所有人所做之事的缘由？或许你会告诉我：因为他们学会了解剖学或实验的艺术，这两者都包含在医生的教育中，只有少数人获得，而其他人拒绝获取这些知识，尽管他们当然可以，如果他们愿意的话。那么，我暂且不谈为什么许多人试图获取这些知识却不能，因为他们迟钝的智力（这其实是一件非常奇怪的事）阻碍了他们从他人那里学习他们自身和自己内里所发生之事。但现在我要问一个非常重要的问题：为什么我无需技艺就知道天上有太阳、月亮和其他星辰；但必须借助技艺才能知道，当我移动手指时，动作是从心脏、大脑、还是两者，或都不是开始的？为什么我不需要老师就知道那远高于我的事物；却必须等待他人才能知道那由我而发、在我之内所行之事？因为，虽然我们被认为是用心思考，虽然我们知道自己的思想，无需任何他人知道，但我们却不知道我们用来思考的心脏本身位于身体的哪个部位，除非某人教我们，而此人却不知道我们在想什么。我并非不知道，当我们听到应全心全意爱天主时，这并不是指我们肋骨下的那块肉，而是指产生我们思想的那股能力。而这正是用这个名称来恰当命名的，因为，正如运动不会在心脏（脉搏由此向四面八方辐射）中停止一样，在思想过程中，我们也不会停留在行动本身而不进一步思索。但是，尽管每一种感觉甚至是由灵魂传递给身体，为何我们能在黑暗中闭着眼睛，通过称为\Quote{触觉}的官能数清自己的外部肢体，却对灵魂本身中心区域的内在功能一无所知？在那里，正是那股能力赋予一切生命和活力——这一奥秘，据我所知，没有任何类型的医生，无论是经验派、解剖学家、教条派、方法派，还是任何活着的人，能有所知？\par

% structure: p0018 source=h2 target=subsection reason=relative-heading-rank; base=h2

\subsection{第八章——我们不记得我们的受造}

任何试图探究此类知识的人，可能不无道理地被我们之前引用的话语所告诫：\Quote{不要探寻高于你的事，也不要搜索超过你能力的事。}现在这不仅是一个高度问题，如超出我们身高的高度；而是我们的智力无法达到的高度，以及我们的精神力量无法应付的强度。然而，这既不是诸天之上的天，也不是星辰的尺度，不是海洋和陆地的范围，也不是最深的地狱；而是我们自己，我们无法理解；是我们自己，在我们过高过强的状态下，超越了我们自身知识的谦卑界限；是我们自己，我们无法拥抱，尽管我们当然没有离开自己。但我们并不因为不完全发现我们自己是什么，就被比作牲畜；然而你认为，如果我们忘记了自己曾是什么，即使我们一度知道，我们就应得这种羞辱性的比较。我的灵魂现在不是从我的父母衍生而来，现在也不是从天主那里接受嘘气。无论祂使用了这两种过程哪一种，祂在创造我时使用了它；祂现在没有对我或在我内使用它。它已经过去并消失了——对我而言，不是现在的事，也不是近期的事。我甚至不知道我是否曾意识到它然后又忘记了；还是在那事发生时，我甚至无法感觉到并知道它。\par

% structure: p0020 source=h2 target=subsection reason=relative-heading-rank; base=h2

\subsection{第九章——以\Person{辛普里丘}的非凡记忆为例说明我们对自己的无知}

现在请看，当我们活着，当我们知道我们活着，当我们确信我们拥有记忆、理解和意志时；我们吹嘘自己对自己的本性有丰富知识——请看，我再说一次，我们对自己的记忆、理解或意志有何用处，是多么完全无知。有一位从年轻时就是我的朋友的人，名叫\Person{辛普里丘}，拥有精确而惊人的记忆力。有一次我请他告诉我\Person{维吉尔}的全部作品倒数第二行是什么；他毫不犹豫、完全准确地回答了我的问题。然后我请他重复前面的各行；他这样做了。我确实相信他能一行一行倒背\Person{维吉尔}。因为无论我在哪里考他，他都做到了。同样地，在散文方面，对于\Person{西塞罗}的任何一篇他已经背诵过的演说，他也会应我们的要求，按我们指定的程度倒背出来。我们表示惊讶时，他呼求天主作证，在这次试验之前，他根本不知道自己有这种能力。因此，就记忆力而言，他的心智当时才得知自己的潜力；而这种发现除非通过试验和实践，否则永远不会发生。此外，他在试验自己的能力之前当然就是同一个人；那么，他怎么会不认识自己呢？\par

% structure: p0022 source=h2 target=subsection reason=relative-heading-rank; base=h2

\subsection{第十章——记忆的可靠性；记忆的不可探寻的宝藏；无人能充分理解一个人的理解能力}

我们常常以为能将某事记在记忆中；因此认为如此，便不写下来。但后来，当我们想要回想起它时，它却不浮现于脑海；于是我们懊悔曾认为它会返回记忆，或没有把它写下来以防它逃逸；看，忽然间，并未刻意寻求，它却来临。那么，当我们如此思想时，难道我们不是我们自己吗？当我们不能再思想它时，我们就不再是原来的我们吗？然而，这是如何发生的呢？我不知道我们如何远离自己、否认自己；同样也不知道我们如何恢复、回归自己。仿佛我们是另一个人，在别处，当我们寻找存放在记忆中的东西却找不到时；我们自己无法回归自己，仿佛置身他处；但后来，当我们找到自己时，又再次返回。因为我们在何处寻找，岂非在我们自身内？我们寻找的，岂非就是我们自己？仿佛我们并不真正居于己身，而是去了某处。你难道不惊惧地观察到如此深奥的奥秘吗？这一切岂非就是我们的本性——不是它过去的样子，而是它现在的样子？看，我们寻求的远超过我们理解的。我常相信，若我用心思考，就能理解他人向我提出的问题。然而，我思考了，却未能解决；许多次我并未如此相信，却得以断定。可见，我对自己理解力的力量并不真正了解；我想，你也是如此。\par

% structure: p0024 source=h2 target=subsection reason=relative-heading-rank; base=h2

\subsection{第十一章——宗徒伯多禄说他愿为主舍命时并未说谎，只是不知主的旨意。}

然而，也许你因我承认这一切而鄙视我，便因此将我比作\Quote{牲畜}。但就我而言，我不会停止劝告你，或（你若不听我）至少警告你：宁可承认这共同的软弱，德行在其中得以成全；免得你假设未知之事为已知，而未能抵达真理。因为我认为有些事连你也想了解却不能；你绝不会寻求了解，除非你希望有朝一日能研究成功。因此，你也对自己理解力的力量无知，却宣称完全了解自己的本性，不愿跟随我承认无知。还有，意志呢？我该对它说什么？我们显然夸耀自由意志。的确，真福宗徒伯多禄愿意为主舍命。他无疑真心愿意，许下承诺时也未背叛主。但他的意志完全不知自己的力量。因此，这位伟大的宗徒，曾认出他的师傅是天主子，却对自己一无所知。因此，我们清楚知道自己愿意或\Quote{不愿}某事；但尽管我们的意志是好的，我亲爱的孩子，除非我们自欺，我们却不知道它的力量、它的资源、它在怎样的诱惑下会屈服或抵抗。\par

% structure: p0026 source=h2 target=subsection reason=relative-heading-rank; base=h2

\subsection{第十二章【VIII】——宗徒保禄能知晓第三层天与乐园，却不知自己是否在身内。}

请看，我们对自己本性的多少事实，不是过去而是现在的，不仅关于身体，也关于内在的人，一无所知，却不应被比作牲畜。然而，这就是你认为我该当的羞辱比较，因为我对自己灵魂的过去起源没有全备的知识——尽管我并非完全无知，因我知道是天主赐予我灵魂，且它并非出自天主。但我何时能数算尽所有关于我们神魂和灵魂本性的细节，而我们对它们无知？我们反倒应在天主面前发出圣咏作者所发出的呼喊：\BibleQuote{关于祢的知识，为我太奇妙；它极其困难，我不能达到。}为什么他加上\Emph{为我}一词呢？岂不是因为他推测天主的知识对他是如何不可理解，因为他连自己也无法理解？宗徒被提到第三层天，听见不可言传的话，是人不可说的；这事是发生在身内还是在身外，他自称不能说；\BibleRef{格后 12:4}但他却不怕你把他比作牲畜。他的神魂知道自己在第三层天，在乐园；却不知道自己是否在身体内。当然，第三层天和乐园不是宗徒保禄自己；但他的身体、灵魂和神魂才是他自己。看，这奇妙的事：他知道那些伟大、崇高、神圣的非己之事；但对属于自己本性的事，他却无知。在如此深奥的知识中，谁能不对他自己存在的极大无知感到惊讶呢？简言之，若不是那位从不犯错者告诉我们\BibleQuote{我们不知道该怎样祈祷才算合适}，谁会相信呢？\BibleRef{罗 8:26}那么，我们的倾向和目标主要应成为什么？岂不是\BibleQuote{向着前面的事伸展}？然而，你仍把我比作牲畜，如果我在身后的事中忘记了自己起源的任何事——尽管你听见同一宗徒说：\BibleQuote{忘记背后，努力面前的，向着标竿直跑，要得天主在基督耶稣内从高天召我来得的奖赏。}\BibleRef{斐 3:13}\par

% structure: p0028 source=h2 target=subsection reason=relative-heading-rank; base=h2

\subsection{第十三章【IX】——圣神为我们代求是什么意思。}

\Quote{我们不知道该如何祈祷才对}？这也许并不是那么不可容忍。因为按照正直公义判断的指示，我们更看重未来而非过去；既然我们的祈祷所关涉的不是我们过去如何，而是我们将要如何，那么不知道该如何祈祷，比起不知道我们灵魂起源的方式，当然更为有害。但请回想我重复的是谁的话，或者你自己再读一遍，并思考它们从何而来；不要用你的责备投掷我，免得你扔出的石头落在一个你不想击中的人头上。因为正是外邦人的伟大导师，宗徒保禄自己说了：\BibleQuote{因为我们不知道该如何祈祷才对}\BibleRef{罗 8:26}。他不只用言语教导这一教训，还用他的榜样加以说明。因为他曾出于无知，祈求\BibleQuote{那根肉身上的刺离开他}，并说这刺赐给他\BibleQuote{是为了防止他因所受启示的丰盛而过于自高}\BibleRef{格后 12:7}。但主爱他，因此没有按他所祈求的去做。然而，当宗徒说：\BibleQuote{因为我们不知道该如何祈祷才对}时，他立即补充说：\BibleQuote{但是圣神亲自以无可言喻的叹息为我们转求。那洞察人心的知道圣神的意愿是什么，因为祂按照天主的旨意为圣人转求}\BibleRef{罗 8:26}——也就是说，祂使圣人献上转求。当然，祂就是那\BibleQuote{天主派遣到我们心中，呼喊：\OriginalTerm{阿爸，父}{Abba, Father}}的圣神\BibleRef{迦 4:6}以及\BibleQuote{我们借着祂呼喊：阿爸，父}\BibleRef{罗 8:15}。因为宗徒用了两种说法——既说我们领受了\Emph{那呼喊阿爸，父的圣神}，也说我们借着祂\Emph{呼喊阿爸，父}。他的目的是通过这些不同的表述解释他使用\Quote{ \Emph{呼喊} }一词的含义：他意指\Emph{使呼喊}，因此是我们受祂的激励和推动而呼喊。因此，愿祂也教导我这事，无论何时祂愿意，如果祂知道对我有益，让我知道我的灵魂源自何处。但要由那洞察天主深奥之事的圣神来教导我，而不是一个对吹胀口袋的气息一无所知的人。不过，我绝不愿因这一无知而将你比作牲畜，因为这并非来自不可治愈的无能，而是出于纯粹的疏忽。\par

% structure: p0030 source=h2 target=subsection reason=relative-heading-rank; base=h2

\subsection{第十四章——认识肉体将复活并永远活着，比学习科学家关于其本性所能教导我们的任何知识更为卓越。}

但是，尽管关于灵魂起源的问题无疑\Quote{更高深}，胜过探讨我们吸入呼出之气来源的问题，你却仍然相信那些从圣经中学到的东西是\Quote{更高深}的，而我们从圣经中藉信德学到这些；这些是任何人的心智都无法追溯的。当然，认识肉体将复活并永远活着，胜过科学家通过仔细检验在肉体中所能发现的任何东西——尽管灵魂的临在激活了它所有无知的事物，但灵魂却不通过外在感官感知这些。同样，知道那在基督内重生和更新的灵魂将永远蒙福，也胜过发现我们对其记忆、理解和意志所无知的一切。现在，这些我称为更卓越和更好的主题，我们绝不可能发现，除非我们因默感圣经的见证而相信它们。这些圣经你也许认为自己彻底相信，以致毫不犹豫地从其中引出关于灵魂起源的明确理论。那么，首先，如果正如你所假设，你绝不应将人通过讨论和探究关于自己本性和品质所知道的归于人性本身，而应归于天主的恩赐。现在你问：\Quote{如果人对此无知，他与牲畜有何区别？}但是，如果我们应当仅仅因我们与牲畜不同就知道这事，为何还需要阅读任何东西来认识它？因为正如你读任何东西给我听都不是为了教我活着（我自己的本性使我无法不知道这一事实），同样，如果知道其他事是本性的属性，你为何要引证圣经经文让我相信关于这事？那么，难道只有那些阅读经文的人才与牲畜不同？我们不是在被创造时就与牲畜不同，甚至在我们学会阅读之前？请问，你怎么能为我们的本性提出如此高的要求，以至于仅仅凭它不同于牲畜，就已经知道如何讨论和探究灵魂的起源；而同时你却使它在认识这事上如此不熟练，以致没有人的天赋能力就无法知道这事，除非相信神圣的见证？\par

% structure: p0032 source=h2 target=subsection reason=relative-heading-rank; base=h2

\subsection{第十五章——我们不可过于智慧而超出所记载的。}

然而，你在此事上又错了；因为你为解答这个问题而选用的圣经章节，并不能证明那一点。它们所证明的是另一件事，没有它我们便无法真正度虔敬的生活，即：我们在天主那里拥有我们灵魂的赐予者、创造者和塑造者。但祂如何为它们这样做，是通过吹入新的灵魂，还是从父母遗传而来，它们并未告诉我们——除了赐给第一个人的那一个灵魂。请仔细阅读我写给天主的仆人、我们的弟兄\Person{Renatus}的信；因为我在那里已经向他指出了这一切，所以没有必要在此重复我的证明。但你喜欢我仿效你的确定理论，从而将自己陷入你所身处的那些困难之中。陷入这些困难后，你说了许多反对天主教信仰的强硬言辞；然而，如果你忠信而谦卑地反省并思考，你一定会看到，若你只知道如何在自己的无知中保持自然和一致，这对你将有多大的益处；而如果你现在还能保持这样的得体，这个益处仍然向你敞开。既然理解如此取悦于人的本性（因为确实，如果我们没有理解，就灵魂而言，我们与畜类无异），我恳求你，理解你所不理解的是什么，免得你什么也不理解；并且不要轻视任何人，他为了真正理解，而明白自己不理解他所不理解的事。至于那受默感的圣咏中的话：\Quote{人在尊贵中而不悟，如同畜类，变成它们那样；}请阅读并理解这些话，好使你以谦卑的心为自己防备这责难，而不是傲慢地将它丢给别人。这话适用于那些只把在肉身中度过的一生视为有价值的生活的人——死后毫无盼望——就如\Quote{畜类}；它不适用于那些从不否认自己实际所知的事，并总是承认自己实际不知之事的无知的人；事实上，他们意识到自己的软弱，而非自信于自己的力量。\par

% structure: p0034 source=h2 target=subsection reason=relative-heading-rank; base=h2

\subsection{第十六章——无知胜于错误。永生与永死的预定。}

我儿，不要让老年的怯懦令你少年的自信不悦。就我自己而言，如果我在天主的教导或某位属灵导师的教导下，无力理解目前关于灵魂起源的问题，我宁愿维护天主公义的旨意，即我们应当在此事上保持无知，正如在许多其他事上一样，也不愿鲁莽说出如此晦涩、以至于我既不能使他人的理智明白、自己也不理解的话；或者甚至助长异端者的主张，他们竭力说服我们，婴孩的灵魂完全无罪，其理由无非是：这种罪只会归咎于天主作为作者，因为祂明知没有为婴孩准备任何重生的助洗，却迫使无辜的灵魂（祂预先知道没有再生之洗为它们预备）变得有罪，将它们安置在有罪的内身中，却不提供那应防止它们遭受永罚的圣洗恩宠。因为事实无疑是，无数婴孩的灵魂在受洗前就离开了肉身。但愿我绝不寻求任何徒劳的努力来稀释这一严酷的事实，并说出你自己说过的话：\Quote{灵魂是注定被肉身玷污而成为有罪的，尽管它先前无罪，由于这罪它可以说得了这种报应。}又说：\Quote{即使没有圣洗，原罪也可以被赦免。}又说：\Quote{甚至天国最终也赐给了那些未受洗的人。}如果我不怕说出这些以及类似的反对信仰的有毒言论，我大概就不怕就这个问题提出某种确定的理论。那么，更好的是，我不单独就我所不知的灵魂之事争论并肯定，而只是坚持宗徒最清楚教导我们的：由于一人，所有凡从亚当而生的人都进入定罪，除非他们在基督内重生，正如祂所预定的那样，在肉身死亡之前，祂预定他们得永生，作为最仁慈的恩宠赐予者；而对于祂预定永死的人，祂也是最公义的惩罚施行者，不仅因他们在自己意志的放纵中所加的罪，也因他们的原罪，即使像婴孩那样，他们丝毫没有加添什么。这就是我对那个问题的确定看法，这样，天主的隐秘之事就保持其秘密，而不损害我自己的信仰。\par

% structure: p0036 source=h2 target=subsection reason=relative-heading-rank; base=h2

\subsection{第十七章[第十二篇]——关于灵魂的双重问题：它是\Quote{形体}吗？它是\Quote{灵}吗？什么是形体。}

现在，按主所恩准我做的，我还必须回应你那一主张，你在其中谈到灵魂时又提及我的名字，说：\Quote{我们不像那位非常能干和博学的主教\Person{奥古斯定}所声称的那样，承认灵魂既是无形体又是一个精神体。} 因此，我们首先要讨论一个问题：灵魂是否应被视为\OriginalTerm{无形体}{incorporeal}，如我所说；还是有形体\OriginalTerm{有形体}{corporeal}，如你所主张。其次，在我们的圣经中它是否被称为\OriginalTerm{精神体}{spirit}——虽然并非整体，但其自身的部分也恰当地被称为精神体。首先，我想知道你是如何定义\Emph{形体}的。因为，凡不由血肉肢体组成的东西就不是\Quote{形体}，那么地球就不能是形体，天空、石头、水、星辰以及任何这类东西也都不能是形体。然而，如果\Quote{形体}是由部分构成的东西，无论部分大小，它们占据或大或小的空间，那么我刚才提到的所有事物都是形体；空气是形体；可见光是形体；凡宗徒所看到的东西也是如此，当他说：\BibleQuote{有属天的形体，也有属地的形体。}\BibleRef{格前 15:40}\par

% structure: p0038 source=h2 target=subsection reason=relative-heading-rank; base=h2

\subsection{第十八章——第一个问题，灵魂是否为有形体；气息与风，无非是运动中的空气。}

现在，灵魂是否是这样的实体，是一个非常微妙和精细的问题。你确实以我极其称许的敏捷断言天主不是形体。但随后，当你说：\Quote{如果灵魂缺乏形体，以致成为（如某些人所喜欢假设的那样）虚空、轻飘虚幻的实体。} 你却又让我有些担忧。从这些话看来，你似乎相信，一切缺乏形体的东西都是虚空实体。如果真是这样，你怎敢说天主缺乏形体，而不担心祂是虚空实体的后果？然而，如果天主没有形体，正如你刚才所承认的；并且说祂是虚空实体是亵渎的；那么，并非所有缺乏形体的东西都是虚空实体。因此，主张灵魂是无形体的人并不必然意味着它是虚空虚浮的实体；因为他承认天主，祂不是虚空的存在，同时又是无形体。但请注意我的实际断言和你认为我所说的之间有多大的差别。我并没有说灵魂是一种空气般的实体；如果我那样说，我就要承认它是一个形体。因为空气是形体；凡懂得他们所言的人都这样宣称，每当我们谈论形体实体时。但是你，因为我称灵魂为无形体，就以为我不仅主张它纯粹虚空，而且由于这种主张，说它是\Quote{一种空气般的实体；} 而我其实一定说过，它既没有空气所具有的形体性，而且充满空气的东西也不可能是空虚无物的。你自己的皮囊比喻未能让你想起这一点。因为当皮囊充气时，除空气被压入其中外还有什么呢？它们远非空虚无物，反而因膨胀而变得沉重。但也许在你看来，气息与空气不同；虽然你的气息无非是运动中的空气；这一点可以从扇子的摇动看出。至于任何中空的器皿，你可能认为它们是空的，但你可以通过将器皿口朝下垂直放入水中，确切地确认它们实际上是满的。你看到水无法进入，因为其中充满空气。然而，如果以相反的方式（口朝上）或斜着放入，它们就会充满，水从同一开口进入，空气则从那里排出逸走。这当然通过动手实验比用文字描述更容易证明。不过，现在不是在此话题上继续耽搁的时刻；因为无论你对空气的性质有何了解，它是否有形体性，你绝不应该认为我说过灵魂是空气般的东西，而是绝对无形体。这一点甚至你承认天主是，你不致描述祂是虚空实体，同时你不得不承认祂具有不变和全能的本质。那么，如果灵魂是无形体，当我们承认天主是无形体且同时否认祂是空虚无物时，我们为何要害怕灵魂是虚空呢？因此，一个无形体之存在能够创造一个无形体的灵魂，正如生活的天主创造了生活的人；虽然，作为不变和全能的，祂并没有将这些属性传递给可变的、远为低劣的受造物。\par

% structure: p0040 source=h2 target=subsection reason=relative-heading-rank; base=h2

\subsection{第十九章【十三】——灵魂是否为\OriginalTerm{精神体}{Spirit}。}

然而，我仍不明白，你为何愿意灵魂是身体，却拒绝把它当作神魂。因为倘若它不是神魂，理由是宗徒将它与神魂分开提及，当他说\Quote{我祈求天主，愿你们整个神魂、灵魂和身体得蒙保守}（\BibleRef{得前 5:23}），同样也是它并非身体的好理由，因为他提及身体时也与它分开。如果你肯定灵魂是身体，尽管两者都被分开提及，那么你也应当允许它是神魂，尽管两者也是分开提及的。的确，灵魂更有理由被你视为神魂而非身体；因为你承认神魂和灵魂属于同一实体，却否认灵魂和身体属于同一实体。那么，凭什么原则，灵魂是身体，而它的本性不同于身体？它又不是神魂，尽管它的本性与神魂是同一的？根据你的论证，你岂不必须承认就连神魂也是身体吗？否则，如果神魂不是身体，而灵魂是身体，那么灵魂和神魂就不是同一实体。然而，你允许两者（尽管你相信它们是两件不同的事物）拥有同一实体。因此，如果灵魂是身体，神魂也是身体；因为只有在其他条件下，它们才能被视为具有同一本性。所以，根据你自己的原则，宗徒所说\Quote{你们的神魂、灵魂和身体}必须意指三个身体；然而身体，也称为肉身，具有不同的本性。在这三个身体（如你所称）中，一个本性不同，另外两个本性相同，整个人由它们组成——一件事、一个存在。现在，尽管你这样断言，你却不允许那两个本性相同的，即灵魂和神魂，具有神魂这一共同称谓；而那两个本性不同的，如你所认为，却应当具有身体这一共同称谓。\par

% structure: p0042 source=h2 target=subsection reason=relative-heading-rank; base=h2

\subsection{第二十章 [XIV.]——身体不接受天主的肖像}

但我跳过这一切，免得我们之间的讨论蜕变成名称之争而非事实之争。那么，让我们来看内人究竟是灵魂、神魂，还是两者皆是。然而，我注意到你已在文字中表达了你的意见，称内人为灵魂；因为你在说到\Quote{当那不能理解的实体凝结时，它会在身体内产生另一个身体，由它本性的力量和旋转而变得圆润和堆积，于是内人便开始显现，他在有形的鞘中被塑造，其轮廓将与外人的肖像相似}时，你谈论的就是它。由此你得出以下结论：\Quote{因此，天主的嘘气造了灵魂；是的，那来自天主的嘘气成了灵魂，一个肖像，有实体的，按其本性是有形的，像它自己的身体，并按照它的肖像而成形。}之后你接着谈论神魂，说道：\Quote{这源于天主嘘气的灵魂，不能没有其内在的感觉和理智而存在；而神魂正是如此。}那么，根据我对你陈述的理解，你意指内人是灵魂，至内人是神魂；好像后者比灵魂更低级，正如灵魂比身体低级。由此发生，正如身体接收另一个渗入自己内部空洞的身体（如你所认为，那就是灵魂）一样，反过来灵魂也必须被视为有它自己的内部空虚，在那里它可以接收第三个身体，即神魂；这样整个人由三个部分组成：外人、内人和至内人。现在，你难道还没有看到，当你试图断言灵魂是有形的时候，你身后会跟随多大的荒谬吗？我请求你告诉我，两者中哪一个要\Quote{因知识而更新，正如按创造者的肖像}（\BibleRef{哥 3:10}）？是内人，还是至内人？就我而言，的确，我不认为宗徒除了内人和外人之外，还知道内人里面另一个人，即至内人。但你必须决定，你希望哪一个按天主的肖像更新。已经拥有了外人肖像的，如何接受这个肖像？因为如果内人已经遍布外人的四肢，并且\Emph{凝结}（这是你所用的词；好像熔化的形状是从软泥中成形，从尘土中变厚而成），那么，如果这个印在它上面的图形，或者更确切地说，从身体中\Emph{表达}出来的图形，要保持它的位置，它如何能按天主的肖像重新塑造？它会有两个肖像——一个从上而来的天主的肖像，一个从下而来的身体的肖像——正如在硬币上所说，\Quote{正面和反面}。你也许会说，灵魂接受了身体的肖像，而神魂接受了天主的肖像，好像前者邻接身体，后者邻接天主；因此，实际上是至内人按天主的肖像重新塑造，而不是内人？好吧，但这种借口是无用的。因为如果至内人完全弥漫于灵魂的所有肢体，正如内人（灵魂）弥漫于身体的肢体一样；那么它现在也通过灵魂接受了身体的肖像，正如灵魂塑造了它一样；因此结果是，它没有任何办法接受天主的肖像，因为前述身体的肖像仍然印在它上面；除非像我刚才引用的硬币那样，上面有一个形状，下面有另一个形状。当你用物质的身体实体的概念来考虑灵魂时，无论你愿意与否，你都会被逼到这些荒谬的地步。但是，正如连你也十分恰当地承认，天主不是身体。那么，身体如何能接受祂的肖像？\Quote{我劝你，弟兄，不要与此世同化，反而要以心思的更新而变化}（\BibleRef{罗 12:1}）；并且不要怀有\Quote{肉性的心思，就是死亡}（\BibleRef{罗 8:6}）。\par

% structure: p0044 source=h2 target=subsection reason=relative-heading-rank; base=h2

\subsection{第二十一章 [XV.]——认识与形式同属灵魂与身体。}

但你说：\Quote{如果灵魂是无形的，那么富人在地狱里看见的是什么？他肯定认出了拉匝禄；他[不]认识亚巴郎。他对亚巴郎的认识从何而来，亚巴郎不是早已死了吗？} 用这些话，我想你认为人没有身体的形式就不能被认出和知道。因此，我想你常常站在镜子前认识自己，以免忘记自己的特征而认不出自己。但请问你，有谁比他自己更认识另一个人？有谁的脸比自己的脸更少看见？但谁能认识天主——就连你也不怀疑祂是无形的——如果认识不能像你所想象的那样没有形体形状（即只有身体才能被认出）？然而，有哪个基督徒在讨论如此重大和困难的问题时，会如此不留意圣神感动的话语，以至说：\Quote{如果灵魂是无形的，它必定缺乏形式？} 你难道忘记了，在那话语中你读到过\BibleQuote{教义的形式}？\BibleRef{罗 6:17} 你也忘记了，论及耶稣基督在祂披上人性之前，经上写着祂是\BibleQuote{在天主的形式内}？\BibleRef{斐 2:6} 那么，你怎么能说：\Quote{如果灵魂是无形的，它必定缺乏形式；} 而你却听到\Quote{天主的形式}——你承认祂是无形的——并且如此表达，好像形式只能存在于身体中？\par

% structure: p0046 source=h2 target=subsection reason=relative-heading-rank; base=h2

\subsection{第二十二章——名称不暗示形体性。}

你也说：\Quote{当形式不被区分时，名称就不再被赋予；而没有人的指定，就没有名称的赋予。} 你的目的是证明亚巴郎的灵魂是形体的，因为祂可以被称呼为\Quote{父亲亚巴郎}。我们已经说过，即使没有身体也有形式。然而，如果你认为没有身体就没有名称的指定，我必须请你数一数这段圣经中出现名称的次数：\BibleQuote{但圣神的效果是：仁爱、喜乐、平安、忍耐、良善、温和、忠信、柔和、节制，}\BibleRef{迦 5:22} 并告诉我你是否不认出这些名称所指的事物本身；或者你认出它们以至于能描绘出一些身体的轮廓。来吧，告诉我，只以爱德为例，它的肢体、形状、颜色是什么？因为如果你不是头脑空虚，这些属性绝不能被视为空虚的东西。然后你接着说：\Quote{那被恳求帮助者的容貌和形式当然必须是形体的。} 好吧，让人们听到你的话；愿没有人恳求天主的帮助，因为没有人在祂身上能看到任何有形体的东西。\par

% structure: p0048 source=h2 target=subsection reason=relative-heading-rank; base=h2

\subsection{第二十三章 [XVI.]——比喻言语不可按字面理解。}

\Quote{总之，}你说，\Quote{在这个比喻中，肢体被归于灵魂，好像它真的是一个身体。} 你坚持认为\Quote{通过眼睛理解整个头，}因为经上说\BibleQuote{他举起了他的眼睛}。你又说了，\Quote{通过舌头指颚，通过手指指手，}因为经上说\BibleQuote{请打发拉匝禄，用指尖蘸点水，来凉凉我的舌头。}\BibleRef{路 16:24} 然而，为了避免将有形体的属性归于天主的矛盾，你说\Quote{通过这些术语必须理解为无形的功能和能力；}因为你非常恰当地坚持认为天主不是有形体的。那么，为什么这些肢体的名称在天主身上不证明形体性，而在灵魂上却证明呢？是因为当论及受造物时，这些术语必须按字面理解，而当论及造物主时，则仅作为隐喻和形象说法？那么你将不得不给我们有实际物质体质的翅膀，因为不是造物主，而是人类受造物说了：\BibleQuote{如果我像鸽子一样取了翅膀。} 此外，如果比喻中的富人有一个有形的舌头，因为他喊道\BibleQuote{让他凉凉我的舌头}，那么看起来我们的舌头，即使我们还带有肉身，它本身就有物质的手，因为经上写着：\BibleQuote{生死在舌头的\Emph{手里}。}我想甚至你自己也清楚，罪既不是受造物，也不是物质实体；那么它为什么有脸呢？难道你没有听到圣咏作者说：\BibleQuote{在我的骨头里没有平安，在我罪的\Emph{面前}？}\par

% structure: p0050 source=h2 target=subsection reason=relative-heading-rank; base=h2

\subsection{第二十四章——亚巴郎的怀抱——其意义为何。}

至于你认为\Quote{所提到的亚巴郎的怀抱是有形的，} 并进一步断言\Quote{通过它意指他的整个身体，} 我担心你会被认为（即使在这样的主题中）试图开玩笑和引人发笑，而不是庄重严肃地行事。因为你不可能愚蠢到认为一个人的物质怀抱能容纳那么多灵魂；不，用你自己的话说，\Quote{承载那么多有功之人的身体，就像天使把拉匝禄带到那里一样。} 除非你恰好认为只有他的灵魂才配进入那怀抱。那么，如果你不是在开玩笑，也不希望犯幼稚的错误，你必须通过\Quote{亚巴郎的怀抱}来理解那遥远而分离的安息与平安的居所，亚巴郎现在就在其中；并且对亚巴郎所说的话不仅指他个人，而且指向他被立为万民之父\BibleRef{创 17:5}，他被呈现给众人作为信德的首要榜样来效法；正如天主愿意被称为\BibleQuote{亚巴郎的天主、依撒格的天主、雅各伯的天主，}尽管祂是无数之人的天主。\par

% structure: p0052 source=h2 target=subsection reason=relative-heading-rank; base=h2

\subsection{第二十五章 [XVII.]——脱离肉身的灵魂可能以形体形态思考自己。}

但切不可以为我说这一切，是好像否认真理：死者的灵魂，如睡眠中的人，可以借其身体的相似形象，思想或善或恶的事。因为在梦中，当我们遭受任何严酷和烦扰的事时，我们当然仍是自己；若这苦楚在醒来时还不消失，我们便会经历极大的痛苦。然而，认为我们在梦中匆匆或飘荡于各处时所凭借的是真实的身体，那是只对此类题目作过粗略思考的人的想法；事实上，主要正是通过这些幻象，灵魂被证明是非形体的；除非你连我们在梦中除自身外也常看见的东西，如天、地、海、日、月、星、河、山、树、或动物，也称之为形体。谁若把这些幻影当作形体，就极其愚蠢，虽然它们的确很像形体。那些确证为具有神圣意义的现象，无论见于梦中还是神魂超拔中，也属于这类性质。谁能追究或描述它们的来源或构成材料呢？它是精神的，而非形体的，这毫无疑问。这类看似形体而实非形体的东西，在清醒的人的思想中形成，并深藏于记忆之中，然后以某种奇妙而不可言喻的方式，从这些隐秘之处显现于记忆的运作中，仿佛可触摸地呈现在眼前。因此，若灵魂是物质形体，它绝不可能在思想范围和记忆空间中容纳如此众多而巨大的物质形体的形状；因为，按照你自己的定义，\Quote{它在本身形体方面并不超越这个外在的身体。}因此，既然本身没有量度，它有什么能力去把握广大形体、空间和区域的模样呢？那么，即使灵魂无身体时，它本身在自身身体的相似形象中显现，又有什么奇怪呢？因为灵魂在梦中从不以自己的身体显现；但仍在其自身身体的相似形象中，穿越熟悉和不熟悉的地方，看见许多或悲或喜的景象。然而，我想你自己也不会大胆主张，灵魂在梦中似乎拥有的肢体和身体的形态具有真正的形体性。若按此推论，它似乎攀登的山就是真的山，它似乎进入的房屋就是物质的房屋，它似乎倚靠的树就是真正的树，有真实的木头和体积，它想象自己饮用的水就是真的水。如果灵魂本身是形体的，当它以身体的相似形象穿梭其中时，所有与它打交道的、如同形体的事物，都将是无可置疑的形体。\par

% structure: p0054 source=h2 target=subsection reason=relative-heading-rank; base=h2

\subsection{第26章 [XVIII.]——\Person{圣裴伯乐雅}在梦中似乎变成了男人，并与某个埃及人角力。}

必须注意一下殉道者的一些神视记叙，因为你以为可以从中引出一些证据。例如，\Person{圣裴伯乐雅}在梦中似乎变成了男人，并与一个埃及人角力。谁能否认，在那显见的身体形象中，那是她的灵魂，而不是她的身体呢？她的身体当然仍保持其女性的性别，躺在床上，感官沉入睡眠，而她的灵魂则在男人身体的相似形象中搏斗。对此你有什么可说？那男性形象是真实的形体，还是虽有身体外表却根本不是形体？请选择其一。如果它是形体，为何不保持其性别的完整？因为在那个女人的肉身上找不到男性的生殖功能，由此不能通过你所称的\OriginalTerm{凝固}{congelation}过程，塑造出这个男性身体的相似形象。那么，我们若愿意就可断定，当她的身体在她睡觉时仍然活着，尽管她的灵魂在角力，她仍保持自己自然的性别，当然，包含在她活着时属于她的一切适当肢体中，并且仍拥有她最初形成的身体形状和轮廓。她没有像死亡那样交出关节和四肢；也没有因死亡力量的运作，从转变能力中撤出任何已经定型的肢体。那么，她的灵魂从哪里得到那个男性身体，在其中她似乎与对手角力呢？然而，如果这[男性形象]不是形体，尽管它有这样一副模样，其中却可感受到真实的搏斗或真实的喜悦，那么你难道现在还不明白，就所能理解的而言，灵魂中可以有某种与物质形体相似的形态，而灵魂本身却不是形体吗？\par

% structure: p0056 source=h2 target=subsection reason=relative-heading-rank; base=h2

\subsection{第27章。——身体受伤时，灵魂是否受伤？}

那么，如果在已亡者中显现类似的事，而灵魂在其中认出彼此，不是靠身体，而是靠身体的幻象，又如何呢？当我们受苦时，即使在梦中，尽管活动的只是身体肢体的幻象，而非肢体本身，痛苦却不仅是幻象，而是真实的；快乐的感觉也是如此。然而，由于圣伯尔伯杜亚尚未去世，你或许不愿从这一事例（尽管它对此问题有很强说服力）为自己定下确切规则，即你认为我们在梦中所见的身体幻象具有何种性质。如果你承认它们像身体，但并非实际的身体，那么整个问题就解决了。但她的弟弟狄诺克拉特已经死了；她看见他带着生前所受并导致他死亡的伤口。你曾竭力论证：当肢体被切除时，灵魂不应被认为因截肢而遭受同等的损失，这种论证的根据何在？请看，伤口是加在狄诺克拉特的灵魂上，当灵魂居住在那身体中时，伤口的力量将其逐出身体。那么，你的观点如何能正确呢？你说：\Quote{当身体的肢体被割断时，灵魂会从打击中撤离，凝结后退缩到其他部位，以致没有一部分随着身体所受的伤口被切除}，即使那人在失去肢体时正在沉睡且毫无知觉。你赋予灵魂如此大的警觉性，以致即使打击落在肉体的任何部位而灵魂毫无所知（当它沉浸在梦境中时），它会立刻凭着上智本能撤离，从而使得任何打击、伤害或截肢都无法加在它身上。然而，你可以随意绞尽脑汁回答这个自然的问题：灵魂如何撤走自身存在的部分，退回到自身之内，以致每当身体肢体被割断或骨折时，灵魂自身不会遭受任何截肢或骨折；但我必须请你看看狄诺克拉特的事例，并向我解释：为什么他的灵魂没有从他身体受致命伤的部位撤离，从而避免自身受到在他身体死后仍然明显可见于他脸上的伤害？或许你乐意让我们认为这些现象是身体的幻象而非实在；这样，正如实际上不是伤口的看起来像伤口，实际上不是身体的也呈现出身体的形状。如果灵魂能被伤害身体的人所伤，我们岂非有充分理由担心它也能被杀害身体的人所杀？然而，主自己最清楚地宣称这是不可能的。\BibleRef{玛 10:28} 狄诺克拉特的灵魂无论如何也不会因杀死他身体的打击而死亡：它的伤口也只是表面的；因为灵魂不是有形体的，所以它并未像身体那样真正受伤；它拥有身体的相似，也分享了其伤口的相似。不过还可以说，在它虚幻的身体中，灵魂感到了真实的痛苦，这痛苦由身体伤口的影子所表示。正是从这真实的痛苦中，他因他圣洁姐姐的祈祷而获得了释放。\par

% structure: p0058 source=h2 target=subsection reason=relative-heading-rank; base=h2

\subsection{第二十八章 —— 灵魂是否因身体的缺陷而变形？}

再者，你说\Quote{灵魂从身体获得形状，随着身体的增长而增长和扩展}，却不考虑如果一个人的手臂在婴儿时被截肢，那么他青年或老年时的灵魂会变成什么怪物？你说\Quote{灵魂的手会收缩，以便不与身体的手一同被切除，并通过凝结缩入身体的其他部分}。照此，上述灵魂的手臂无论停留在何处，都保持最初接受身体形状时的长度，因为它失去了通过同等扩张而增长的形式。因此，那个在婴儿时失去手的青年或老人的灵魂确实有两只手（因为收缩的那只逃脱了肢体截肢），但根据假设，其中一只是成年青年或老人的手，而另一只只是婴儿的手，与截肢发生时一样。相信我，这样的灵魂并非按照身体的模型和形状被造，而是在错误的变形印记下被虚构出来。在我看来，除非靠天主的助佑，你彻底而平静地审视做梦者的异象，并由此说服自己有些形式不是实在的身体，只是身体的幻象，否则你无法从这错误中被救出。现在，尽管我们以为是身体的对象也属于同一类，但关于已亡者，我们可以从睡着的人那里事后推断。因为圣经描述已亡者为\Quote{睡着}并非徒然，\BibleRef{得前 4:13} 仅是因为在某种意义上\Quote{睡与死相似}。\par

% structure: p0060 source=h2 target=subsection reason=relative-heading-rank; base=h2

\subsection{第二十九章 [XIX.]——灵魂是否也将身体的衣物带走？}

如果灵魂确实是身体，而形式也确实是它在梦中看见自己的一个有形模样，因为它从包裹它的身体那里获得了自己的表现：那么，一个失去肢体的人，在梦中绝不会看到自己失去了被截去的部位，尽管实际上他已经失去了它；相反，他总会显得完整无缺，因为灵魂本身没有任何部分被切除。但是，既然人们有时看见自己完整，有时看见肢体残缺，而这正是他们实际的情况，那么，这一事实除了表明：灵魂，无论是对于它在梦中看见的其他事物，还是对于身体，它所携带的并非它们的实在，而只是它们的相似性之外，还有什么呢？然而，灵魂的喜悦或悲伤、快乐或痛苦，无论是发生在实在的身体中还是出现在表象的身体中，都是真实的感受。你自己不是也说过（而且完全正确）：\Quote{灵魂不需要食物和衣服，只有身体才需要}？那么，为什么地狱中的富翁渴望一滴水呢？\BibleRef{路 16:24}为什么\Person{圣撒慕尔}死后（如你自己注意到的）仍然穿着他平日的衣服显现呢？\BibleRef{撒上 28:14}前者难道是想用水来修复灵魂的亏损，就像修复肉体的亏损一样吗？后者是穿着衣服离世的吗？在前一种情形中，确实有一种真实的痛苦折磨着灵魂；但并没有一个需要食物的真实身体。而后一种情形，他看起来穿着衣服，并不是因为他是一个真实的身体，而只是灵魂，具有身体的表象并穿着衣服。因为，尽管灵魂伸展和收缩以适应身体的肢体，但它并不能同样地适应衣服，以使自己的形式与衣服相吻合。\par

% structure: p0062 source=h2 target=subsection reason=relative-heading-rank; base=h2

\subsection{第30章——身体性是认知所必需的吗？}

但有谁能够探究出，即使是不善良的灵魂，在脱离可朽坏的身体之后，拥有何等认识能力，以致能够凭内在感官观察并认出那些像它们一样邪恶的灵魂，甚至是善良的灵魂——无论是在那些并非实在是身体、而是身体相似的状态中，还是在灵魂的喜怒哀乐等情感中，而这些情感并不包含任何肢体轮廓？由此产生了比喻中的富翁，尽管在痛苦中，却认出了\Quote{亚巴郎}，他从未见过他的面容和身形，但灵魂虽无形体，却能领会他身体的相似。\BibleRef{路 16:23}但谁能正确地说，他曾认识某个人，除非他知道了他的生活和性情——这些当然既无物质实体，也无颜色？我们正是这样认识自己，比别人认识我们更确定，因为我们的自我意识和性情完全呈现在我们面前。我们清楚觉察到这一点，但在其中却看不到任何形体的相似。然而，我们并不能在另一个人身上觉察到自己本性的这种内在特质，即使他就在我们眼前；尽管在他不在时，我们能够回忆、认出并想起他的面容。但我们的面容却不能同样地被回忆、认出和想起；然而，我们却完全真实地说，我们自己比他人更了解自己，这清楚地表明，哪一方面的知识更强、更真实。\par

% structure: p0064 source=h2 target=subsection reason=relative-heading-rank; base=h2

\subsection{第31章[第二十章]——论灵魂中知识的区分。}

那么，既然灵魂有一种功能，我们通过它凭借五种肉体感官感知实在的物体；另一种功能，使我们能够辨别非形体的、类似身体的东西（借此我们也能看见自己，如同仅仅类似身体）；第三种功能，我们借此对其能力更确定、更深刻地洞见其对象，这些对象既非形体也非类似于形体的东西——例如\OriginalTerm{信德}{faith}、\OriginalTerm{望德}{hope}、\OriginalTerm{爱德}{charity}——这些事物既无肤色、也无情感、也无诸如此类的东西：在所有这些功能中，我们应当更专注、甚至更熟悉地专注于哪一个，并在认识天主、按那创造我们的肖像而更新？岂不是在我刚才列为第三的那个功能上吗？在那里我们当然不会经历性别差异或其类似物。\par

% structure: p0066 source=h2 target=subsection reason=relative-heading-rank; base=h2

\subsection{第32章——赋予灵魂所有性器官却无性别的矛盾。}

因为灵魂的那种形式，无论是男性的还是女性的，具有男女特有的肢体区别，这不仅仅是身体的相似，而是实际的身体，那么它要么是男性，要么是女性，不管你愿不愿意，正如它看起来是男人或女人一样。但如果你的意见正确，灵魂是一个身体，甚至是一个活的身体，那么它既有肿胀下垂的乳房，又没有胡须，它有子宫和女人的一切生殖器官，却终究不是女人。那么，我的说法不是更符合真理吗：灵魂确实有眼睛、舌头、手指，以及所有其他类似身体肢体的部分，但整个只是身体的相似，而不是真实的身体？我的说法可以经受普遍检验；每个人都能在自己身上验证，当他回忆不在场的朋友的形象时；他也能确凿地验证，当他回忆自己和他人在梦中出现的形象时。然而，就你而言，在整个自然界中，你无法举出你所想象的那种怪物的例子：有一个真实活着的女人身体，却没有女人的性别。\par

% structure: p0068 source=h2 target=subsection reason=relative-heading-rank; base=h2

\subsection{第33章——死后复生的凤凰。}

如今，你关于凤凰的说法与我们当前的主题毫无关系。因为凤凰象征身体的复活，它并不否定灵魂的性别；如果确实如人们所想，它在死后重生。然而，我想您认为，除非像年轻人那样大谈特谈凤凰，否则您的论述就缺乏充分的可信度。如今，你在你的鸟的身体里发现了雄性的生殖器官，却并非雄鸟；或者发现了雌性的生殖器官，却并非雌鸟吗？但，我恳求你，反省一下你所说的——你试图构建并推荐给我们接受的理论是什么。你说，灵魂伸展到身体的各个肢体，通过凝固变得僵硬，接受了整个身体的完整形状，从头冠到脚底，从最内部的骨髓到皮肤的外表面。照此推论，在女性身体的情况下，它必定接受了女人身体的所有内在附属物，却仍然不是女人！请问，为什么在一个真实的活体中，所有成员都是女性的，而整体却非女人？为什么所有都是男性的，结果却不是男人？谁能如此鲁莽，竟相信、宣称并教导这一切？难道灵魂从不生育？那么，当然，公骡和母骡就没有雄雌之分了。难道没有血肉身体的灵魂就不能交合？然而，被阉割的男人也享有这种匮乏；虽然他们的过程和运动都被剥夺，但他们的性别并未被移除——他们的男性成员还留下一些微小的残余。从来没有人说阉人不是男性。那么你的意见——即使阉人的灵魂也拥有完好的生殖器官，而且根据你的原则，这些器官即使从他们的身体结构上完全移除，也将在他们的灵魂中保持完整——现在又当如何呢？因为你说，当那块肉开始被割除时，灵魂知道如何退出，使得被切除的部分形状不会丢失；虽然它通过凝聚而覆盖在上面，但以极快的速度退却，并如此深藏于内而得以保全；然而，那携带了全套男性生殖器官附属物到那里去的灵魂，恐怕在来世不能算是男性，因为即使它在身体上没有其他标记，仅凭那些器官它就是男性。这些意见，我儿，毫无真理可言；如果你不承认灵魂中有性别，那么也没有身体可言了。\par

% structure: p0070 source=h2 target=subsection reason=relative-heading-rank; base=h2

\subsection{第三十四章 [XXI.]——先知异象}

并非每一个身体的外貌本身就是一个身体。你入睡时会看到这点；但当你再次醒来时，要仔细辨别你所看到的。因为在梦中，你似乎拥有一个身体；但那实在不是你的身体，而是你的灵魂；也不是真实的身体，而是身体的外貌。你的身体躺在床榻上，但灵魂却在行走；你身体的舌头沉默不语，但在梦中你的灵魂却在说话；你的眼睛紧闭，但你的灵魂却是清醒的；当然，你身体伸展在床上的肢体是活的，而非死的。因此，你所认为的灵魂那种凝固的形态，还没有像从鞘中抽出一样被提取出来；然而在其中却看到了你肉身的完整无缺的外貌。属于这一类身体相似物的——它们不是真实的身体，虽然看似如此——就是你在圣经中读到过的先知异象中的一切显现，然而你并不理解；这些异象也标志着所有时间——现在、过去和未来——所要发生的事。你弄错了这些，不是因为它们本身具有欺骗性，而是因为你没有按照应有的方式来接受它们。因为在同一默示性的异象中，看到\Quote{殉道者的灵魂}（\BibleRef{默 6:9}），也看到\Quote{一只像是被宰杀的羔羊，有七只角}（\BibleRef{默 5:6}）；还有马和其他动物，以比喻的方式前后一致地描述；最后，有星辰坠落，大地像书卷般卷起（\BibleRef{默 6:13}）；尽管如此，世界并没有真正崩塌。因此，如果我们明智地理解这一切，虽然我们说它们是真实的显现，但我们并不称它们为真实的身体。\par

% structure: p0072 source=h2 target=subsection reason=relative-heading-rank; base=h2

\subsection{第三十五章——天使是以真实的身体向人显现吗？}

然而，要非常仔细地讨论这类身体相似物——无论天使，无论是善的还是恶的，只要他们以人的形象或任何身体的形象显现，是否以此方式显现？还是他们拥有真实的身体，并以此真实的身体状态显示自己？或者，当人做梦或出神时，他们以这些形式——不是身体，而是身体的相似物——被感知，而对清醒的人，他们则呈现出可见的、必要时可触摸的真实身体？我认为在这本书中完全不必探索和充分处理这些问题。此时关于灵魂的非物质性已经论述得足够多了。如果你宁愿坚持你的意见，认为灵魂是物质性的，那么你首先必须定义什么是\Quote{\Emph{身体}}；以免，也许结果是我们对事情本身意见一致，却徒劳地在名称上纠缠。然而，如果你在灵魂中设想那种被所有有学问的人称为\Quote{身体}的物质——我指的是占据空间部分的物质，较小的部分占据较小的空间，较大的部分占据较大的空间——通过长、宽、高的不同关系，我认为你现在已经能够明智地观察到那些荒谬的结论了。\par

% structure: p0074 source=h2 target=subsection reason=relative-heading-rank; base=h2

\subsection{第三十六章 [XXII.]——他转而讨论关于灵魂的第二个问题：它是否被称为灵。}

现在我要说明的是，虽然\Emph{灵}这一称呼正确地用于灵魂的一部分，而非全部——正如宗徒所说：\BibleQuote{你们的整个灵、灵魂和身体}（\BibleRef{得前 5:23}），或者按照\WorkTitle{约伯记}中更为生动的说法：\BibleQuote{你将我的灵魂与我的灵分开}（\BibleRef{约 7:15}）——然而整个灵魂也被称为这名；尽管这个问题似乎更是一个名称问题，而非实质问题。因为事实上，在灵魂中确实有某种东西被正确地称为 \Quote{灵}，而（这暂且不论）它也同样被恰当地称为 \Quote{灵魂}；我们目前的争论并非关于这些事物本身；主要因为我自己肯定承认，而你也同样认为，那个我们用以推理和理解的东西被正确地称为灵，然而这些事物又被区分性称呼，正如宗徒所说：\BibleQuote{你们的整个灵、灵魂和身体}。然而，这位宗徒似乎也将这个灵描述为 \Emph{理智}；如他所说：\BibleQuote{这样看来，我用理智服事天主的法律，但用肉身服事罪的法律}（\BibleRef{罗 7:25}）。这意思正是他在另一处所表达的：\BibleQuote{因为肉体私欲与灵相争，灵与肉体相争}（\BibleRef{迦 5:17}）。他在前一节称为 \Emph{理智} 的东西，在后一节中应理解为称为 \Emph{灵}。并非如你解释的那样：\Quote{强调的是整个理智，它由灵魂和灵组成}——我不知道你从何处得到这个观点。我们通常用 \Quote{理智} 仅指我们的理性官能和智力官能；因此，当宗徒说：\BibleQuote{更新你们心神的灵}（\BibleRef{弗 4:23}），不正是说更新你们的理智吗？因此，\Quote{心神的灵} 无非就是理智，正如 \Quote{肉身的身体} 无非就是肉身；经上记载：\BibleQuote{脱去肉身的身体}（\BibleRef{哥 2:11}），宗徒称肉体为 \Quote{肉身的身体}。的确，他从另一个角度将其称为人的灵，这与理智有明显区别：他说：\BibleQuote{如果我用舌头祈祷，我的灵祈祷，但我的理智却没有结果}（\BibleRef{格前 14:14}）。然而我们现在所讨论的并非那与理智有别的灵；这涉及一个确实困难的问题。因为圣经在很多地方、以不同意义提到灵；但就我们现在所讨论的、我们用以推理、理解和智慧的灵而言，我们双方都同意它被称（而且正确地称为）\Quote{灵}，其意义并不包括整个灵魂，而只是它的一部分。然而，如果你争辩说灵魂不是灵，理由是理智被明确称为 \Quote{灵}，你同样可以否认雅各的整个后裔被称为以色列，因为除了犹大支派外，同样的称呼也明确地、单独地由当时在撒玛利亚组织起来的十个支派所承担。但我们何必在这里继续耽搁呢？\par

% structure: p0076 source=h2 target=subsection reason=relative-heading-rank; base=h2

\subsection{第37章 [XXIII.]——\Quote{灵}一词的广义和狭义}

但现在，为了更易于阐明，我请求你注意：在圣经叙述我们主死亡的事件中，灵魂也被称为灵，如此写着：\BibleQuote{祂低下头，交付了灵魂。}\BibleRef{若 19:30} 现在，当你听到或读到这些话时，你愿意将它们理解为以部分代表全体，而不是因为灵魂也可以称为灵。但我为了能更迅速证明我所说的，实际上要立刻方便地召唤你自己作为证人。因为你如此定义灵，以至于牲畜似乎没有灵，只有灵魂。非理性的动物之所以如此称呼，是因为它们没有理智和理性的能力。因此，当你劝诫人自身认识自己的本性时，你这样说：\Quote{如今，既然善良的天主没有无目的地创造任何事物，祂便造了人自身作为理性的动物，能够理解，具有理性，为感觉所赋予活力，从而能够明智地将所有缺乏理性的事物进行安排。}在你的这些话中，你已经清楚地肯定了确实最真实的事：人具有理性并能够理解，而这当然是缺乏理性的动物所没有的。并且你根据这个观点，引用了圣经的一段话，并采用其语言，将无理解力的人比作牲畜，它们当然没有理智。类似的陈述出现在圣经的另一段话中：\BibleQuote{不要像马或骡子一般，它们没有理解。}既然如此，我也希望你观察你如何定义和描述灵，当你试图区分它和灵魂时：\Quote{这个灵魂，}你说，\Quote{它源自天主的嘘气，不可能没有它自身的内感和理智；这就是灵。}稍后你加上：\Quote{虽然灵魂使身体有生命，但因它有感觉、智慧和精力，就必定有灵。}再稍后你说：\Quote{灵魂是一回事，灵——即灵魂的智慧和感觉——则是另一回事。}在这些话中，你足够清楚地表明了你如何理解人的灵：它甚至是我们理性的能力，灵魂借此行使感觉和智力——当然，不是身体感官所感到的感觉，而是最内在感觉的运作，由此产生出\Quote{情感}一词。正因如此，毫无疑问，我们被置于野兽之上，因为它们没有理性。因此，这些动物没有\Emph{灵}——即理智、理性和智慧的感觉——而只有\Emph{灵魂}。因为关于它们曾说过：\BibleQuote{水要滋生有灵魂的爬行生物；}\BibleRef{创 1:20} 又说：\BibleQuote{地要生出有灵魂的生物。}\BibleRef{创 1:24} 的确，为使你有最充分和最清楚的保证，灵魂在圣经用语中也称为灵，野兽的灵魂有灵的称谓。当然，牲畜没有你，我亲爱的弟兄，所定义的与灵魂不同的那种灵。因此，很明显，野兽的灵魂在一般意义上可以恰当地称为\Quote{灵}，正如我们在\WorkTitle{训道篇}中读到：\BibleQuote{谁知道人的灵是否上升；\Emph{野兽的灵}是否下降入地？}\BibleRef{训 3:21} 同样，关于洪水的毁灭，圣经作证：\BibleQuote{凡在地上行动的有血肉的，飞鸟、牲畜、走兽，和在地上爬行的各种爬虫，以及所有的人，都死了；凡有生命之灵的也都死了。}\BibleRef{创 7:21} 这里，如果我们除去一切怀疑的争议，我们理解\Emph{灵}这个术语在一般意义上与\Emph{灵魂}同义。这个术语的意义如此广泛，甚至天主被称为\BibleQuote{灵；}\BibleRef{若 4:24} 空气的狂风，虽有物质实体，也被圣咏作者称为暴风的\Quote{灵。}因此，基于所有这些理由，你将不再否认被称为灵魂的也称为灵；我认为我已经从圣经的篇章中引用了足够多来确保你的同意，在那些野兽的灵魂（它没有理解）被称为灵的段落中。如果你接纳并明智地思考我们讨论中关于灵魂的非形体性的内容，就没有进一步的理由让你因我曾说我相信灵魂不是身体而是灵而感到冒犯——既因为它被证明不是物质的，又因为在一般意义上它被称为灵。\par

% structure: p0078 source=h2 target=subsection reason=relative-heading-rank; base=h2

\subsection{第38章 ——\Person{维克托}的主要错误再次指出。}

因此，如果你拿起这些为了回应你的意见而出于真诚与挚爱所写的书卷，并以彼此相爱的态度来阅读它们；如果你留意你在第一卷开始时自己曾宣称的，\Quote{你急于不坚持任何你个人认为不可靠的意见}，那么我恳求你特别警惕我在本论著前一卷中向你指出的那十一个错误。不要说：\Quote{灵魂属于天主的含义，是天主并非从无中、也非从其他事物中，而是从祂自己的本性中创造了它}；或说：\Quote{既然施予的天主本身永存，祂也永远在无限的时间中施予灵魂}；或说：\Quote{灵魂因肉体而失去了一些它早在肉体之前就拥有的功德}；或说：\Quote{灵魂借着肉体恢复它古老的状态，并通过那曾使其遭受污染的同一肉体而重生}；或说：\Quote{灵魂甚至在犯罪之前就已经该当有罪}；或说：\Quote{未能领受圣洗重生的婴儿仍可获得其原罪的赦免}；或说：\Quote{主预定要受洗的人，可以被带离祂的预定，或在全能者所预定的事在他们身上成就之前就死去}；或说：\Quote{那些在领洗前去世的人，圣经论及他们说：\BibleQuote{他迅速被提去，免得恶念改变他的悟性}}，以及其余类似的段落；或说：\Quote{在天主的国以外，有一些属于主所说的\InnerQuote{许多}住所，即在祂父的家里}；或说：\Quote{应当为那些未领洗而离世的人奉献基督体血的祭献}；或说：\Quote{任何未领受基督的圣洗而死的人，暂时被接纳进入乐园，日后甚至得以享天国的福乐}。\newline{}我的孩子，尤其要警惕这些意见，并且，既然你渴望成为错误的战胜者，不要因\Quote{文森休斯}这个姓氏而沾沾自喜。当你对任何事无知时，不要自以为知道；为了获得真知，要学会如何不知。因为我们在\Quote{天主的隐秘事}中假装无知，对未知之事随意构建理论并将其当作已知，产生并捍卫错误如同真理，便是犯罪。至于我自己对灵魂是每人生时被新造，还是由父母传生这一问题的无知（然而这种无知被一种不可亵渎的信念所修正，即确信灵魂确实由神圣创造者所造，虽非出自祂的实体），我相信你这位充满爱心的人现在已被说服：它要么根本不应受到指责，要么如果需要指责，那也应由一位有学问足以彻底消除它的人来执行；对于我的其他见解也是如此：灵魂虽然具有身体的无形相似，但它们本身不是身体；并且，在不损害灵魂与精神之间自然区别的情况下，灵魂在广义上实际上被称为精神。如果我不幸未能说服你，那我宁可让读者来判断我所提出的观点是否未能使你信服。\par

% structure: p0080 source=h2 target=subsection reason=relative-heading-rank; base=h2

\subsection{第三十九章——结束劝言。}

如果你（或许有此可能）想知道是否还有许多其他我认为需要在你著作中修正的地方，那么对你而言，来找我并非难事——当然，不是作为学生来找老师，而是作为盛年之人来找年迈者，作为强壮之人来找孱弱者。虽然你本不该出版你的著作，但一个人若亲口承认其受纠正的公正性，比起从任何为错误辩护者口中获得赞扬，更有真正而更大的荣耀。我不愿相信所有那些聆听你朗读上述著作并慷慨赞美你的人，要么自己原先就持有正统教义所不认同的意见，要么被你诱导而接受了它们；我仍不禁认为，他们敏锐的心智被你滔滔不绝而持续不断的雄辩所钝化，因而未能对你的讲论内容给予足够关注；或者，当他们在某些情况下能够理解你所说的内容时，他们赞美你与其说是由于任何对真理的非常清晰的陈述，不如说是由于你语言的丰富、心智的敏捷和才智的充沛。因为对一个青年人的口才，往往因着对未来的期待而给予赞美、声誉和善意的关注，尽管此时它还不具备一位完全成熟的导师那种醇熟的完美和可靠性。因此，为使你自己获得真智慧，并使你的言谈不仅能愉悦，甚至能造益他人，在你摒除他人危险的喝彩之后，应当以良心审慎地看管自己的言语。\par

\SourceRecord{Translated by Peter Holmes and Robert Ernest Wallis, and revised by Benjamin B. Warfield.；Nicene and Post-Nicene Fathers, First Series；Vol. 5.；Edited by Philip Schaff.；Buffalo, NY: Christian Literature Publishing Co.,；1887.；<http://www.newadvent.org/fathers/15084.htm>.}
